\documentclass[11pt]{article}
% 必须显式写 a4paper：geometry 不指定纸张时沿用 article 的 letterpaper 默认，
% 成品实测是 612x792pt（US Letter），\textheight 只有 656pt——既是排版警告的来源，
% 也是投稿格式错误。A4 下 \textwidth≈459pt、\textheight≈706pt。
\usepackage[a4paper,margin=2.4cm]{geometry}
\usepackage{graphicx}
\usepackage{amsmath,amssymb}
\usepackage{booktabs}
\usepackage[hidelinks]{hyperref}
\usepackage{microtype}
\usepackage{float}
\usepackage{longtable}
\usepackage{array}
\usepackage{fontspec}
\usepackage{xcolor}
\newfontfamily\cjkfont{Noto Sans SC}[Scale=0.92]
% 中文段落只有 \allowbreak 断点、没有可伸缩空白，默认 tolerance 下 TeX 找不到
% 满意解就只能让整行溢出（实测最大 93pt）。放宽 tolerance 并允许 emergency
% 伸缩——它们只在原本无解的位置生效，不会把已经断好的行改差。
\tolerance=3000
\emergencystretch=1.5em
% longtable（表 2 有 25 行）与图注、节标题交替出现时，TeX 在 \output 阶段会尝试
% 组装一个高于 \textheight 的页面盒并报 "vbox too high"，实测中文版 232pt。加了
% \raggedbottom 后底部不强行对齐，该警告消失；实际渲染本就没有内容越界（18 页逐页
% 量过正文块底边），所以这是消噪而非修错。
\raggedbottom
\title{CAIEMS{\cjkfont ：\allowbreak 面\allowbreak 向\allowbreak 施\allowbreak 工\allowbreak 事\allowbreak 故\allowbreak 叙\allowbreak 述\allowbreak 结\allowbreak 构\allowbreak 化\allowbreak 处\allowbreak 理\allowbreak 的\allowbreak 探\allowbreak 索\allowbreak 性} LLM {\cjkfont 数\allowbreak 据\allowbreak 工\allowbreak 程\allowbreak 管\allowbreak 道}}
\author{Tianhao Sun\textsuperscript{1}, [Supervisor name]\textsuperscript{1,*}}
\date{2026-09-14}
\begin{document}
\maketitle

\section*{{\cjkfont 摘\allowbreak 要}}
{\cjkfont 建\allowbreak 筑\allowbreak 施\allowbreak 工\allowbreak 事\allowbreak 故\allowbreak 报\allowbreak 告\allowbreak 是\allowbreak 安\allowbreak 全\allowbreak 知\allowbreak 识\allowbreak 中\allowbreak 规\allowbreak 模\allowbreak 大\allowbreak 而\allowbreak 利\allowbreak 用\allowbreak 不\allowbreak 足\allowbreak 的\allowbreak 资\allowbreak 源\allowbreak ，\allowbreak 将\allowbreak 其\allowbreak 转\allowbreak 化\allowbreak 为\allowbreak 可\allowbreak 分\allowbreak 析\allowbreak 记\allowbreak 录\allowbreak 至\allowbreak 今\allowbreak 仍\allowbreak 依\allowbreak 赖\allowbreak 大\allowbreak 量\allowbreak 人\allowbreak 工\allowbreak 。\allowbreak 本\allowbreak 文\allowbreak 提\allowbreak 出} CAIEMS------{\cjkfont 一\allowbreak 条\allowbreak 探\allowbreak 索\allowbreak 性} LLM {\cjkfont 数\allowbreak 据\allowbreak 工\allowbreak 程\allowbreak 管\allowbreak 道\allowbreak ，\allowbreak 将\allowbreak 公\allowbreak 开\allowbreak 可\allowbreak 获\allowbreak 取\allowbreak 的\allowbreak 中\allowbreak 文\allowbreak 施\allowbreak 工\allowbreak 事\allowbreak 故\allowbreak 叙\allowbreak 述\allowbreak 加\allowbreak 工\allowbreak 为\allowbreak 可\allowbreak 查\allowbreak 询\allowbreak 的\allowbreak 数\allowbreak 据\allowbreak 库\allowbreak 原\allowbreak 型\allowbreak ：}2,623 {\cjkfont 条\allowbreak 记\allowbreak 录\allowbreak 、}33 {\cjkfont 个\allowbreak 字\allowbreak 段\allowbreak 、}16,323 {\cjkfont 条\allowbreak 致\allowbreak 因\allowbreak 标\allowbreak 注\allowbreak 。\allowbreak 标\allowbreak 注\allowbreak 采\allowbreak 用\allowbreak 三\allowbreak 层\allowbreak 证\allowbreak 据\allowbreak 分\allowbreak 级\allowbreak 方\allowbreak 案\allowbreak ：\allowbreak 逐\allowbreak 字\allowbreak 证\allowbreak 据\allowbreak 编\allowbreak 码\allowbreak 、\allowbreak 语\allowbreak 义\allowbreak 归\allowbreak 纳\allowbreak 编\allowbreak 码\allowbreak 与\allowbreak 领\allowbreak 域\allowbreak 先\allowbreak 验\allowbreak 推\allowbreak 断\allowbreak 编\allowbreak 码\allowbreak ，\allowbreak 三\allowbreak 者\allowbreak 分\allowbreak 字\allowbreak 段\allowbreak 存\allowbreak 储\allowbreak 、\allowbreak 分\allowbreak 口\allowbreak 径\allowbreak 统\allowbreak 计\allowbreak 。\allowbreak 管\allowbreak 道\allowbreak 集\allowbreak 成\allowbreak 了\allowbreak 网\allowbreak 络\allowbreak 门\allowbreak 户\allowbreak 与\allowbreak 事\allowbreak 故\allowbreak 汇\allowbreak 编\allowbreak 数\allowbreak 字\allowbreak 化\allowbreak 构\allowbreak 成\allowbreak 的\allowbreak 双\allowbreak 轨\allowbreak 采\allowbreak 集\allowbreak 、}LLM {\cjkfont 文\allowbreak 本\allowbreak 分\allowbreak 类\allowbreak 、\allowbreak 结\allowbreak 构\allowbreak 化\allowbreak 字\allowbreak 段\allowbreak 抽\allowbreak 取\allowbreak 、\allowbreak 基\allowbreak 于} CAICS {\cjkfont 操\allowbreak 作\allowbreak 性\allowbreak 编\allowbreak 码\allowbreak 方\allowbreak 案\allowbreak （}5 {\cjkfont 层} 19 {\cjkfont 类\allowbreak ）\allowbreak 的\allowbreak 致\allowbreak 因\allowbreak 编\allowbreak 码\allowbreak 、}MECE {\cjkfont 字\allowbreak 段\allowbreak 标\allowbreak 准\allowbreak 化\allowbreak 与\allowbreak 交\allowbreak 互\allowbreak 式} Web/REST {\cjkfont 界\allowbreak 面\allowbreak 。\allowbreak 三\allowbreak 项\allowbreak 支\allowbreak 撑\allowbreak 实\allowbreak 验\allowbreak 均\allowbreak 以\allowbreak 明\allowbreak 确\allowbreak 限\allowbreak 定\allowbreak 的\allowbreak 口\allowbreak 径\allowbreak 报\allowbreak 告\allowbreak 。\allowbreak 其\allowbreak 一\allowbreak ，}4 {\cjkfont 个} LLM {\cjkfont 在} 200 {\cjkfont 条\allowbreak 记\allowbreak 录\allowbreak 上\allowbreak 的\allowbreak 约\allowbreak 束\allowbreak 性\allowbreak 比\allowbreak 较\allowbreak ：\allowbreak 所\allowbreak 选\allowbreak 模\allowbreak 型\allowbreak 基\allowbreak 础\allowbreak 字\allowbreak 段\allowbreak 准\allowbreak 确\allowbreak 率} 91.2\%{\cjkfont ，\allowbreak 且\allowbreak 为\allowbreak 候\allowbreak 选\allowbreak 模\allowbreak 型\allowbreak 中\allowbreak 成\allowbreak 本\allowbreak 最\allowbreak 低\allowbreak 者\allowbreak ；\allowbreak 需\allowbreak 说\allowbreak 明\allowbreak 的\allowbreak 是\allowbreak 各\allowbreak 模\allowbreak 型\allowbreak 运\allowbreak 行\allowbreak 条\allowbreak 件\allowbreak 并\allowbreak 不\allowbreak 对\allowbreak 等\allowbreak （\allowbreak 温\allowbreak 度\allowbreak 约\allowbreak 束\allowbreak 与\allowbreak 输\allowbreak 出\allowbreak 预\allowbreak 算\allowbreak 不\allowbreak 同\allowbreak ）\allowbreak ，\allowbreak 因\allowbreak 此\allowbreak 这\allowbreak 是\allowbreak 面\allowbreak 向\allowbreak 选\allowbreak 型\allowbreak 的\allowbreak 系\allowbreak 统\allowbreak 级\allowbreak 比\allowbreak 较\allowbreak ，\allowbreak 而\allowbreak 非\allowbreak 一\allowbreak 般\allowbreak 性\allowbreak 能\allowbreak 力\allowbreak 排\allowbreak 名\allowbreak 。\allowbreak 其\allowbreak 二\allowbreak ，}100 {\cjkfont 条\allowbreak 记\allowbreak 录\allowbreak 、\allowbreak 三\allowbreak 专\allowbreak 家\allowbreak 多\allowbreak 数\allowbreak 投\allowbreak 票\allowbreak 真\allowbreak 值\allowbreak 上\allowbreak 的\allowbreak 六\allowbreak 策\allowbreak 略\allowbreak 提\allowbreak 示\allowbreak 对\allowbreak 比\allowbreak （\allowbreak 少\allowbreak 样\allowbreak 本\allowbreak 最\allowbreak 优\allowbreak ，}F1=0.669{\cjkfont ）\allowbreak ，\allowbreak 期\allowbreak 间\allowbreak 观\allowbreak 察\allowbreak 到\allowbreak 一\allowbreak 种\allowbreak 提\allowbreak 示\allowbreak 措\allowbreak 辞\allowbreak 敏\allowbreak 感\allowbreak 性}------{\cjkfont 我\allowbreak 们\allowbreak 称\allowbreak 之\allowbreak 为\allowbreak 词\allowbreak 汇\allowbreak 门\allowbreak 控\allowbreak 效\allowbreak 应\allowbreak ：\allowbreak 疑\allowbreak 问\allowbreak 式\allowbreak 条\allowbreak 件\allowbreak 检\allowbreak 查\allowbreak 措\allowbreak 辞\allowbreak 抑\allowbreak 制\allowbreak 了\allowbreak 隐\allowbreak 性\allowbreak 类\allowbreak 别\allowbreak 的\allowbreak 同\allowbreak 义\allowbreak 证\allowbreak 据\allowbreak （}M1 {\cjkfont 召\allowbreak 回} 0.141{\cjkfont ）\allowbreak ，\allowbreak 同\allowbreak 一\allowbreak 设\allowbreak 置\allowbreak 下\allowbreak 两\allowbreak 步\allowbreak 判\allowbreak 断\allowbreak 改\allowbreak 写\allowbreak 将} M1 {\cjkfont 召\allowbreak 回\allowbreak 提\allowbreak 升\allowbreak 至} 0.941{\cjkfont 。\allowbreak 我\allowbreak 们\allowbreak 将\allowbreak 其\allowbreak 报\allowbreak 告\allowbreak 为\allowbreak 带\allowbreak 实\allowbreak 用\allowbreak 补\allowbreak 救\allowbreak 的\allowbreak 观\allowbreak 察\allowbreak 性\allowbreak 现\allowbreak 象\allowbreak ，\allowbreak 而\allowbreak 非\allowbreak 已\allowbreak 分\allowbreak 离\allowbreak 的\allowbreak 因\allowbreak 果\allowbreak 机\allowbreak 制\allowbreak 。\allowbreak 其\allowbreak 三\allowbreak ，\allowbreak 协\allowbreak 议\allowbreak 研\allowbreak 究\allowbreak （}n=289{\cjkfont ）\allowbreak ：}LLM {\cjkfont 与\allowbreak 一\allowbreak 名\allowbreak 受\allowbreak 训\allowbreak 标\allowbreak 注\allowbreak 者\allowbreak 在\allowbreak 六\allowbreak 项\allowbreak 基\allowbreak 础\allowbreak 字\allowbreak 段\allowbreak 上\allowbreak 达\allowbreak 到\allowbreak 高\allowbreak 一\allowbreak 致\allowbreak 性\allowbreak （\allowbreak 分\allowbreak 类\allowbreak 字\allowbreak 段} Cohen's kappa 0.91--0.97{\cjkfont ；\allowbreak 计\allowbreak 数\allowbreak 字\allowbreak 段} ICC 0.985{\cjkfont ；\allowbreak 日\allowbreak 期\allowbreak 年}-{\cjkfont 月\allowbreak 匹\allowbreak 配\allowbreak 率} 99.7\%{\cjkfont ）\allowbreak 。\allowbreak 全\allowbreak 文\allowbreak 对\allowbreak 结\allowbreak 果\allowbreak 的\allowbreak 边\allowbreak 界\allowbreak 作\allowbreak 显\allowbreak 式\allowbreak 声\allowbreak 明\allowbreak ：\allowbreak 银\allowbreak 标\allowbreak 准\allowbreak 部\allowbreak 分\allowbreak 由\allowbreak 模\allowbreak 型\allowbreak 输\allowbreak 出\allowbreak 构\allowbreak 建\allowbreak ，\allowbreak 提\allowbreak 示\allowbreak 开\allowbreak 发\allowbreak 与\allowbreak 最\allowbreak 终\allowbreak 测\allowbreak 试\allowbreak 共\allowbreak 用\allowbreak 同\allowbreak 一\allowbreak 样\allowbreak 本\allowbreak ，\allowbreak 数\allowbreak 据\allowbreak 库\allowbreak 覆\allowbreak 盖\allowbreak 的\allowbreak 是\allowbreak 公\allowbreak 开\allowbreak 可\allowbreak 检\allowbreak 索\allowbreak 的\allowbreak 事\allowbreak 故\allowbreak 报\allowbreak 道\allowbreak 而\allowbreak 非\allowbreak 全\allowbreak 部\allowbreak 施\allowbreak 工\allowbreak 事\allowbreak 故\allowbreak 的\allowbreak 代\allowbreak 表\allowbreak 性\allowbreak 样\allowbreak 本\allowbreak 。}

\textbf{{\cjkfont 关\allowbreak 键\allowbreak 词}}{\cjkfont ：\allowbreak 大\allowbreak 语\allowbreak 言\allowbreak 模\allowbreak 型\allowbreak ；\allowbreak 建\allowbreak 筑\allowbreak 施\allowbreak 工\allowbreak 安\allowbreak 全\allowbreak ；\allowbreak 信\allowbreak 息\allowbreak 抽\allowbreak 取\allowbreak ；\allowbreak 提\allowbreak 示\allowbreak 工\allowbreak 程\allowbreak ；\allowbreak 事\allowbreak 故\allowbreak 数\allowbreak 据\allowbreak 库\allowbreak ；\allowbreak 致\allowbreak 因\allowbreak 标\allowbreak 注\allowbreak ；\allowbreak 标\allowbreak 注\allowbreak 者\allowbreak 一\allowbreak 致\allowbreak 性}

\section{{\cjkfont 引\allowbreak 言}}
{\cjkfont 建\allowbreak 筑\allowbreak 业\allowbreak 仍\allowbreak 是\allowbreak 全\allowbreak 球\allowbreak 最\allowbreak 危\allowbreak 险\allowbreak 的\allowbreak 行\allowbreak 业\allowbreak 之\allowbreak 一\allowbreak 。\allowbreak 每\allowbreak 起\allowbreak 施\allowbreak 工\allowbreak 亡\allowbreak 人\allowbreak 事\allowbreak 故\allowbreak 都\allowbreak 会\allowbreak 产\allowbreak 生\allowbreak 调\allowbreak 查\allowbreak 报\allowbreak 告\allowbreak 与\allowbreak 新\allowbreak 闻\allowbreak 报\allowbreak 道\allowbreak ，\allowbreak 记\allowbreak 录\allowbreak 发\allowbreak 生\allowbreak 了\allowbreak 什\allowbreak 么\allowbreak 、\allowbreak 如\allowbreak 何\allowbreak 发\allowbreak 生\allowbreak 、\allowbreak 为\allowbreak 何\allowbreak 发\allowbreak 生\allowbreak 。\allowbreak 这\allowbreak 些\allowbreak 文\allowbreak 本\allowbreak 承\allowbreak 载\allowbreak 着\allowbreak 数\allowbreak 十\allowbreak 年\allowbreak 间\allowbreak 积\allowbreak 累\allowbreak 的\allowbreak 致\allowbreak 因\allowbreak 信\allowbreak 息\allowbreak ，\allowbreak 却\allowbreak 以\allowbreak 非\allowbreak 结\allowbreak 构\allowbreak 化\allowbreak 自\allowbreak 然\allowbreak 语\allowbreak 言\allowbreak 存\allowbreak 在\allowbreak ，\allowbreak 难\allowbreak 以\allowbreak 系\allowbreak 统\allowbreak 分\allowbreak 析\allowbreak 。\allowbreak 在\allowbreak 中\allowbreak 国\allowbreak ，\allowbreak 事\allowbreak 故\allowbreak 叙\allowbreak 述\allowbreak 分\allowbreak 散\allowbreak 于\allowbreak 安\allowbreak 全\allowbreak 信\allowbreak 息\allowbreak 门\allowbreak 户\allowbreak 、\allowbreak 地\allowbreak 方\allowbreak 应\allowbreak 急\allowbreak 管\allowbreak 理\allowbreak 局\allowbreak 、\allowbreak 新\allowbreak 闻\allowbreak 媒\allowbreak 体\allowbreak 与\allowbreak 公\allowbreak 开\allowbreak 事\allowbreak 故\allowbreak 汇\allowbreak 编\allowbreak ，\allowbreak 人\allowbreak 工\allowbreak 汇\allowbreak 编\allowbreak 成\allowbreak 本\allowbreak 高\allowbreak 昂\allowbreak ，\allowbreak 自\allowbreak 动\allowbreak 汇\allowbreak 编\allowbreak 则\allowbreak 是\allowbreak 一\allowbreak 个\allowbreak 反\allowbreak 复\allowbreak 出\allowbreak 现\allowbreak 的\allowbreak 工\allowbreak 程\allowbreak 挑\allowbreak 战\allowbreak （}Kong and Ahn, 2024{\cjkfont ）\allowbreak 。}

{\cjkfont 针\allowbreak 对\allowbreak 这\allowbreak 一\allowbreak 转\allowbreak 化\allowbreak 问\allowbreak 题\allowbreak 的\allowbreak 既\allowbreak 有\allowbreak 计\allowbreak 算\allowbreak 路\allowbreak 径\allowbreak 可\allowbreak 分\allowbreak 为\allowbreak 三\allowbreak 类\allowbreak 。}BERT{\cjkfont （}Devlin et al., 2019{\cjkfont ）\allowbreak 及\allowbreak 其\allowbreak 蒸\allowbreak 馏\allowbreak 变\allowbreak 体\allowbreak 等\allowbreak 监\allowbreak 督\allowbreak 分\allowbreak 类\allowbreak 器\allowbreak 依\allowbreak 赖\allowbreak 标\allowbreak 注\allowbreak 语\allowbreak 料\allowbreak ，\allowbreak 而\allowbreak 该\allowbreak 领\allowbreak 域\allowbreak 的\allowbreak 标\allowbreak 注\allowbreak 语\allowbreak 料\allowbreak 生\allowbreak 产\allowbreak 成\allowbreak 本\allowbreak 高\allowbreak 昂\allowbreak （}Togan et al., 2025{\cjkfont ）\allowbreak 。\allowbreak 基\allowbreak 于} LLM {\cjkfont 的\allowbreak 研\allowbreak 究\allowbreak 已\allowbreak 演\allowbreak 示\allowbreak 了\allowbreak 从} OSHA {\cjkfont 叙\allowbreak 述\allowbreak 中\allowbreak 提\allowbreak 取\allowbreak 事\allowbreak 故\allowbreak 属\allowbreak 性\allowbreak （}Ahmadi et al., 2024{\cjkfont ）\allowbreak 、\allowbreak 对\allowbreak 韩\allowbreak 国\allowbreak 坠\allowbreak 落\allowbreak 事\allowbreak 故\allowbreak 做} HFACS {\cjkfont 分\allowbreak 类\allowbreak （}Kim and Han, 2024{\cjkfont ）\allowbreak 、\allowbreak 从\allowbreak 施\allowbreak 工\allowbreak 报\allowbreak 告\allowbreak 构\allowbreak 建\allowbreak 知\allowbreak 识\allowbreak 图\allowbreak 谱\allowbreak （}Liu and Luo, 2025{\cjkfont ）\allowbreak 的\allowbreak 可\allowbreak 行\allowbreak 性\allowbreak 。\allowbreak 但\allowbreak 这\allowbreak 些\allowbreak 成\allowbreak 果\allowbreak 大\allowbreak 多\allowbreak 止\allowbreak 步\allowbreak 于\allowbreak 单\allowbreak 项\allowbreak 研\allowbreak 究\allowbreak ：\allowbreak 抽\allowbreak 取\allowbreak 服\allowbreak 务\allowbreak 于\allowbreak 某\allowbreak 一\allowbreak 次\allowbreak 分\allowbreak 析\allowbreak 或\allowbreak 静\allowbreak 态\allowbreak 图\allowbreak 谱\allowbreak ，\allowbreak 而\allowbreak 非\allowbreak 可\allowbreak 检\allowbreak 查\allowbreak 、\allowbreak 可\allowbreak 更\allowbreak 新\allowbreak 的\allowbreak 数\allowbreak 据\allowbreak 库\allowbreak 。}

{\cjkfont 本\allowbreak 文\allowbreak 报\allowbreak 告\allowbreak 的\allowbreak 是} CAIEMS------{\cjkfont 一\allowbreak 项\allowbreak 探\allowbreak 索\allowbreak 性\allowbreak 数\allowbreak 据\allowbreak 工\allowbreak 程\allowbreak 工\allowbreak 作\allowbreak ，\allowbreak 而\allowbreak 非\allowbreak 理\allowbreak 论\allowbreak 证\allowbreak 实\allowbreak 性\allowbreak 研\allowbreak 究\allowbreak 。\allowbreak 唯\allowbreak 一\allowbreak 的\allowbreak 组\allowbreak 织\allowbreak 性\allowbreak 问\allowbreak 题\allowbreak 是\allowbreak 实\allowbreak 践\allowbreak 性\allowbreak 的\allowbreak ：\allowbreak 一\allowbreak 个\allowbreak 小\allowbreak 团\allowbreak 队\allowbreak 能\allowbreak 否\allowbreak 利\allowbreak 用\allowbreak 通\allowbreak 用} LLM API {\cjkfont 与\allowbreak 有\allowbreak 限\allowbreak 算\allowbreak 力\allowbreak ，\allowbreak 将\allowbreak 公\allowbreak 开\allowbreak 可\allowbreak 获\allowbreak 取\allowbreak 的\allowbreak 中\allowbreak 文\allowbreak 事\allowbreak 故\allowbreak 叙\allowbreak 述\allowbreak 转\allowbreak 化\allowbreak 为\allowbreak 可\allowbreak 查\allowbreak 询\allowbreak 的\allowbreak 数\allowbreak 据\allowbreak 库\allowbreak ，\allowbreak 并\allowbreak 使\allowbreak 标\allowbreak 签\allowbreak 在\allowbreak 明\allowbreak 确\allowbreak 界\allowbreak 定\allowbreak 的\allowbreak 置\allowbreak 信\allowbreak 层\allowbreak 级\allowbreak 上\allowbreak 可\allowbreak 追\allowbreak 溯\allowbreak 到\allowbreak 原\allowbreak 文\allowbreak ？\allowbreak 我\allowbreak 们\allowbreak 报\allowbreak 告\allowbreak 管\allowbreak 道\allowbreak 产\allowbreak 出\allowbreak 了\allowbreak 什\allowbreak 么\allowbreak 、\allowbreak 我\allowbreak 们\allowbreak 对\allowbreak 其\allowbreak 可\allowbreak 靠\allowbreak 性\allowbreak 测\allowbreak 量\allowbreak 了\allowbreak 什\allowbreak 么\allowbreak 、\allowbreak 以\allowbreak 及\allowbreak 我\allowbreak 们\allowbreak 尚\allowbreak 不\allowbreak 能\allowbreak 确\allowbreak 立\allowbreak 什\allowbreak 么\allowbreak 。}

{\cjkfont 这\allowbreak 一\allowbreak 框\allowbreak 架\allowbreak 对\allowbreak 应\allowbreak 三\allowbreak 项\allowbreak 贡\allowbreak 献\allowbreak 。}(C1) {\cjkfont 数\allowbreak 据\allowbreak 工\allowbreak 程\allowbreak 管\allowbreak 道\allowbreak 与\allowbreak 数\allowbreak 据\allowbreak 库\allowbreak 原\allowbreak 型\allowbreak ：}2,623 {\cjkfont 条\allowbreak 记\allowbreak 录\allowbreak ，\allowbreak 采\allowbreak 集\allowbreak 自\allowbreak 网\allowbreak 络\allowbreak 爬\allowbreak 取\allowbreak 公\allowbreak 开\allowbreak 门\allowbreak 户\allowbreak 与\allowbreak 公\allowbreak 开\allowbreak 事\allowbreak 故\allowbreak 汇\allowbreak 编\allowbreak 数\allowbreak 字\allowbreak 化\allowbreak 构\allowbreak 成\allowbreak 的\allowbreak 双\allowbreak 轨\allowbreak 方\allowbreak 案\allowbreak ，\allowbreak 含} 33 {\cjkfont 字\allowbreak 段\allowbreak 与} 16,323 {\cjkfont 条\allowbreak 三\allowbreak 层\allowbreak 分\allowbreak 级\allowbreak 致\allowbreak 因\allowbreak 标\allowbreak 注\allowbreak ，\allowbreak 配} Web {\cjkfont 界\allowbreak 面\allowbreak 与} REST API{\cjkfont ；\allowbreak 数\allowbreak 据\allowbreak 与\allowbreak 代\allowbreak 码\allowbreak 于\allowbreak 论\allowbreak 文\allowbreak 发\allowbreak 表\allowbreak 后\allowbreak 开\allowbreak 放\allowbreak 检\allowbreak 查\allowbreak 。}(C2) CAICS {\cjkfont 操\allowbreak 作\allowbreak 性\allowbreak 编\allowbreak 码\allowbreak 方\allowbreak 案}------{\cjkfont 面\allowbreak 向\allowbreak 施\allowbreak 工\allowbreak 事\allowbreak 故\allowbreak 致\allowbreak 因\allowbreak 的} 5 {\cjkfont 层} 19 {\cjkfont 类\allowbreak 操\allowbreak 作\allowbreak 定\allowbreak 义\allowbreak ，\allowbreak 明\allowbreak 确\allowbreak 将\allowbreak 其\allowbreak 定\allowbreak 位\allowbreak 为\allowbreak 本\allowbreak 数\allowbreak 据\allowbreak 库\allowbreak 的\allowbreak 操\allowbreak 作\allowbreak 性\allowbreak 方\allowbreak 案\allowbreak ，\allowbreak 而\allowbreak 非\allowbreak 新\allowbreak 的\allowbreak 致\allowbreak 因\allowbreak 理\allowbreak 论\allowbreak 。}(C3) {\cjkfont 一\allowbreak 个\allowbreak 经\allowbreak 验\allowbreak 记\allowbreak 录\allowbreak 的\allowbreak 提\allowbreak 示\allowbreak 措\allowbreak 辞\allowbreak 敏\allowbreak 感\allowbreak 性}------{\cjkfont 词\allowbreak 汇\allowbreak 门\allowbreak 控\allowbreak 效\allowbreak 应}------{\cjkfont 在\allowbreak 提\allowbreak 示\allowbreak 开\allowbreak 发\allowbreak 过\allowbreak 程\allowbreak 中\allowbreak 观\allowbreak 察\allowbreak 到\allowbreak ，\allowbreak 并\allowbreak 配\allowbreak 有\allowbreak 带\allowbreak 边\allowbreak 界\allowbreak 声\allowbreak 明\allowbreak 的\allowbreak 两\allowbreak 步\allowbreak 判\allowbreak 断\allowbreak 补\allowbreak 救\allowbreak 。}

{\cjkfont 三\allowbreak 项\allowbreak 支\allowbreak 撑\allowbreak 实\allowbreak 验\allowbreak 锚\allowbreak 定\allowbreak 可\allowbreak 靠\allowbreak 性\allowbreak 讨\allowbreak 论\allowbreak 。\allowbreak 实\allowbreak 验\allowbreak 一\allowbreak 在} 200 {\cjkfont 条\allowbreak 记\allowbreak 录\allowbreak 上\allowbreak 比\allowbreak 较} 4 {\cjkfont 个} LLM {\cjkfont 以\allowbreak 支\allowbreak 持\allowbreak 骨\allowbreak 干\allowbreak 选\allowbreak 型\allowbreak ，\allowbreak 并\allowbreak 显\allowbreak 式\allowbreak 声\allowbreak 明\allowbreak 运\allowbreak 行\allowbreak 条\allowbreak 件\allowbreak 不\allowbreak 对\allowbreak 等\allowbreak 的\allowbreak 保\allowbreak 留\allowbreak 。\allowbreak 实\allowbreak 验\allowbreak 二\allowbreak 在} 100 {\cjkfont 条\allowbreak 记\allowbreak 录\allowbreak 、\allowbreak 三\allowbreak 专\allowbreak 家\allowbreak 多\allowbreak 数\allowbreak 投\allowbreak 票\allowbreak 真\allowbreak 值\allowbreak 上\allowbreak 比\allowbreak 较\allowbreak 六\allowbreak 种\allowbreak 提\allowbreak 示\allowbreak 策\allowbreak 略\allowbreak ；\allowbreak 该\allowbreak 实\allowbreak 验\allowbreak 发\allowbreak 现\allowbreak 了\allowbreak 上\allowbreak 述\allowbreak 措\allowbreak 辞\allowbreak 敏\allowbreak 感\allowbreak 性\allowbreak 。\allowbreak 实\allowbreak 验\allowbreak 三\allowbreak 在} 289 {\cjkfont 条\allowbreak 记\allowbreak 录\allowbreak 上\allowbreak 测\allowbreak 量} LLM {\cjkfont 与\allowbreak 一\allowbreak 名\allowbreak 受\allowbreak 训\allowbreak 标\allowbreak 注\allowbreak 者\allowbreak 在\allowbreak 六\allowbreak 项\allowbreak 基\allowbreak 础\allowbreak 字\allowbreak 段\allowbreak 上\allowbreak 的\allowbreak 一\allowbreak 致\allowbreak 性\allowbreak 。\allowbreak 各\allowbreak 实\allowbreak 验\allowbreak 所\allowbreak 用\allowbreak 样\allowbreak 本\allowbreak 的\allowbreak 已\allowbreak 知\allowbreak 局\allowbreak 限}------{\cjkfont 银\allowbreak 标\allowbreak 准\allowbreak 部\allowbreak 分\allowbreak 由\allowbreak 模\allowbreak 型\allowbreak 产\allowbreak 出\allowbreak 、\allowbreak 策\allowbreak 略\allowbreak 样\allowbreak 本\allowbreak 同\allowbreak 时\allowbreak 承\allowbreak 担\allowbreak 开\allowbreak 发\allowbreak 与\allowbreak 最\allowbreak 终\allowbreak 报\allowbreak 告}------{\cjkfont 均\allowbreak 在\allowbreak 每\allowbreak 次\allowbreak 使\allowbreak 用\allowbreak 处\allowbreak 声\allowbreak 明\allowbreak ，\allowbreak 而\allowbreak 非\allowbreak 集\allowbreak 中\allowbreak 于\allowbreak 一\allowbreak 段\allowbreak 笼\allowbreak 统\allowbreak 的\allowbreak 保\allowbreak 留\allowbreak 说\allowbreak 明\allowbreak 。}

{\cjkfont 论\allowbreak 文\allowbreak 其\allowbreak 余\allowbreak 部\allowbreak 分\allowbreak 组\allowbreak 织\allowbreak 如\allowbreak 下\allowbreak ：\allowbreak 第} 2 {\cjkfont 节\allowbreak 综\allowbreak 述\allowbreak 相\allowbreak 关\allowbreak 工\allowbreak 作\allowbreak ；\allowbreak 第} 3 {\cjkfont 节\allowbreak 描\allowbreak 述\allowbreak 数\allowbreak 据\allowbreak 来\allowbreak 源\allowbreak 与\allowbreak 语\allowbreak 料\allowbreak 构\allowbreak 建\allowbreak 、\allowbreak 证\allowbreak 据\allowbreak 分\allowbreak 级\allowbreak 方\allowbreak 案\allowbreak 、}CAICS {\cjkfont 操\allowbreak 作\allowbreak 定\allowbreak 义\allowbreak 、\allowbreak 管\allowbreak 道\allowbreak 架\allowbreak 构\allowbreak 与\allowbreak 数\allowbreak 据\allowbreak 库\allowbreak ；\allowbreak 第} 4 {\cjkfont 节\allowbreak 按\allowbreak 声\allowbreak 明\allowbreak 口\allowbreak 径\allowbreak 报\allowbreak 告\allowbreak 三\allowbreak 项\allowbreak 实\allowbreak 验\allowbreak ；\allowbreak 第} 5 {\cjkfont 节\allowbreak 讨\allowbreak 论\allowbreak 意\allowbreak 义\allowbreak 与\allowbreak 局\allowbreak 限\allowbreak ；\allowbreak 第} 6 {\cjkfont 节\allowbreak 总\allowbreak 结\allowbreak 。}

\section{{\cjkfont 相\allowbreak 关\allowbreak 工\allowbreak 作}}
\subsection{{\cjkfont 施\allowbreak 工\allowbreak 事\allowbreak 故\allowbreak 文\allowbreak 档\allowbreak 的\allowbreak 文\allowbreak 本\allowbreak 分\allowbreak 析}}
{\cjkfont 施\allowbreak 工\allowbreak 安\allowbreak 全\allowbreak 文\allowbreak 档\allowbreak 的\allowbreak 计\allowbreak 算\allowbreak 分\allowbreak 析\allowbreak 早\allowbreak 于} LLM{\cjkfont 。\allowbreak 对\allowbreak 人\allowbreak 工\allowbreak 编\allowbreak 码\allowbreak 记\allowbreak 录\allowbreak 做\allowbreak 关\allowbreak 联\allowbreak 规\allowbreak 则\allowbreak 挖\allowbreak 掘\allowbreak 识\allowbreak 别\allowbreak 了\allowbreak 事\allowbreak 故\allowbreak 类\allowbreak 型\allowbreak 与\allowbreak 致\allowbreak 因\allowbreak 因\allowbreak 素\allowbreak 间\allowbreak 的\allowbreak 可\allowbreak 解\allowbreak 释\allowbreak 关\allowbreak 联\allowbreak （}Rafindadi et al., 2023{\cjkfont ）\allowbreak ；\allowbreak 对} OSHA {\cjkfont 屋\allowbreak 面\allowbreak 工\allowbreak 人\allowbreak 记\allowbreak 录\allowbreak 做\allowbreak 决\allowbreak 策\allowbreak 树\allowbreak 建\allowbreak 模\allowbreak 产\allowbreak 出\allowbreak 了\allowbreak 伤\allowbreak 害\allowbreak 预\allowbreak 测\allowbreak 规\allowbreak 则\allowbreak （}Onuchukwu et al., 2024{\cjkfont ）\allowbreak 。}Ayhan and Tokdemir (2019) {\cjkfont 在\allowbreak 专\allowbreak 家\allowbreak 编\allowbreak 码\allowbreak 事\allowbreak 件\allowbreak 集\allowbreak 上\allowbreak 结\allowbreak 合\allowbreak 回\allowbreak 归\allowbreak 与\allowbreak 模\allowbreak 糊\allowbreak 逻\allowbreak 辑\allowbreak 。\allowbreak 这\allowbreak 些\allowbreak 研\allowbreak 究\allowbreak 的\allowbreak 共\allowbreak 同\allowbreak 约\allowbreak 束\allowbreak 是\allowbreak 循\allowbreak 环\allowbreak 性\allowbreak 的\allowbreak ：\allowbreak 分\allowbreak 析\allowbreak 质\allowbreak 量\allowbreak 受\allowbreak 制\allowbreak 于\allowbreak 结\allowbreak 构\allowbreak 化\allowbreak 输\allowbreak 入\allowbreak ，\allowbreak 而\allowbreak 结\allowbreak 构\allowbreak 化\allowbreak 输\allowbreak 入\allowbreak 本\allowbreak 身\allowbreak 产\allowbreak 自\allowbreak 昂\allowbreak 贵\allowbreak 的\allowbreak 人\allowbreak 工\allowbreak 编\allowbreak 码\allowbreak 。}Togan et al. (2025) {\cjkfont 直\allowbreak 面\allowbreak 转\allowbreak 化\allowbreak 环\allowbreak 节\allowbreak ，\allowbreak 在} 5,224 {\cjkfont 条\allowbreak 施\allowbreak 工\allowbreak 事\allowbreak 故\allowbreak 案\allowbreak 例\allowbreak 上\allowbreak 对\allowbreak 比} TF-IDF{\cjkfont 、}BERT{\cjkfont 、}RoBERTa{\cjkfont 、}XLNet {\cjkfont 与} DistilBERT {\cjkfont 的\allowbreak 安\allowbreak 全\allowbreak 属\allowbreak 性\allowbreak 提\allowbreak 取\allowbreak ，\allowbreak 发\allowbreak 现} DistilBERT {\cjkfont 以\allowbreak 约\allowbreak 一\allowbreak 半\allowbreak 的\allowbreak 计\allowbreak 算\allowbreak 成\allowbreak 本\allowbreak 取\allowbreak 得\allowbreak 相\allowbreak 近\allowbreak 性\allowbreak 能\allowbreak 。\allowbreak 其\allowbreak 结\allowbreak 果\allowbreak 是\allowbreak 监\allowbreak 督\allowbreak 范\allowbreak 式\allowbreak 的\allowbreak 典\allowbreak 型\allowbreak 写\allowbreak 照\allowbreak ：\allowbreak 可\allowbreak 用\allowbreak 的\allowbreak 准\allowbreak 确\allowbreak 率\allowbreak （\allowbreak 等\allowbreak 级\allowbreak 分\allowbreak 类} F1=0.72{\cjkfont ）\allowbreak 以\allowbreak 大\allowbreak 规\allowbreak 模\allowbreak 标\allowbreak 注\allowbreak 语\allowbreak 料\allowbreak 为\allowbreak 代\allowbreak 价\allowbreak 。}

\subsection{LLM {\cjkfont 事\allowbreak 故\allowbreak 信\allowbreak 息\allowbreak 抽\allowbreak 取}}
LLM {\cjkfont 通\allowbreak 过\allowbreak 无\allowbreak 需\allowbreak 任\allowbreak 务\allowbreak 特\allowbreak 定\allowbreak 训\allowbreak 练\allowbreak 、\allowbreak 直\allowbreak 接\allowbreak 从\allowbreak 文\allowbreak 本\allowbreak 推\allowbreak 断\allowbreak 结\allowbreak 构\allowbreak 改\allowbreak 变\allowbreak 了\allowbreak 瓶\allowbreak 颈\allowbreak 位\allowbreak 置\allowbreak （}Brown et al., 2020; OpenAI, 2023{\cjkfont ）\allowbreak 。}Ahmadi et al. (2024) {\cjkfont 评\allowbreak 测} GPT-3.5{\cjkfont 、}GPT-4.0{\cjkfont 、}Gemini Pro {\cjkfont 与} LLaMA 3.1 {\cjkfont 对} OSHA {\cjkfont 施\allowbreak 工\allowbreak 报\allowbreak 告\allowbreak 的\allowbreak 分\allowbreak 类\allowbreak 。}Zhang, Hou, et al. (2025a) {\cjkfont 构\allowbreak 建\allowbreak 了\allowbreak 基\allowbreak 于} GPT {\cjkfont 的\allowbreak 地\allowbreak 铁\allowbreak 施\allowbreak 工\allowbreak 报\allowbreak 告\allowbreak 分\allowbreak 析\allowbreak 对\allowbreak 话\allowbreak 系\allowbreak 统} AIR Agent{\cjkfont 。}Liu, Ding, et al. (2025b) {\cjkfont 将\allowbreak 零\allowbreak 样\allowbreak 本} GPT-4o {\cjkfont 抽\allowbreak 取\allowbreak 应\allowbreak 用\allowbreak 于} 1,097 {\cjkfont 条\allowbreak 高\allowbreak 处\allowbreak 坠\allowbreak 落\allowbreak 报\allowbreak 告\allowbreak ，\allowbreak 实\allowbreak 体\allowbreak 识\allowbreak 别} F1=0.94{\cjkfont 。\allowbreak 相\allowbreak 邻\allowbreak 领\allowbreak 域\allowbreak 中\allowbreak ，}Huang et al. (2026) {\cjkfont 为\allowbreak 海\allowbreak 事\allowbreak 事\allowbreak 故\allowbreak 构\allowbreak 建\allowbreak 了\allowbreak 带\allowbreak 质\allowbreak 量\allowbreak 评\allowbreak 估\allowbreak 框\allowbreak 架\allowbreak 的} LLM {\cjkfont 知\allowbreak 识\allowbreak 图\allowbreak 谱\allowbreak 管\allowbreak 道\allowbreak ，}Jiang et al. (2026) {\cjkfont 构\allowbreak 建\allowbreak 了\allowbreak 带\allowbreak 本\allowbreak 体\allowbreak 建\allowbreak 模\allowbreak 与\allowbreak 交\allowbreak 互\allowbreak 可\allowbreak 视\allowbreak 化\allowbreak 的\allowbreak 煤\allowbreak 矿\allowbreak 事\allowbreak 故\allowbreak 知\allowbreak 识\allowbreak 图\allowbreak 谱\allowbreak 。\allowbreak 这\allowbreak 些\allowbreak 系\allowbreak 统\allowbreak 已\allowbreak 在\allowbreak 不\allowbreak 同\allowbreak 程\allowbreak 度\allowbreak 上\allowbreak 组\allowbreak 合\allowbreak 了\allowbreak 采\allowbreak 集\allowbreak 、}LLM {\cjkfont 抽\allowbreak 取\allowbreak 与\allowbreak 图\allowbreak 谱\allowbreak 存\allowbreak 储\allowbreak ；\allowbreak 差\allowbreak 异\allowbreak 在\allowbreak 于\allowbreak 数\allowbreak 据\allowbreak 质\allowbreak 量\allowbreak 报\allowbreak 告\allowbreak 的\allowbreak 深\allowbreak 度\allowbreak 、\allowbreak 是\allowbreak 否\allowbreak 存\allowbreak 在\allowbreak 人\allowbreak 工\allowbreak 验\allowbreak 证\allowbreak 的\allowbreak 参\allowbreak 考\allowbreak 标\allowbreak 准\allowbreak 、\allowbreak 以\allowbreak 及\allowbreak 产\allowbreak 物\allowbreak 是\allowbreak 否\allowbreak 作\allowbreak 为\allowbreak 可\allowbreak 更\allowbreak 新\allowbreak 服\allowbreak 务\allowbreak 维\allowbreak 护\allowbreak 。}

\subsection{{\cjkfont 基\allowbreak 于} HFACS {\cjkfont 的} LLM {\cjkfont 致\allowbreak 因\allowbreak 分\allowbreak 类}}
{\cjkfont 人\allowbreak 因\allowbreak 分\allowbreak 析\allowbreak 与\allowbreak 分\allowbreak 类\allowbreak 系\allowbreak 统\allowbreak （}Shappell and Wiegmann, 2000{\cjkfont ）\allowbreak 主\allowbreak 导\allowbreak 高\allowbreak 风\allowbreak 险\allowbreak 行\allowbreak 业\allowbreak 的\allowbreak 致\allowbreak 因\allowbreak 分\allowbreak 类\allowbreak 。\allowbreak 思\allowbreak 维\allowbreak 链\allowbreak 提\allowbreak 示\allowbreak （}Wei et al., 2022{\cjkfont ）\allowbreak 及\allowbreak 其\allowbreak 变\allowbreak 体}------{\cjkfont 零\allowbreak 样\allowbreak 本\allowbreak 推\allowbreak 理\allowbreak （}Kojima et al., 2022{\cjkfont ）\allowbreak 、\allowbreak 自\allowbreak 洽\allowbreak 解\allowbreak 码\allowbreak （}Wang et al., 2023{\cjkfont ）\allowbreak 与\allowbreak 从\allowbreak 易\allowbreak 到\allowbreak 难\allowbreak 分\allowbreak 解\allowbreak （}Zhou et al., 2023{\cjkfont ）}------{\cjkfont 构\allowbreak 成\allowbreak 基\allowbreak 于\allowbreak 提\allowbreak 示\allowbreak 的\allowbreak 分\allowbreak 类\allowbreak 方\allowbreak 法\allowbreak 的\allowbreak 方\allowbreak 法\allowbreak 论\allowbreak 基\allowbreak 座\allowbreak 。}Liu et al. (2025a) {\cjkfont 提\allowbreak 出} HFACS-CoT{\cjkfont ，\allowbreak 在} 113 {\cjkfont 条\allowbreak 通\allowbreak 用\allowbreak 航\allowbreak 空\allowbreak 叙\allowbreak 述\allowbreak 上\allowbreak 以} GPT-4o {\cjkfont 比\allowbreak 较\allowbreak 六\allowbreak 种\allowbreak 提\allowbreak 示\allowbreak 策\allowbreak 略\allowbreak ，\allowbreak 发\allowbreak 现\allowbreak 领\allowbreak 域} CoT {\cjkfont 更\allowbreak 优\allowbreak （\allowbreak 召\allowbreak 回} 74.9\%{\cjkfont 、\allowbreak 精\allowbreak 确} 46.2\%{\cjkfont ）\allowbreak 。}Kim and Han (2024) {\cjkfont 将} HFACS {\cjkfont 与} Claude 3.5 Sonnet {\cjkfont 结\allowbreak 合\allowbreak 分\allowbreak 析} 1,816 {\cjkfont 条\allowbreak 韩\allowbreak 国\allowbreak 施\allowbreak 工\allowbreak 坠\allowbreak 落\allowbreak 案\allowbreak 例\allowbreak 。}Yan et al. (2025) {\cjkfont 跨\allowbreak 七\allowbreak 个} LLM {\cjkfont 应\allowbreak 用} HFACS {\cjkfont 引\allowbreak 导\allowbreak 提\allowbreak 示\allowbreak 。}Ahmadi et al. (2025) {\cjkfont 探\allowbreak 索\allowbreak 了} Llama-3.1 8B {\cjkfont 的} GRPO {\cjkfont 微\allowbreak 调\allowbreak 路\allowbreak 线\allowbreak 。\allowbreak 纵\allowbreak 观\allowbreak 该\allowbreak 文\allowbreak 献\allowbreak ，\allowbreak 策\allowbreak 略\allowbreak 排\allowbreak 名\allowbreak 均\allowbreak 在\allowbreak 单\allowbreak 一\allowbreak 模\allowbreak 型}-{\cjkfont 任\allowbreak 务\allowbreak 组\allowbreak 合\allowbreak 内\allowbreak 建\allowbreak 立\allowbreak ，\allowbreak 提\allowbreak 示\allowbreak 设\allowbreak 计\allowbreak 被\allowbreak 报\allowbreak 告\allowbreak 为}"{\cjkfont 可\allowbreak 用\allowbreak 配\allowbreak 置}"{\cjkfont ，\allowbreak 鲜\allowbreak 有\allowbreak 系\allowbreak 统\allowbreak 考\allowbreak 察\allowbreak 嵌\allowbreak 入\allowbreak 知\allowbreak 识\allowbreak 的\allowbreak 措\allowbreak 辞\allowbreak 如\allowbreak 何\allowbreak 约\allowbreak 束\allowbreak 模\allowbreak 型\allowbreak 行\allowbreak 为\allowbreak 。}

\subsection{{\cjkfont 知\allowbreak 识\allowbreak 图\allowbreak 谱\allowbreak 与\allowbreak 事\allowbreak 故\allowbreak 数\allowbreak 据\allowbreak 库}}
Kong and Ahn (2024) {\cjkfont 系\allowbreak 统\allowbreak 综\allowbreak 述\allowbreak 了\allowbreak 施\allowbreak 工\allowbreak 安\allowbreak 全\allowbreak 管\allowbreak 理\allowbreak 知\allowbreak 识\allowbreak 图\allowbreak 谱\allowbreak ，\allowbreak 并\allowbreak 指\allowbreak 出\allowbreak 研\allowbreak 究\allowbreak 原\allowbreak 型\allowbreak 与\allowbreak 运\allowbreak 营\allowbreak 平\allowbreak 台\allowbreak 之\allowbreak 间\allowbreak 的\allowbreak 差\allowbreak 距\allowbreak 。}Liu and Luo (2025) {\cjkfont 以\allowbreak 迭\allowbreak 代} CoT {\cjkfont 提\allowbreak 示\allowbreak 生\allowbreak 成\allowbreak 施\allowbreak 工\allowbreak 事\allowbreak 故\allowbreak 知\allowbreak 识\allowbreak 图\allowbreak 谱\allowbreak 。}Yang and Park (2025) {\cjkfont 构\allowbreak 建\allowbreak 了\allowbreak 道\allowbreak 路\allowbreak 事\allowbreak 故\allowbreak 新\allowbreak 闻\allowbreak 知\allowbreak 识\allowbreak 图\allowbreak 谱\allowbreak 。}Zhang, Wei, et al. (2025b) {\cjkfont 开\allowbreak 发\allowbreak 了\allowbreak 面\allowbreak 向\allowbreak 流\allowbreak 程\allowbreak 工\allowbreak 业\allowbreak 事\allowbreak 故\allowbreak 知\allowbreak 识\allowbreak 的\allowbreak 图\allowbreak 检\allowbreak 索\allowbreak 增\allowbreak 强\allowbreak 问\allowbreak 答\allowbreak 系\allowbreak 统} GRAG-ProSafe{\cjkfont 。}Shi et al. (2025) {\cjkfont 将\allowbreak 故\allowbreak 障\allowbreak 树\allowbreak 、\allowbreak 贝\allowbreak 叶\allowbreak 斯\allowbreak 网\allowbreak 络\allowbreak 与\allowbreak 本\allowbreak 体\allowbreak 推\allowbreak 理\allowbreak 结\allowbreak 合\allowbreak 用\allowbreak 于\allowbreak 基\allowbreak 坑\allowbreak 坍\allowbreak 塌\allowbreak 分\allowbreak 析\allowbreak 。}Xiao et al. (2026) {\cjkfont 将\allowbreak 知\allowbreak 识\allowbreak 图\allowbreak 谱\allowbreak 作\allowbreak 为} LLM {\cjkfont 人\allowbreak 因\allowbreak 可\allowbreak 靠\allowbreak 性\allowbreak 分\allowbreak 析\allowbreak 的\allowbreak 检\allowbreak 索\allowbreak 层\allowbreak 集\allowbreak 成\allowbreak 。\allowbreak 这\allowbreak 些\allowbreak 系\allowbreak 统\allowbreak 各\allowbreak 自\allowbreak 构\allowbreak 建\allowbreak 了\allowbreak 某\allowbreak 种\allowbreak 知\allowbreak 识\allowbreak 产\allowbreak 物\allowbreak ；\allowbreak 本\allowbreak 工\allowbreak 作\allowbreak 的\allowbreak 差\allowbreak 异\allowbreak 在\allowbreak 于\allowbreak 侧\allowbreak 重\allowbreak 而\allowbreak 非\allowbreak 类\allowbreak 别}------{\cjkfont 优\allowbreak 先\allowbreak 可\allowbreak 追\allowbreak 溯\allowbreak 的\allowbreak 分\allowbreak 级\allowbreak 证\allowbreak 据\allowbreak 标\allowbreak 注\allowbreak 与\allowbreak 可\allowbreak 更\allowbreak 新\allowbreak 的\allowbreak 查\allowbreak 询\allowbreak 接\allowbreak 口\allowbreak ，\allowbreak 并\allowbreak 以\allowbreak 测\allowbreak 量\allowbreak 的\allowbreak 边\allowbreak 界\allowbreak 报\allowbreak 告\allowbreak 标\allowbreak 注\allowbreak 的\allowbreak 可\allowbreak 靠\allowbreak 性\allowbreak 。}

\section{{\cjkfont 系\allowbreak 统}}
\subsection{{\cjkfont 问\allowbreak 题\allowbreak 形\allowbreak 式\allowbreak 化\allowbreak 与\allowbreak 证\allowbreak 据\allowbreak 分\allowbreak 级\allowbreak 标\allowbreak 注\allowbreak 方\allowbreak 案}}
{\cjkfont 令} $T = \{t_1, t_2, \dots, t_N\}$ {\cjkfont 表\allowbreak 示\allowbreak 非\allowbreak 结\allowbreak 构\allowbreak 化\allowbreak 事\allowbreak 故\allowbreak 叙\allowbreak 述\allowbreak 语\allowbreak 料\allowbreak （\allowbreak 本\allowbreak 语\allowbreak 料\allowbreak 中\allowbreak 位\allowbreak 数} 517 {\cjkfont 字\allowbreak 符\allowbreak ）\allowbreak 。\allowbreak 管\allowbreak 道\allowbreak 将\allowbreak 每\allowbreak 条} $t_i$ {\cjkfont 转\allowbreak 化\allowbreak 为\allowbreak 结\allowbreak 构\allowbreak 化\allowbreak 记\allowbreak 录} $r_i = (b_i, c_i)${\cjkfont ：\allowbreak 基\allowbreak 础\allowbreak 信\allowbreak 息\allowbreak 向\allowbreak 量} $b_i$ {\cjkfont 覆\allowbreak 盖} 33 {\cjkfont 个\allowbreak 预\allowbreak 定\allowbreak 义\allowbreak 字\allowbreak 段\allowbreak ；\allowbreak 致\allowbreak 因\allowbreak 标\allowbreak 注\allowbreak 集} $c_i = \{(x_{ij}, e_{ij}, \iota_{ij}, q_{ij})\}${\cjkfont ，\allowbreak 其\allowbreak 中} $x_{ij}$ {\cjkfont 为} CAICS {\cjkfont 方\allowbreak 案\allowbreak 中\allowbreak 的\allowbreak 致\allowbreak 因\allowbreak 编\allowbreak 码\allowbreak ，}$e_{ij}$ {\cjkfont 为\allowbreak 支\allowbreak 撑\allowbreak 文\allowbreak 本\allowbreak ，}$\iota_{ij}$ {\cjkfont 为\allowbreak 推\allowbreak 理\allowbreak 链\allowbreak ，}$q_{ij} \in \{\text{explicit}, \text{inferred}, \text{prior}\}$ {\cjkfont 为\allowbreak 证\allowbreak 据\allowbreak 层\allowbreak 级\allowbreak 。}

{\cjkfont 三\allowbreak 个\allowbreak 层\allowbreak 级\allowbreak 承\allowbreak 担\allowbreak 不\allowbreak 同\allowbreak 的\allowbreak 认\allowbreak 识\allowbreak 论\allowbreak 承\allowbreak 诺\allowbreak ，\allowbreak 数\allowbreak 据\allowbreak 库\allowbreak 分\allowbreak 开\allowbreak 存\allowbreak 储\allowbreak 、\allowbreak 分\allowbreak 开\allowbreak 报\allowbreak 告\allowbreak ：}

\begin{itemize}
\item \textbf{explicit{\cjkfont （\allowbreak 逐\allowbreak 字\allowbreak 证\allowbreak 据\allowbreak ）}}{\cjkfont ：\allowbreak 原\allowbreak 文\allowbreak 含\allowbreak 直\allowbreak 接\allowbreak 支\allowbreak 撑\allowbreak 该\allowbreak 编\allowbreak 码\allowbreak 的\allowbreak 表\allowbreak 述\allowbreak ，}$e_{ij}$ {\cjkfont 为\allowbreak 逐\allowbreak 字\allowbreak 引\allowbreak 文\allowbreak 。}Web {\cjkfont 界\allowbreak 面\allowbreak 中\allowbreak 仅\allowbreak 此\allowbreak 类\allowbreak 编\allowbreak 码\allowbreak 呈\allowbreak 现\allowbreak 为}"{\cjkfont 有\allowbreak 证\allowbreak 据\allowbreak 支\allowbreak 撑}"{\cjkfont 。}
\item \textbf{inferred{\cjkfont （\allowbreak 语\allowbreak 义\allowbreak 归\allowbreak 纳\allowbreak ）}}{\cjkfont ：\allowbreak 编\allowbreak 码\allowbreak 由\allowbreak 叙\allowbreak 述\allowbreak 经\allowbreak 同\allowbreak 义\allowbreak 改\allowbreak 写\allowbreak 归\allowbreak 纳\allowbreak 得\allowbreak 出\allowbreak （\allowbreak 如}"{\cjkfont 安\allowbreak 全\allowbreak 管\allowbreak 理\allowbreak 不\allowbreak 到\allowbreak 位}"{\cjkfont 在\allowbreak 未\allowbreak 引\allowbreak 用\allowbreak 任\allowbreak 何\allowbreak 固\allowbreak 定\allowbreak 关\allowbreak 键\allowbreak 词\allowbreak 的\allowbreak 情\allowbreak 况\allowbreak 下\allowbreak 支\allowbreak 撑} M1{\cjkfont ）\allowbreak 。\allowbreak 推\allowbreak 理\allowbreak 链\allowbreak 已\allowbreak 记\allowbreak 录\allowbreak 。}
\item \textbf{prior{\cjkfont （\allowbreak 领\allowbreak 域\allowbreak 先\allowbreak 验\allowbreak ）}}{\cjkfont ：\allowbreak 编\allowbreak 码\allowbreak 来\allowbreak 自\allowbreak 提\allowbreak 示\allowbreak 中\allowbreak 授\allowbreak 权\allowbreak 的\allowbreak 领\allowbreak 域\allowbreak 假\allowbreak 设\allowbreak （\allowbreak 如\allowbreak 亡\allowbreak 人\allowbreak 事\allowbreak 故\allowbreak 视\allowbreak 为\allowbreak 隐\allowbreak 含\allowbreak 监\allowbreak 管\allowbreak 缺\allowbreak 失\allowbreak ，\allowbreak 即\allowbreak 使\allowbreak 文\allowbreak 本\allowbreak 未\allowbreak 提\allowbreak ）\allowbreak 。\allowbreak 此\allowbreak 类\allowbreak 编\allowbreak 码\allowbreak 带\allowbreak 先\allowbreak 验\allowbreak 标\allowbreak 签\allowbreak 展\allowbreak 示\allowbreak ，\allowbreak 并\allowbreak 排\allowbreak 除\allowbreak 在}"{\cjkfont 有\allowbreak 证\allowbreak 据\allowbreak 支\allowbreak 撑}"{\cjkfont 统\allowbreak 计\allowbreak 之\allowbreak 外\allowbreak 。}
\end{itemize}
{\cjkfont 该\allowbreak 方\allowbreak 案\allowbreak 取\allowbreak 代\allowbreak 了\allowbreak 早\allowbreak 期\allowbreak 草\allowbreak 稿\allowbreak 中}"{\cjkfont 每\allowbreak 个\allowbreak 编\allowbreak 码\allowbreak 均\allowbreak 带\allowbreak 可\allowbreak 引\allowbreak 用\allowbreak 证\allowbreak 据\allowbreak 、\allowbreak 无\allowbreak 证\allowbreak 据\allowbreak 即\allowbreak 幻\allowbreak 觉}"{\cjkfont 的\allowbreak 更\allowbreak 强\allowbreak 声\allowbreak 明\allowbreak 。\allowbreak 分\allowbreak 层\allowbreak 方\allowbreak 案\allowbreak 才\allowbreak 是\allowbreak 对\allowbreak 生\allowbreak 产\allowbreak 提\allowbreak 示\allowbreak 实\allowbreak 际\allowbreak 行\allowbreak 为\allowbreak 的\allowbreak 准\allowbreak 确\allowbreak 描\allowbreak 述\allowbreak ：\allowbreak 直\allowbreak 接\allowbreak 抽\allowbreak 取\allowbreak 、\allowbreak 同\allowbreak 义\allowbreak 归\allowbreak 纳\allowbreak 与\allowbreak 领\allowbreak 域\allowbreak 先\allowbreak 验\allowbreak 推\allowbreak 断\allowbreak 对\allowbreak 下\allowbreak 游\allowbreak 分\allowbreak 析\allowbreak 都\allowbreak 有\allowbreak 用\allowbreak ，\allowbreak 但\allowbreak 混\allowbreak 为\allowbreak 一\allowbreak 谈\allowbreak 会\allowbreak 把\allowbreak 模\allowbreak 型\allowbreak 先\allowbreak 验\allowbreak 呈\allowbreak 现\allowbreak 为\allowbreak 经\allowbreak 验\allowbreak 发\allowbreak 现\allowbreak 。\allowbreak 全\allowbreak 库} 16,323 {\cjkfont 条\allowbreak 标\allowbreak 注\allowbreak 中\allowbreak ，}explicit 14,784 {\cjkfont 条\allowbreak （}90.6\%{\cjkfont ）\allowbreak 、}inferred 1,117 {\cjkfont 条\allowbreak （}6.8\%{\cjkfont ）\allowbreak 、}prior 422 {\cjkfont 条\allowbreak （}2.6\%{\cjkfont ）\allowbreak ；\allowbreak 本\allowbreak 文\allowbreak 报\allowbreak 告\allowbreak 总\allowbreak 量\allowbreak ，\allowbreak 并\allowbreak 在\allowbreak 相\allowbreak 关\allowbreak 处\allowbreak 给\allowbreak 出\allowbreak 分\allowbreak 层\allowbreak 细\allowbreak 目\allowbreak 。}

\subsection{{\cjkfont 数\allowbreak 据\allowbreak 来\allowbreak 源\allowbreak 与\allowbreak 语\allowbreak 料\allowbreak 构\allowbreak 建}}
{\cjkfont 本\allowbreak 文\allowbreak 报\allowbreak 告\allowbreak 的\allowbreak 每\allowbreak 一\allowbreak 项\allowbreak 结\allowbreak 果\allowbreak 都\allowbreak 受\allowbreak 语\allowbreak 料\allowbreak 内\allowbreak 容\allowbreak 的\allowbreak 约\allowbreak 束\allowbreak ，\allowbreak 而\allowbreak 该\allowbreak 语\allowbreak 料\allowbreak 是}\textbf{{\cjkfont 汇\allowbreak 集}}{\cjkfont 而\allowbreak 成\allowbreak 、\allowbreak 并\allowbreak 非\allowbreak 抽\allowbreak 样\allowbreak 得\allowbreak 来\allowbreak 。\allowbreak 本\allowbreak 节\allowbreak 说\allowbreak 明} 2,623 {\cjkfont 条\allowbreak 记\allowbreak 录\allowbreak 的\allowbreak 来\allowbreak 源\allowbreak ，\allowbreak 以\allowbreak 及\allowbreak 建\allowbreak 模\allowbreak 之\allowbreak 前\allowbreak 对\allowbreak 其\allowbreak 施\allowbreak 加\allowbreak 的\allowbreak 四\allowbreak 项\allowbreak 处\allowbreak 理}------{\cjkfont 采\allowbreak 集\allowbreak 、\allowbreak 清\allowbreak 洗\allowbreak 、\allowbreak 相\allowbreak 关\allowbreak 性\allowbreak 筛\allowbreak 选\allowbreak 与\allowbreak 整\allowbreak 合\allowbreak 。\allowbreak 这\allowbreak 些\allowbreak 流\allowbreak 程\allowbreak 细\allowbreak 化\allowbreak 了\allowbreak 管\allowbreak 道\allowbreak 架\allowbreak 构\allowbreak 的\allowbreak 阶\allowbreak 段} 1{\cjkfont （}3.4 {\cjkfont 节\allowbreak ）\allowbreak 及\allowbreak 其\allowbreak 原\allowbreak 本\allowbreak 仅\allowbreak 以\allowbreak 一\allowbreak 句\allowbreak 话\allowbreak 概\allowbreak 括\allowbreak 的\allowbreak 预\allowbreak 处\allowbreak 理\allowbreak 。\allowbreak 我\allowbreak 们\allowbreak 不\allowbreak 惜\allowbreak 篇\allowbreak 幅\allowbreak 加\allowbreak 以\allowbreak 说\allowbreak 明\allowbreak ，\allowbreak 是\allowbreak 因\allowbreak 为\allowbreak 在\allowbreak 端\allowbreak 到\allowbreak 端\allowbreak 数\allowbreak 据\allowbreak 工\allowbreak 程\allowbreak 管\allowbreak 道\allowbreak 中\allowbreak ，\allowbreak 输\allowbreak 入\allowbreak 的\allowbreak 性\allowbreak 质\allowbreak 会\allowbreak 传\allowbreak 播\allowbreak 至\allowbreak 每\allowbreak 一\allowbreak 项\allowbreak 下\allowbreak 游\allowbreak 统\allowbreak 计\allowbreak ：\allowbreak 读\allowbreak 者\allowbreak 若\allowbreak 不\allowbreak 了\allowbreak 解\allowbreak 采\allowbreak 集\allowbreak 了\allowbreak 什\allowbreak 么\allowbreak 、\allowbreak 丢\allowbreak 弃\allowbreak 了\allowbreak 什\allowbreak 么\allowbreak ，\allowbreak 便\allowbreak 无\allowbreak 从\allowbreak 判\allowbreak 断\allowbreak 第} 4 {\cjkfont 节\allowbreak 的\allowbreak 结\allowbreak 果\allowbreak 。}

\subsubsection{{\cjkfont 双\allowbreak 轨\allowbreak 语\allowbreak 料}}
{\cjkfont 语\allowbreak 料\allowbreak 由\allowbreak 两\allowbreak 条\allowbreak 结\allowbreak 构\allowbreak 不\allowbreak 同\allowbreak 的\allowbreak 采\allowbreak 集\allowbreak 轨\allowbreak 道\allowbreak 组\allowbreak 成\allowbreak 。}\textbf{{\cjkfont 网\allowbreak 络\allowbreak 采\allowbreak 集\allowbreak 轨}}{\cjkfont 贡\allowbreak 献} 1,728 {\cjkfont 条\allowbreak 记\allowbreak 录\allowbreak （}65.9\%{\cjkfont ）\allowbreak ，\allowbreak 由\allowbreak 专\allowbreak 用\allowbreak 爬\allowbreak 虫\allowbreak 自\allowbreak 公\allowbreak 开\allowbreak 门\allowbreak 户\allowbreak 检\allowbreak 索\allowbreak 获\allowbreak 得\allowbreak 。}\textbf{{\cjkfont 文\allowbreak 献\allowbreak 数\allowbreak 字\allowbreak 化\allowbreak 轨}}{\cjkfont 贡\allowbreak 献} 895 {\cjkfont 条\allowbreak 记\allowbreak 录\allowbreak （}34.1\%{\cjkfont ）\allowbreak ，\allowbreak 来\allowbreak 自\allowbreak 以} PDF {\cjkfont 形\allowbreak 式\allowbreak 发\allowbreak 行\allowbreak 的\allowbreak 公\allowbreak 开\allowbreak 事\allowbreak 故\allowbreak 汇\allowbreak 编\allowbreak 。\allowbreak 表} 1 {\cjkfont 报\allowbreak 告\allowbreak 两\allowbreak 条\allowbreak 轨\allowbreak 道\allowbreak 的\allowbreak 构\allowbreak 成\allowbreak 。}

\textbf{{\cjkfont 表} 1. {\cjkfont 按\allowbreak 来\allowbreak 源\allowbreak 轨\allowbreak 道\allowbreak 划\allowbreak 分\allowbreak 的\allowbreak 语\allowbreak 料\allowbreak 构\allowbreak 成}}

\begingroup\setlength{\tabcolsep}{2pt}\small
\begin{longtable}{>{\raggedright\arraybackslash}p{0.1003\textwidth}>{\raggedright\arraybackslash}p{0.3689\textwidth}>{\raggedright\arraybackslash}p{0.1836\textwidth}>{\raggedright\arraybackslash}p{0.1692\textwidth}>{\raggedright\arraybackslash}p{0.1287\textwidth}}
\hline
{\cjkfont 轨\allowbreak 道} & {\cjkfont 来\allowbreak 源} & {\cjkfont 记\allowbreak 录\allowbreak 数} & {\cjkfont 占\allowbreak 比} & {\cjkfont 事\allowbreak 故\allowbreak 年\allowbreak 代} \\
\hline\endfirsthead
\hline
{\cjkfont 轨\allowbreak 道} & {\cjkfont 来\allowbreak 源} & {\cjkfont 记\allowbreak 录\allowbreak 数} & {\cjkfont 占\allowbreak 比} & {\cjkfont 事\allowbreak 故\allowbreak 年\allowbreak 代} \\
\hline\endhead
{\cjkfont 网\allowbreak 络\allowbreak 采\allowbreak 集} & {\cjkfont 安\allowbreak 全\allowbreak 管\allowbreak 理\allowbreak 网\allowbreak （}safehoo.com{\cjkfont ）\allowbreak ，\allowbreak 安\allowbreak 全\allowbreak 新\allowbreak 闻\allowbreak 门\allowbreak 户} & 1,247 & 47.5\% & 2005--2026 \\
 & esafety{\cjkfont （}esafety.cn{\cjkfont ）\allowbreak ，\allowbreak 安\allowbreak 全\allowbreak 行\allowbreak 业\allowbreak 门\allowbreak 户} & 389 & 14.8\% & 2005--2026 \\
 & {\cjkfont 省\allowbreak 级\allowbreak 应\allowbreak 急\allowbreak 管\allowbreak 理\allowbreak 部\allowbreak 门\allowbreak （}12 {\cjkfont 个\allowbreak 有\allowbreak 产\allowbreak 出\allowbreak 的\allowbreak 源\allowbreak ）} & 89 & 3.4\% & 2005--2026 \\
 & {\cjkfont 应\allowbreak 急\allowbreak 管\allowbreak 理\allowbreak 部\allowbreak （}mem.gov.cn{\cjkfont ）} & 3 & 0.1\% & 2005--2026 \\
 & \textit{{\cjkfont 小\allowbreak 计}} & \textit{1,728} & \textit{65.9\%} &  \\
{\cjkfont 文\allowbreak 献\allowbreak 数\allowbreak 字\allowbreak 化} & {\cjkfont 樊\allowbreak 锡\allowbreak 仁\allowbreak 《\allowbreak 建\allowbreak 筑\allowbreak 施\allowbreak 工\allowbreak 重\allowbreak 大\allowbreak 伤\allowbreak 亡\allowbreak 事\allowbreak 故} 800 {\cjkfont 例\allowbreak 》} & 703 & 26.8\% & 1978--1990 \\
 & {\cjkfont 住\allowbreak 建\allowbreak 部\allowbreak 事\allowbreak 故\allowbreak 案\allowbreak 例\allowbreak 汇\allowbreak 编\allowbreak （}3 {\cjkfont 卷\allowbreak ）} & 192 & 7.3\% & 2005--2026 \\
 & \textit{{\cjkfont 小\allowbreak 计}} & \textit{895} & \textit{34.1\%} &  \\
\textbf{{\cjkfont 合\allowbreak 计}} &  & \textbf{2,623} & \textbf{100\%} &  \\
\hline
\end{longtable}\endgroup
{\cjkfont 两\allowbreak 条\allowbreak 轨\allowbreak 道\allowbreak 不\allowbreak 可\allowbreak 互\allowbreak 换\allowbreak ，\allowbreak 将\allowbreak 其\allowbreak 呈\allowbreak 现\allowbreak 为\allowbreak 单\allowbreak 一\allowbreak 同\allowbreak 质\allowbreak 来\allowbreak 源\allowbreak 会\allowbreak 曲\allowbreak 解\allowbreak 数\allowbreak 据\allowbreak 库\allowbreak 。\allowbreak 文\allowbreak 献\allowbreak 数\allowbreak 字\allowbreak 化\allowbreak 轨\allowbreak 承\allowbreak 载\allowbreak 语\allowbreak 料\allowbreak 的\allowbreak 历\allowbreak 史\allowbreak 纵\allowbreak 深\allowbreak ：\allowbreak 取\allowbreak 自\allowbreak 樊\allowbreak 锡\allowbreak 仁\allowbreak 汇\allowbreak 编\allowbreak 的} 703 {\cjkfont 条\allowbreak 记\allowbreak 录\allowbreak 是} 1978--1990 {\cjkfont 时\allowbreak 间\allowbreak 块\allowbreak 的\allowbreak 唯\allowbreak 一\allowbreak 来\allowbreak 源\allowbreak ，\allowbreak 无\allowbreak 任\allowbreak 何\allowbreak 网\allowbreak 络\allowbreak 来\allowbreak 源\allowbreak 贡\allowbreak 献\allowbreak 该\allowbreak 时\allowbreak 期\allowbreak 数\allowbreak 据\allowbreak 。\allowbreak 网\allowbreak 络\allowbreak 轨\allowbreak 提\allowbreak 供\allowbreak 当\allowbreak 代\allowbreak 记\allowbreak 录\allowbreak ，\allowbreak 包\allowbreak 括\allowbreak 全\allowbreak 部\allowbreak 政\allowbreak 府\allowbreak 调\allowbreak 查\allowbreak 报\allowbreak 告}------{\cjkfont 这\allowbreak 些\allowbreak 是\allowbreak 语\allowbreak 料\allowbreak 中\allowbreak 最\allowbreak 详\allowbreak 尽\allowbreak 、\allowbreak 最\allowbreak 权\allowbreak 威\allowbreak 的\allowbreak 文\allowbreak 档\allowbreak 。\allowbreak 两\allowbreak 条\allowbreak 轨\allowbreak 道\allowbreak 在\allowbreak 行\allowbreak 文\allowbreak 语\allowbreak 体\allowbreak 、\allowbreak 篇\allowbreak 幅\allowbreak 与\allowbreak 字\allowbreak 段\allowbreak 完\allowbreak 整\allowbreak 度\allowbreak 上\allowbreak 存\allowbreak 在\allowbreak 系\allowbreak 统\allowbreak 性\allowbreak 差\allowbreak 异\allowbreak ，\allowbreak 因\allowbreak 此\allowbreak 管\allowbreak 道\allowbreak 对\allowbreak 两\allowbreak 者\allowbreak 施\allowbreak 加\allowbreak 相\allowbreak 同\allowbreak 的\allowbreak 抽\allowbreak 取\allowbreak 与\allowbreak 编\allowbreak 码\allowbreak 提\allowbreak 示\allowbreak ，\allowbreak 但\allowbreak 在\allowbreak 报\allowbreak 告\allowbreak 总\allowbreak 量\allowbreak 统\allowbreak 计\allowbreak 的\allowbreak 同\allowbreak 时\allowbreak 报\allowbreak 告\allowbreak 来\allowbreak 源\allowbreak 构\allowbreak 成\allowbreak ，\allowbreak 并\allowbreak 在} Web {\cjkfont 界\allowbreak 面\allowbreak 中\allowbreak 将\allowbreak 来\allowbreak 源\allowbreak 轨\allowbreak 道\allowbreak 作\allowbreak 为\allowbreak 可\allowbreak 筛\allowbreak 选\allowbreak 属\allowbreak 性\allowbreak 开\allowbreak 放\allowbreak 。}

{\cjkfont 语\allowbreak 料\allowbreak 横\allowbreak 跨} 1978 {\cjkfont 至} 2026 {\cjkfont 年\allowbreak ，\allowbreak 且\allowbreak 存\allowbreak 在\allowbreak 明\allowbreak 显\allowbreak 断\allowbreak 代\allowbreak 。}2,623 {\cjkfont 条\allowbreak 记\allowbreak 录\allowbreak 中\allowbreak ，}698 {\cjkfont 条\allowbreak （}26.6\%{\cjkfont ）\allowbreak 落\allowbreak 在} 1978--1990 {\cjkfont 年\allowbreak ，}1,895 {\cjkfont 条\allowbreak （}72.2\%{\cjkfont ）\allowbreak 落\allowbreak 在} 2005--2026 {\cjkfont 年\allowbreak ，}30 {\cjkfont 条\allowbreak （}1.1\%{\cjkfont ）\allowbreak 无\allowbreak 可\allowbreak 辨\allowbreak 识\allowbreak 日\allowbreak 期\allowbreak 。}\textbf{1991--2004 {\cjkfont 年\allowbreak 完\allowbreak 全\allowbreak 缺\allowbreak 失\allowbreak 。}} {\cjkfont 该\allowbreak 断\allowbreak 代\allowbreak 反\allowbreak 映\allowbreak 的\allowbreak 是\allowbreak 数\allowbreak 字\allowbreak 来\allowbreak 源\allowbreak 的\allowbreak 可\allowbreak 得\allowbreak 性\allowbreak ，\allowbreak 而\allowbreak 非\allowbreak 施\allowbreak 工\allowbreak 安\allowbreak 全\allowbreak 本\allowbreak 身\allowbreak 的\allowbreak 断\allowbreak 裂\allowbreak ：\allowbreak 历\allowbreak 史\allowbreak 汇\allowbreak 编\allowbreak 覆\allowbreak 盖\allowbreak 改\allowbreak 革\allowbreak 开\allowbreak 放\allowbreak 初\allowbreak 期\allowbreak ，\allowbreak 而\allowbreak 事\allowbreak 故\allowbreak 调\allowbreak 查\allowbreak 报\allowbreak 告\allowbreak 的\allowbreak 持\allowbreak 续\allowbreak 性\allowbreak 公\allowbreak 共\allowbreak 网\allowbreak 络\allowbreak 发\allowbreak 布\allowbreak 始\allowbreak 于} 2000 {\cjkfont 年\allowbreak 代\allowbreak 中\allowbreak 期\allowbreak 。\allowbreak 因\allowbreak 此\allowbreak 由\allowbreak 本\allowbreak 数\allowbreak 据\allowbreak 库\allowbreak 得\allowbreak 出\allowbreak 的\allowbreak 纵\allowbreak 向\allowbreak 结\allowbreak 论\allowbreak 是\allowbreak 两\allowbreak 个\allowbreak 分\allowbreak 离\allowbreak 时\allowbreak 间\allowbreak 窗\allowbreak 之\allowbreak 间\allowbreak 的\allowbreak 比\allowbreak 较\allowbreak ，\allowbreak 而\allowbreak 非\allowbreak 连\allowbreak 续\allowbreak 趋\allowbreak 势\allowbreak ；\allowbreak 我\allowbreak 们\allowbreak 也\allowbreak 不\allowbreak 将\allowbreak 本\allowbreak 语\allowbreak 料\allowbreak 视\allowbreak 为\allowbreak 任\allowbreak 何\allowbreak 总\allowbreak 体\allowbreak 的\allowbreak 代\allowbreak 表\allowbreak 性\allowbreak 样\allowbreak 本\allowbreak 。}

\subsubsection{{\cjkfont 采\allowbreak 集\allowbreak 方\allowbreak 案}}
\textbf{{\cjkfont 网\allowbreak 络\allowbreak 采\allowbreak 集\allowbreak 轨\allowbreak 。}} {\cjkfont 二\allowbreak 十\allowbreak 五\allowbreak 个\allowbreak 采\allowbreak 集\allowbreak 器\allowbreak 获\allowbreak 取\allowbreak 网\allowbreak 络\allowbreak 部\allowbreak 分\allowbreak 。\allowbreak 每\allowbreak 个\allowbreak 采\allowbreak 集\allowbreak 器\allowbreak 实\allowbreak 现\allowbreak 统\allowbreak 一\allowbreak 接\allowbreak 口\allowbreak 并\allowbreak 注\allowbreak 册\allowbreak 于\allowbreak 中\allowbreak 央\allowbreak 注\allowbreak 册\allowbreak 表\allowbreak ；\allowbreak 具\allowbreak 体\allowbreak 实\allowbreak 现\allowbreak 分\allowbreak 为\allowbreak 四\allowbreak 族}------{\cjkfont 一\allowbreak 个\allowbreak 搜\allowbreak 索\allowbreak 引\allowbreak 擎\allowbreak 聚\allowbreak 合\allowbreak 器\allowbreak ，\allowbreak 发\allowbreak 出\allowbreak 预\allowbreak 定\allowbreak 义\allowbreak 查\allowbreak 询\allowbreak 并\allowbreak 仅\allowbreak 接\allowbreak 纳\allowbreak 域\allowbreak 名\allowbreak 白\allowbreak 名\allowbreak 单\allowbreak 内\allowbreak 的\allowbreak 结\allowbreak 果\allowbreak ；\allowbreak 十\allowbreak 九\allowbreak 个\allowbreak 省\allowbreak 级\allowbreak 应\allowbreak 急\allowbreak 管\allowbreak 理\allowbreak 部\allowbreak 门\allowbreak 采\allowbreak 集\allowbreak 器\allowbreak ，\allowbreak 解\allowbreak 析} HTML {\cjkfont 列\allowbreak 表\allowbreak 页\allowbreak 与\allowbreak 详\allowbreak 情\allowbreak 页\allowbreak ，\allowbreak 并\allowbreak 在\allowbreak 多\allowbreak 数\allowbreak 情\allowbreak 况\allowbreak 下\allowbreak 下\allowbreak 载\allowbreak 、\allowbreak 抽\allowbreak 取\allowbreak 随\allowbreak 附\allowbreak 的} PDF {\cjkfont 调\allowbreak 查\allowbreak 报\allowbreak 告\allowbreak ；\allowbreak 两\allowbreak 个\allowbreak 国\allowbreak 家\allowbreak 级\allowbreak 采\allowbreak 集\allowbreak 器\allowbreak （\allowbreak 应\allowbreak 急\allowbreak 管\allowbreak 理\allowbreak 部\allowbreak 与\allowbreak 住\allowbreak 建\allowbreak 部\allowbreak ）\allowbreak ；\allowbreak 以\allowbreak 及\allowbreak 三\allowbreak 个\allowbreak 辅\allowbreak 助\allowbreak 导\allowbreak 入\allowbreak 器\allowbreak ，\allowbreak 分\allowbreak 别\allowbreak 覆\allowbreak 盖\allowbreak 开\allowbreak 放\allowbreak 政\allowbreak 府\allowbreak 数\allowbreak 据\allowbreak 集\allowbreak 、\allowbreak 安\allowbreak 全\allowbreak 新\allowbreak 闻\allowbreak 专\allowbreak 题\allowbreak 归\allowbreak 档\allowbreak 与\allowbreak 两\allowbreak 个\allowbreak 已\allowbreak 公\allowbreak 开\allowbreak 的\allowbreak 学\allowbreak 术\allowbreak 事\allowbreak 故\allowbreak 数\allowbreak 据\allowbreak 集\allowbreak 。}

HTTP {\cjkfont 检\allowbreak 索\allowbreak 在\allowbreak 所\allowbreak 有\allowbreak 采\allowbreak 集\allowbreak 器\allowbreak 中\allowbreak 统\allowbreak 一\allowbreak 使\allowbreak 用\allowbreak 同\allowbreak 一\allowbreak 客\allowbreak 户\allowbreak 端\allowbreak 库\allowbreak （}httpx{\cjkfont ，\allowbreak 超\allowbreak 时} 15--60 {\cjkfont 秒\allowbreak ）\allowbreak 。}HTML {\cjkfont 以} BeautifulSoup {\cjkfont 解\allowbreak 析\allowbreak ，\allowbreak 依\allowbreak 据\allowbreak 来\allowbreak 源\allowbreak 标\allowbreak 记\allowbreak 选\allowbreak 用} lxml {\cjkfont 或\allowbreak 标\allowbreak 准\allowbreak 库\allowbreak 解\allowbreak 析\allowbreak 器\allowbreak 。\allowbreak 正\allowbreak 文\allowbreak 抽\allowbreak 取\allowbreak 优\allowbreak 先\allowbreak 使\allowbreak 用} trafilatura{\cjkfont （\allowbreak 在\allowbreak 页\allowbreak 面\allowbreak 结\allowbreak 构\allowbreak 允\allowbreak 许\allowbreak 时\allowbreak ）\allowbreak ，\allowbreak 回\allowbreak 退\allowbreak 至\allowbreak 移\allowbreak 除} script{\cjkfont 、}style{\cjkfont 、\allowbreak 导\allowbreak 航\allowbreak 、\allowbreak 页\allowbreak 脚\allowbreak 与\allowbreak 表\allowbreak 单\allowbreak 元\allowbreak 素\allowbreak 的} BeautifulSoup {\cjkfont 流\allowbreak 程\allowbreak ，\allowbreak 再\allowbreak 回\allowbreak 退\allowbreak 至\allowbreak 去\allowbreak 标\allowbreak 签\allowbreak 并\allowbreak 反\allowbreak 转\allowbreak 义\allowbreak 实\allowbreak 体\allowbreak 。\allowbreak 随\allowbreak 附} PDF {\cjkfont 在\allowbreak 存\allowbreak 在\allowbreak 文\allowbreak 本\allowbreak 层\allowbreak 时\allowbreak 以} PyMuPDF {\cjkfont 抽\allowbreak 取\allowbreak ，\allowbreak 不\allowbreak 存\allowbreak 在\allowbreak 时\allowbreak 转\allowbreak 交\allowbreak 具\allowbreak 备} OCR {\cjkfont 能\allowbreak 力\allowbreak 的\allowbreak 文\allowbreak 档\allowbreak 解\allowbreak 析\allowbreak 服\allowbreak 务\allowbreak 。}

{\cjkfont 查\allowbreak 询\allowbreak 构\allowbreak 造\allowbreak 结\allowbreak 合\allowbreak 两\allowbreak 种\allowbreak 机\allowbreak 制\allowbreak 。\allowbreak 聚\allowbreak 合\allowbreak 器\allowbreak 发\allowbreak 出} 17 {\cjkfont 条\allowbreak 预\allowbreak 定\allowbreak 义\allowbreak 查\allowbreak 询\allowbreak 模\allowbreak 板\allowbreak ，\allowbreak 覆\allowbreak 盖\allowbreak 坍\allowbreak 塌\allowbreak 、\allowbreak 高\allowbreak 处\allowbreak 坠\allowbreak 落\allowbreak 、\allowbreak 起\allowbreak 重\allowbreak 、\allowbreak 基\allowbreak 坑\allowbreak 、\allowbreak 脚\allowbreak 手\allowbreak 架\allowbreak 、\allowbreak 触\allowbreak 电\allowbreak 、\allowbreak 火\allowbreak 灾\allowbreak 、\allowbreak 中\allowbreak 毒\allowbreak 、\allowbreak 物\allowbreak 体\allowbreak 打\allowbreak 击\allowbreak 与\allowbreak 机\allowbreak 械\allowbreak 等\allowbreak 情\allowbreak 形\allowbreak ，\allowbreak 并\allowbreak 附\allowbreak 部\allowbreak 委\allowbreak 级\allowbreak 报\allowbreak 告\allowbreak 查\allowbreak 询\allowbreak ；\allowbreak 逐\allowbreak 条\allowbreak 记\allowbreak 录\allowbreak 的\allowbreak 搜\allowbreak 索\allowbreak 增\allowbreak 强\allowbreak （}3.4 {\cjkfont 节\allowbreak ）\allowbreak 随\allowbreak 后\allowbreak 根\allowbreak 据\allowbreak 事\allowbreak 故\allowbreak 自\allowbreak 身\allowbreak 的\allowbreak 地\allowbreak 名\allowbreak 与\allowbreak 关\allowbreak 键\allowbreak 词\allowbreak 再\allowbreak 生\allowbreak 成\allowbreak 至\allowbreak 多\allowbreak 三\allowbreak 条\allowbreak 查\allowbreak 询\allowbreak 。\allowbreak 检\allowbreak 索\allowbreak 结\allowbreak 果\allowbreak 的\allowbreak 接\allowbreak 纳\allowbreak 限\allowbreak 定\allowbreak 于} 15 {\cjkfont 个\allowbreak 新\allowbreak 闻\allowbreak 域\allowbreak 名\allowbreak 白\allowbreak 名\allowbreak 单\allowbreak ，\allowbreak 外\allowbreak 加\allowbreak 任\allowbreak 意} `gov.\allowbreak{}cn` {\cjkfont 主\allowbreak 机\allowbreak ，\allowbreak 从\allowbreak 而\allowbreak 使\allowbreak 命\allowbreak 中\allowbreak 检\allowbreak 索\allowbreak 的\allowbreak 资\allowbreak 料\allowbreak 与\allowbreak 定\allowbreak 向\allowbreak 采\allowbreak 集\allowbreak 器\allowbreak 处\allowbreak 于\allowbreak 同\allowbreak 一\allowbreak 来\allowbreak 源\allowbreak 总\allowbreak 体\allowbreak 之\allowbreak 内\allowbreak 。}

{\cjkfont 礼\allowbreak 貌\allowbreak 性\allowbreak 与\allowbreak 稳\allowbreak 健\allowbreak 性\allowbreak 措\allowbreak 施\allowbreak 刻\allowbreak 意\allowbreak 保\allowbreak 持\allowbreak 常\allowbreak 规\allowbreak ：\allowbreak 固\allowbreak 定\allowbreak 的\allowbreak 桌\allowbreak 面\allowbreak 浏\allowbreak 览\allowbreak 器} User-Agent{\cjkfont 、\allowbreak 完\allowbreak 整\allowbreak 的\allowbreak 请\allowbreak 求\allowbreak 头\allowbreak （\allowbreak 含\allowbreak 语\allowbreak 言\allowbreak 偏\allowbreak 好\allowbreak ，\allowbreak 检\allowbreak 索\allowbreak 类\allowbreak 请\allowbreak 求\allowbreak 另\allowbreak 附} referer{\cjkfont ）\allowbreak 、\allowbreak 每\allowbreak 个\allowbreak 采\allowbreak 集\allowbreak 器} 1--5 {\cjkfont 秒\allowbreak 的\allowbreak 串\allowbreak 行\allowbreak 请\allowbreak 求\allowbreak 间\allowbreak 隔\allowbreak 、\allowbreak 每\allowbreak 个} URL {\cjkfont 三\allowbreak 次\allowbreak 尝\allowbreak 试\allowbreak 并\allowbreak 配\allowbreak 两\allowbreak 秒\allowbreak 退\allowbreak 避\allowbreak 。}\textbf{{\cjkfont 我\allowbreak 们\allowbreak 未\allowbreak 实\allowbreak 现\allowbreak 代\allowbreak 理\allowbreak 轮\allowbreak 换\allowbreak 、}User-Agent {\cjkfont 轮\allowbreak 换\allowbreak 、}JavaScript {\cjkfont 渲\allowbreak 染\allowbreak 或\allowbreak 验\allowbreak 证\allowbreak 码\allowbreak 处\allowbreak 理\allowbreak ，\allowbreak 采\allowbreak 集\allowbreak 完\allowbreak 全\allowbreak 串\allowbreak 行\allowbreak 而\allowbreak 非\allowbreak 分\allowbreak 布\allowbreak 式\allowbreak 。}} {\cjkfont 我\allowbreak 们\allowbreak 明\allowbreak 确\allowbreak 说\allowbreak 明\allowbreak 这\allowbreak 一\allowbreak 点\allowbreak ，\allowbreak 因\allowbreak 为\allowbreak 它\allowbreak 同\allowbreak 时\allowbreak 界\allowbreak 定\allowbreak 了\allowbreak 采\allowbreak 集\allowbreak 速\allowbreak 率\allowbreak 与\allowbreak 当\allowbreak 前\allowbreak 语\allowbreak 料\allowbreak 的\allowbreak 可\allowbreak 扩\allowbreak 展\allowbreak 边\allowbreak 界\allowbreak ：\allowbreak 需\allowbreak 要\allowbreak 会\allowbreak 话\allowbreak 保\allowbreak 持\allowbreak 或\allowbreak 动\allowbreak 态\allowbreak 渲\allowbreak 染\allowbreak 的\allowbreak 来\allowbreak 源\allowbreak 落\allowbreak 在\allowbreak 本\allowbreak 方\allowbreak 案\allowbreak 之\allowbreak 外\allowbreak ，\allowbreak 仅\allowbreak 在\allowbreak 其\allowbreak 报\allowbreak 告\allowbreak 被\allowbreak 接\allowbreak 纳\allowbreak 站\allowbreak 点\allowbreak 转\allowbreak 载\allowbreak 的\allowbreak 范\allowbreak 围\allowbreak 内\allowbreak 有\allowbreak 所\allowbreak 体\allowbreak 现\allowbreak 。}

\textbf{{\cjkfont 文\allowbreak 献\allowbreak 数\allowbreak 字\allowbreak 化\allowbreak 轨\allowbreak 。}} {\cjkfont 每\allowbreak 卷\allowbreak 先\allowbreak 转\allowbreak 换\allowbreak 为\allowbreak 文\allowbreak 本\allowbreak （\allowbreak 抽\allowbreak 取\allowbreak 既\allowbreak 有\allowbreak 文\allowbreak 本\allowbreak 层\allowbreak ，\allowbreak 扫\allowbreak 描\allowbreak 卷\allowbreak 则\allowbreak 经} OCR {\cjkfont 文\allowbreak 档\allowbreak 解\allowbreak 析\allowbreak ）\allowbreak ，\allowbreak 随\allowbreak 后\allowbreak 切\allowbreak 分\allowbreak 为\allowbreak 独\allowbreak 立\allowbreak 案\allowbreak 例\allowbreak 。\allowbreak 切\allowbreak 分\allowbreak 提\allowbreak 供\allowbreak 两\allowbreak 种\allowbreak 程\allowbreak 序\allowbreak ：\allowbreak 基\allowbreak 于\allowbreak 规\allowbreak 则\allowbreak 的\allowbreak 切\allowbreak 分\allowbreak 器\allowbreak 锚\allowbreak 定\allowbreak 此\allowbreak 类\allowbreak 汇\allowbreak 编\allowbreak 使\allowbreak 用\allowbreak 的\allowbreak 结\allowbreak 构\allowbreak 性\allowbreak 标\allowbreak 题\allowbreak 惯\allowbreak 例\allowbreak ；\allowbreak 基\allowbreak 于} LLM {\cjkfont 的\allowbreak 切\allowbreak 分\allowbreak 器\allowbreak 返\allowbreak 回\allowbreak 每\allowbreak 个\allowbreak 案\allowbreak 例\allowbreak 的\allowbreak 起\allowbreak 止\allowbreak 标\allowbreak 记\allowbreak ，\allowbreak 随\allowbreak 后\allowbreak 以\allowbreak 字\allowbreak 符\allowbreak 串\allowbreak 检\allowbreak 索}\textbf{{\cjkfont 程\allowbreak 序\allowbreak 化}}{\cjkfont 切\allowbreak 分\allowbreak ，\allowbreak 而\allowbreak 非\allowbreak 接\allowbreak 受\allowbreak 模\allowbreak 型\allowbreak 生\allowbreak 成\allowbreak 的\allowbreak 字\allowbreak 符\allowbreak 偏\allowbreak 移\allowbreak 量}------{\cjkfont 因\allowbreak 此\allowbreak 所\allowbreak 得\allowbreak 跨\allowbreak 度\allowbreak 可\allowbreak 精\allowbreak 确\allowbreak 回\allowbreak 溯\allowbreak 至\allowbreak 源\allowbreak 文\allowbreak 档\allowbreak 并\allowbreak 可\allowbreak 供\allowbreak 审\allowbreak 计\allowbreak 。\allowbreak 低\allowbreak 于\allowbreak 最\allowbreak 小\allowbreak 长\allowbreak 度\allowbreak （}50--100 {\cjkfont 字\allowbreak 符\allowbreak ，\allowbreak 视\allowbreak 切\allowbreak 分\allowbreak 器\allowbreak 而\allowbreak 定\allowbreak ）\allowbreak 的\allowbreak 片\allowbreak 段\allowbreak 被\allowbreak 丢\allowbreak 弃\allowbreak 。\allowbreak 每\allowbreak 条\allowbreak 保\allowbreak 留\allowbreak 案\allowbreak 例\allowbreak 随\allowbreak 后\allowbreak 进\allowbreak 入\allowbreak 与\allowbreak 网\allowbreak 络\allowbreak 记\allowbreak 录\allowbreak 相\allowbreak 同\allowbreak 的\allowbreak 清\allowbreak 洗\allowbreak 、\allowbreak 筛\allowbreak 选\allowbreak 与\allowbreak 编\allowbreak 码\allowbreak 通\allowbreak 路\allowbreak ，\allowbreak 两\allowbreak 条\allowbreak 轨\allowbreak 道\allowbreak 因\allowbreak 而\allowbreak 在\allowbreak 任\allowbreak 何\allowbreak 信\allowbreak 息\allowbreak 抽\allowbreak 取\allowbreak 发\allowbreak 生\allowbreak 之\allowbreak 前\allowbreak 即\allowbreak 已\allowbreak 汇\allowbreak 合\allowbreak 。}

\subsubsection{{\cjkfont 清\allowbreak 洗\allowbreak 、\allowbreak 规\allowbreak 范\allowbreak 化\allowbreak 与\allowbreak 规\allowbreak 则\allowbreak 预\allowbreak 抽\allowbreak 取}}
{\cjkfont 清\allowbreak 洗\allowbreak 与\allowbreak 大\allowbreak 部\allowbreak 分\allowbreak 字\allowbreak 段\allowbreak 抽\allowbreak 取\allowbreak 由\allowbreak 确\allowbreak 定\allowbreak 性\allowbreak 规\allowbreak 则\allowbreak 完\allowbreak 成\allowbreak ，\allowbreak 而\allowbreak 非\allowbreak 模\allowbreak 型\allowbreak 调\allowbreak 用\allowbreak 。\allowbreak 三\allowbreak 个\allowbreak 组\allowbreak 件\allowbreak 依\allowbreak 次\allowbreak 作\allowbreak 用\allowbreak 于\allowbreak 每\allowbreak 条\allowbreak 候\allowbreak 选\allowbreak 记\allowbreak 录\allowbreak 。}

{\cjkfont 文\allowbreak 本\allowbreak 清\allowbreak 洗\allowbreak 器\allowbreak 去\allowbreak 除\allowbreak 模\allowbreak 板\allowbreak 化\allowbreak 内\allowbreak 容\allowbreak 与\allowbreak 渲\allowbreak 染\allowbreak 残\allowbreak 留\allowbreak 。\allowbreak 十\allowbreak 条\allowbreak 正\allowbreak 则\allowbreak 模\allowbreak 式\allowbreak 剥\allowbreak 离\allowbreak 反\allowbreak 复\allowbreak 出\allowbreak 现\allowbreak 的\allowbreak 站\allowbreak 点\allowbreak 装\allowbreak 饰}------{\cjkfont 版\allowbreak 权\allowbreak 声\allowbreak 明\allowbreak 、}ICP {\cjkfont 备\allowbreak 案\allowbreak 号\allowbreak 、\allowbreak 技\allowbreak 术\allowbreak 支\allowbreak 持\allowbreak 行\allowbreak 、\allowbreak 主\allowbreak 办\allowbreak 单\allowbreak 位\allowbreak 行\allowbreak 、\allowbreak 浏\allowbreak 览\allowbreak 器\allowbreak 兼\allowbreak 容\allowbreak 提\allowbreak 示\allowbreak 、\allowbreak 免\allowbreak 责\allowbreak 声\allowbreak 明\allowbreak 与\allowbreak 隐\allowbreak 私\allowbreak 政\allowbreak 策}------{\cjkfont 并\allowbreak 删\allowbreak 除\allowbreak 八\allowbreak 个\allowbreak 零\allowbreak 宽\allowbreak 或\allowbreak 不\allowbreak 可\allowbreak 见\allowbreak 码\allowbreak 点\allowbreak 。}HTML {\cjkfont 实\allowbreak 体\allowbreak 被\allowbreak 反\allowbreak 转\allowbreak 义\allowbreak ，\allowbreak 文\allowbreak 本\allowbreak 归\allowbreak 一\allowbreak 化\allowbreak 为} NFKC{\cjkfont （\allowbreak 将\allowbreak 全\allowbreak 角\allowbreak 形\allowbreak 式\allowbreak 折\allowbreak 叠\allowbreak 为\allowbreak 半\allowbreak 角\allowbreak ）\allowbreak ；\allowbreak 连\allowbreak 续\allowbreak 空\allowbreak 白\allowbreak 折\allowbreak 叠\allowbreak ，\allowbreak 两\allowbreak 个\allowbreak 以\allowbreak 上\allowbreak 连\allowbreak 续\allowbreak 换\allowbreak 行\allowbreak 被\allowbreak 压\allowbreak 缩\allowbreak 。\allowbreak 长\allowbreak 度\allowbreak 不\allowbreak 足\allowbreak 四\allowbreak 字\allowbreak 符\allowbreak 且\allowbreak 不\allowbreak 含\allowbreak 数\allowbreak 字\allowbreak 的\allowbreak 行\allowbreak 被\allowbreak 丢\allowbreak 弃}------{\cjkfont 该\allowbreak 过\allowbreak 滤\allowbreak 剔\allowbreak 除\allowbreak 零\allowbreak 散\allowbreak 标\allowbreak 题\allowbreak ，\allowbreak 同\allowbreak 时\allowbreak 保\allowbreak 留\allowbreak 承\allowbreak 载\allowbreak 伤\allowbreak 亡\allowbreak 数\allowbreak 字\allowbreak 的\allowbreak 短\allowbreak 行\allowbreak 。\allowbreak 超\allowbreak 过} 50,000 {\cjkfont 字\allowbreak 符\allowbreak 的\allowbreak 文\allowbreak 档\allowbreak 被\allowbreak 截\allowbreak 断\allowbreak 并\allowbreak 保\allowbreak 留\allowbreak 前\allowbreak 段\allowbreak ，\allowbreak 因\allowbreak 为\allowbreak 事\allowbreak 故\allowbreak 叙\allowbreak 述\allowbreak 通\allowbreak 常\allowbreak 在\allowbreak 前\allowbreak 部\allowbreak 即\allowbreak 陈\allowbreak 述\allowbreak 要\allowbreak 点\allowbreak 。\allowbreak 清\allowbreak 洗\allowbreak 器\allowbreak 对\allowbreak 网\allowbreak 络\allowbreak 文\allowbreak 本\allowbreak 与\allowbreak 数\allowbreak 字\allowbreak 化\allowbreak 案\allowbreak 例\allowbreak 文\allowbreak 本\allowbreak 施\allowbreak 加\allowbreak 完\allowbreak 全\allowbreak 相\allowbreak 同\allowbreak 的\allowbreak 处\allowbreak 理\allowbreak 。}

{\cjkfont 规\allowbreak 则\allowbreak 字\allowbreak 段\allowbreak 抽\allowbreak 取\allowbreak 器\allowbreak 随\allowbreak 后\allowbreak 对\allowbreak 每\allowbreak 条\allowbreak 清\allowbreak 洗\allowbreak 后\allowbreak 文\allowbreak 本\allowbreak 施\allowbreak 加\allowbreak 有\allowbreak 序\allowbreak 正\allowbreak 则\allowbreak 模\allowbreak 式\allowbreak ，\allowbreak 覆\allowbreak 盖\allowbreak 八\allowbreak 个\allowbreak 字\allowbreak 段\allowbreak 组}------{\cjkfont 日\allowbreak 期\allowbreak 、\allowbreak 地\allowbreak 点\allowbreak 、\allowbreak 伤\allowbreak 亡\allowbreak 、\allowbreak 事\allowbreak 故\allowbreak 类\allowbreak 型\allowbreak 、\allowbreak 等\allowbreak 级\allowbreak 、\allowbreak 直\allowbreak 接\allowbreak 经\allowbreak 济\allowbreak 损\allowbreak 失\allowbreak 、\allowbreak 施\allowbreak 工\allowbreak 单\allowbreak 位\allowbreak 与\allowbreak 项\allowbreak 目\allowbreak 名\allowbreak 称\allowbreak 。\allowbreak 模\allowbreak 式\allowbreak 按\allowbreak 特\allowbreak 异\allowbreak 性\allowbreak 排\allowbreak 序\allowbreak 并\allowbreak 携\allowbreak 带\allowbreak 优\allowbreak 先\allowbreak 级\allowbreak 权\allowbreak 重\allowbreak ；\allowbreak 例\allowbreak 如\allowbreak 完\allowbreak 整\allowbreak 日\allowbreak 期\allowbreak （}`YYYY{\cjkfont 年}M{\cjkfont 月}D{\cjkfont 日}`{\cjkfont ）\allowbreak 优\allowbreak 先\allowbreak 于\allowbreak 年\allowbreak 月\allowbreak 形\allowbreak 式\allowbreak ，\allowbreak 年\allowbreak 月\allowbreak 形\allowbreak 式\allowbreak 又\allowbreak 优\allowbreak 先\allowbreak 于\allowbreak 依\allowbreak 上\allowbreak 下\allowbreak 文\allowbreak 回\allowbreak 溯\allowbreak 的\allowbreak 月\allowbreak 日\allowbreak 形\allowbreak 式\allowbreak 。\allowbreak 每\allowbreak 个\allowbreak 抽\allowbreak 取\allowbreak 值\allowbreak 都\allowbreak 通\allowbreak 过\allowbreak 合\allowbreak 理\allowbreak 性\allowbreak 检\allowbreak 查\allowbreak （\allowbreak 年\allowbreak 份\allowbreak 在} 2000--2030 {\cjkfont 之\allowbreak 间\allowbreak 、\allowbreak 死\allowbreak 亡\allowbreak 人\allowbreak 数\allowbreak 低\allowbreak 于} 200{\cjkfont 、\allowbreak 受\allowbreak 伤\allowbreak 人\allowbreak 数\allowbreak 低\allowbreak 于} 500{\cjkfont ）\allowbreak 。\allowbreak 加\allowbreak 权\allowbreak 置\allowbreak 信\allowbreak 度\allowbreak 综\allowbreak 合\allowbreak 日\allowbreak 期\allowbreak 、\allowbreak 城\allowbreak 市\allowbreak 、\allowbreak 事\allowbreak 故\allowbreak 类\allowbreak 型\allowbreak 与\allowbreak 死\allowbreak 亡\allowbreak 人\allowbreak 数\allowbreak 的\allowbreak 命\allowbreak 中\allowbreak 情\allowbreak 况\allowbreak ，\allowbreak 并\allowbreak 对\allowbreak 每\allowbreak 多\allowbreak 回\allowbreak 收\allowbreak 一\allowbreak 个\allowbreak 字\allowbreak 段\allowbreak 给\allowbreak 予\allowbreak 小\allowbreak 幅\allowbreak 加\allowbreak 成\allowbreak 。}

{\cjkfont 规\allowbreak 范\allowbreak 化\allowbreak 器\allowbreak 将\allowbreak 抽\allowbreak 取\allowbreak 值\allowbreak 与\allowbreak 自\allowbreak 由\allowbreak 文\allowbreak 本\allowbreak 值\allowbreak 映\allowbreak 射\allowbreak 至\allowbreak 受\allowbreak 控\allowbreak 词\allowbreak 表\allowbreak ：\allowbreak 事\allowbreak 故\allowbreak 类\allowbreak 型\allowbreak 映\allowbreak 射\allowbreak 至\allowbreak 十\allowbreak 二\allowbreak 类\allowbreak 之\allowbreak 一\allowbreak ；\allowbreak 等\allowbreak 级\allowbreak 映\allowbreak 射\allowbreak 至\allowbreak 中\allowbreak 国\allowbreak 事\allowbreak 故\allowbreak 法\allowbreak 规\allowbreak 规\allowbreak 定\allowbreak 的\allowbreak 四\allowbreak 档}------{\cjkfont 一\allowbreak 般\allowbreak 、\allowbreak 较\allowbreak 大\allowbreak 、\allowbreak 重\allowbreak 大\allowbreak 、\allowbreak 特\allowbreak 别\allowbreak 重\allowbreak 大\allowbreak ，\allowbreak 以} 3{\cjkfont 、}10{\cjkfont 、}30 {\cjkfont 人\allowbreak 死\allowbreak 亡\allowbreak 为\allowbreak 阈\allowbreak 值}------{\cjkfont 且\allowbreak 文\allowbreak 本\allowbreak 中\allowbreak 明\allowbreak 示\allowbreak 的\allowbreak 等\allowbreak 级\allowbreak 词\allowbreak 优\allowbreak 先\allowbreak 于\allowbreak 推\allowbreak 断\allowbreak 档\allowbreak 位\allowbreak ；\allowbreak 省\allowbreak 市\allowbreak 映\allowbreak 射\allowbreak 至\allowbreak 规\allowbreak 范\allowbreak 行\allowbreak 政\allowbreak 区\allowbreak 名\allowbreak ，\allowbreak 依\allowbreak 据\allowbreak 一\allowbreak 份\allowbreak 逾\allowbreak 九\allowbreak 十\allowbreak 条\allowbreak 的\allowbreak 城\allowbreak 市}-{\cjkfont 省\allowbreak 份\allowbreak 对\allowbreak 照\allowbreak 表\allowbreak ；\allowbreak 日\allowbreak 期\allowbreak 映\allowbreak 射\allowbreak 至} ISO {\cjkfont 格\allowbreak 式\allowbreak ，\allowbreak 中\allowbreak 文\allowbreak 数\allowbreak 字\allowbreak 表\allowbreak 述\allowbreak 转\allowbreak 换\allowbreak 为\allowbreak 整\allowbreak 数\allowbreak 。}

{\cjkfont 该\allowbreak 层\allowbreak 的\allowbreak 经\allowbreak 济\allowbreak 考\allowbreak 量\allowbreak 是\allowbreak 显\allowbreak 式\allowbreak 的\allowbreak 。\allowbreak 由\allowbreak 于\allowbreak 它\allowbreak 确\allowbreak 定\allowbreak 性\allowbreak 执\allowbreak 行\allowbreak 且\allowbreak 不\allowbreak 产\allowbreak 生} API {\cjkfont 成\allowbreak 本\allowbreak ，\allowbreak 它\allowbreak 对\allowbreak 每\allowbreak 条\allowbreak 记\allowbreak 录\allowbreak 无\allowbreak 条\allowbreak 件\allowbreak 运\allowbreak 行\allowbreak ，\allowbreak 其\allowbreak 置\allowbreak 信\allowbreak 度\allowbreak 输\allowbreak 出\allowbreak 决\allowbreak 定\allowbreak 后\allowbreak 续\allowbreak 动\allowbreak 作\allowbreak ：\allowbreak 置\allowbreak 信\allowbreak 度\allowbreak 低\allowbreak 于} 0.7{\cjkfont 、\allowbreak 或\allowbreak 缺\allowbreak 失\allowbreak 四\allowbreak 个\allowbreak 关\allowbreak 键\allowbreak 字\allowbreak 段\allowbreak 中\allowbreak 三\allowbreak 个\allowbreak 以\allowbreak 上\allowbreak 的\allowbreak 记\allowbreak 录\allowbreak 被\allowbreak 标\allowbreak 记\allowbreak 为\allowbreak 待\allowbreak 增\allowbreak 强\allowbreak ，\allowbreak 其\allowbreak 余\allowbreak 记\allowbreak 录\allowbreak 则\allowbreak 以\allowbreak 规\allowbreak 则\allowbreak 值\allowbreak 直\allowbreak 接\allowbreak 推\allowbreak 进\allowbreak 。\allowbreak 该\allowbreak 设\allowbreak 计\allowbreak 使\allowbreak 多\allowbreak 数\allowbreak 字\allowbreak 段\allowbreak 级\allowbreak 取\allowbreak 值\allowbreak 独\allowbreak 立\allowbreak 于\allowbreak 任\allowbreak 何\allowbreak 模\allowbreak 型\allowbreak ，\allowbreak 这\allowbreak 也\allowbreak 正\allowbreak 是} 4.3 {\cjkfont 节\allowbreak 得\allowbreak 以\allowbreak 将\allowbreak 规\allowbreak 则\allowbreak 导\allowbreak 出\allowbreak 字\allowbreak 段\allowbreak 与\allowbreak 模\allowbreak 型\allowbreak 导\allowbreak 出\allowbreak 字\allowbreak 段\allowbreak 同\allowbreak 独\allowbreak 立\allowbreak 标\allowbreak 注\allowbreak 者\allowbreak 比\allowbreak 对\allowbreak 的\allowbreak 前\allowbreak 提\allowbreak 。}

\subsubsection{{\cjkfont 相\allowbreak 关\allowbreak 性\allowbreak 筛\allowbreak 选\allowbreak 与\allowbreak 去\allowbreak 重}}
\textbf{{\cjkfont 相\allowbreak 关\allowbreak 性\allowbreak 筛\allowbreak 选\allowbreak 。}} {\cjkfont 并\allowbreak 非\allowbreak 每\allowbreak 个\allowbreak 采\allowbreak 集\allowbreak 到\allowbreak 的\allowbreak 页\allowbreak 面\allowbreak 都\allowbreak 涉\allowbreak 及\allowbreak 施\allowbreak 工\allowbreak 事\allowbreak 故\allowbreak ；\allowbreak 门\allowbreak 户\allowbreak 同\allowbreak 时\allowbreak 刊\allowbreak 载\allowbreak 监\allowbreak 管\allowbreak 通\allowbreak 知\allowbreak 、\allowbreak 培\allowbreak 训\allowbreak 公\allowbreak 告\allowbreak 以\allowbreak 及\allowbreak 采\allowbreak 矿\allowbreak 、\allowbreak 化\allowbreak 工\allowbreak 与\allowbreak 运\allowbreak 输\allowbreak 行\allowbreak 业\allowbreak 的\allowbreak 工\allowbreak 业\allowbreak 事\allowbreak 故\allowbreak 。\allowbreak 筛\allowbreak 选\allowbreak 分\allowbreak 两\allowbreak 道\allowbreak 闸\allowbreak 门\allowbreak ，\allowbreak 后\allowbreak 者\allowbreak 与\allowbreak 管\allowbreak 道\allowbreak 阶\allowbreak 段\allowbreak 一\allowbreak 并\allowbreak 说\allowbreak 明\allowbreak 。}

{\cjkfont 第\allowbreak 一\allowbreak 道\allowbreak 闸\allowbreak 门\allowbreak 作\allowbreak 用\allowbreak 于\allowbreak 标\allowbreak 题\allowbreak 。}26 {\cjkfont 个\allowbreak 施\allowbreak 工\allowbreak 关\allowbreak 键\allowbreak 词}------{\cjkfont 覆\allowbreak 盖\allowbreak 建\allowbreak 筑\allowbreak 、\allowbreak 市\allowbreak 政\allowbreak 、\allowbreak 拆\allowbreak 除\allowbreak 、\allowbreak 脚\allowbreak 手\allowbreak 架\allowbreak 、\allowbreak 塔\allowbreak 吊\allowbreak 、\allowbreak 基\allowbreak 坑\allowbreak 、\allowbreak 模\allowbreak 板\allowbreak 及\allowbreak 标\allowbreak 准\allowbreak 事\allowbreak 故\allowbreak 类\allowbreak 型\allowbreak 术\allowbreak 语}------{\cjkfont 命\allowbreak 中\allowbreak 任\allowbreak 一\allowbreak 即\allowbreak 放\allowbreak 行\allowbreak ；\allowbreak 八\allowbreak 条\allowbreak 否\allowbreak 决\allowbreak 模\allowbreak 式\allowbreak 直\allowbreak 接\allowbreak 剔\allowbreak 除\allowbreak 行\allowbreak 政\allowbreak 文\allowbreak 件\allowbreak （\allowbreak 挂\allowbreak 牌\allowbreak 督\allowbreak 办\allowbreak 、\allowbreak 暂\allowbreak 行\allowbreak 办\allowbreak 法\allowbreak 、\allowbreak 管\allowbreak 理\allowbreak 办\allowbreak 法\allowbreak 、\allowbreak 实\allowbreak 施\allowbreak 细\allowbreak 则\allowbreak 、\allowbreak 培\allowbreak 训\allowbreak 班\allowbreak ，\allowbreak 以\allowbreak 及\allowbreak 事\allowbreak 故\allowbreak 报\allowbreak 告\allowbreak 条\allowbreak 例\allowbreak 本\allowbreak 身\allowbreak ）\allowbreak 。\allowbreak 该\allowbreak 闸\allowbreak 门\allowbreak 刻\allowbreak 意\allowbreak 宽\allowbreak 松\allowbreak ：\allowbreak 其\allowbreak 目\allowbreak 的\allowbreak 是\allowbreak 廉\allowbreak 价\allowbreak 地\allowbreak 丢\allowbreak 弃\allowbreak 明\allowbreak 显\allowbreak 的\allowbreak 非\allowbreak 事\allowbreak 故\allowbreak 内\allowbreak 容\allowbreak ，\allowbreak 而\allowbreak 非\allowbreak 裁\allowbreak 决\allowbreak 边\allowbreak 界\allowbreak 案\allowbreak 例\allowbreak 。}

{\cjkfont 通\allowbreak 过\allowbreak 标\allowbreak 题\allowbreak 闸\allowbreak 门\allowbreak 的\allowbreak 记\allowbreak 录\allowbreak 进\allowbreak 入\allowbreak 规\allowbreak 则\allowbreak 分\allowbreak 类\allowbreak 器\allowbreak ，\allowbreak 该\allowbreak 分\allowbreak 类\allowbreak 器\allowbreak 对\allowbreak 标\allowbreak 题\allowbreak 与\allowbreak 正\allowbreak 文\allowbreak 开\allowbreak 头\allowbreak 评\allowbreak 分\allowbreak 四\allowbreak 张\allowbreak 信\allowbreak 号\allowbreak 词\allowbreak 表\allowbreak 。\allowbreak 第\allowbreak 一\allowbreak 张\allowbreak 列\allowbreak 出\allowbreak 约\allowbreak 三\allowbreak 十\allowbreak 个\allowbreak 几\allowbreak 乎\allowbreak 只\allowbreak 出\allowbreak 现\allowbreak 在\allowbreak 施\allowbreak 工\allowbreak 场\allowbreak 景\allowbreak 的\allowbreak 术\allowbreak 语\allowbreak （\allowbreak 塔\allowbreak 式\allowbreak 起\allowbreak 重\allowbreak 机\allowbreak 、\allowbreak 施\allowbreak 工\allowbreak 升\allowbreak 降\allowbreak 机\allowbreak 、\allowbreak 基\allowbreak 坑\allowbreak 工\allowbreak 程\allowbreak 、\allowbreak 模\allowbreak 板\allowbreak 支\allowbreak 撑\allowbreak 、\allowbreak 混\allowbreak 凝\allowbreak 土\allowbreak 浇\allowbreak 筑\allowbreak 、\allowbreak 钢\allowbreak 筋\allowbreak 绑\allowbreak 扎\allowbreak 、\allowbreak 土\allowbreak 方\allowbreak 开\allowbreak 挖\allowbreak 、\allowbreak 桩\allowbreak 基\allowbreak ）\allowbreak 。\allowbreak 第\allowbreak 二\allowbreak 张\allowbreak 列\allowbreak 出\allowbreak 约\allowbreak 二\allowbreak 十\allowbreak 八\allowbreak 个\allowbreak 同\allowbreak 时\allowbreak 也\allowbreak 出\allowbreak 现\allowbreak 在\allowbreak 工\allowbreak 业\allowbreak 场\allowbreak 景\allowbreak 的\allowbreak 施\allowbreak 工\allowbreak 术\allowbreak 语\allowbreak （\allowbreak 施\allowbreak 工\allowbreak 现\allowbreak 场\allowbreak 、\allowbreak 建\allowbreak 筑\allowbreak 工\allowbreak 程\allowbreak 、\allowbreak 脚\allowbreak 手\allowbreak 架\allowbreak 、\allowbreak 高\allowbreak 处\allowbreak 作\allowbreak 业\allowbreak 、\allowbreak 拆\allowbreak 除\allowbreak 工\allowbreak 程\allowbreak ）\allowbreak 。\allowbreak 第\allowbreak 三\allowbreak 张\allowbreak 列\allowbreak 出\allowbreak 约\allowbreak 五\allowbreak 十\allowbreak 个\allowbreak 工\allowbreak 业\allowbreak 与\allowbreak 采\allowbreak 掘\allowbreak 术\allowbreak 语\allowbreak （\allowbreak 矿\allowbreak 业\allowbreak 、\allowbreak 钢\allowbreak 铁\allowbreak 、\allowbreak 水\allowbreak 泥\allowbreak 、\allowbreak 化\allowbreak 工\allowbreak 、\allowbreak 纺\allowbreak 织\allowbreak 、\allowbreak 电\allowbreak 力\allowbreak ）\allowbreak 。\allowbreak 第\allowbreak 四\allowbreak 张\allowbreak 列\allowbreak 出\allowbreak 约\allowbreak 二\allowbreak 十\allowbreak 个\allowbreak 否\allowbreak 决\allowbreak 术\allowbreak 语\allowbreak （\allowbreak 道\allowbreak 路\allowbreak 交\allowbreak 通\allowbreak 、\allowbreak 高\allowbreak 速\allowbreak 公\allowbreak 路\allowbreak 、\allowbreak 列\allowbreak 车\allowbreak 脱\allowbreak 轨\allowbreak 、\allowbreak 煤\allowbreak 矿\allowbreak 瓦\allowbreak 斯\allowbreak 、\allowbreak 渔\allowbreak 船\allowbreak 、\allowbreak 防\allowbreak 汛\allowbreak 、\allowbreak 危\allowbreak 险\allowbreak 化\allowbreak 学\allowbreak 品\allowbreak ）\allowbreak 。\allowbreak 七\allowbreak 条\allowbreak 有\allowbreak 序\allowbreak 规\allowbreak 则\allowbreak 随\allowbreak 后\allowbreak 组\allowbreak 合\allowbreak 这\allowbreak 些\allowbreak 信\allowbreak 号\allowbreak ：\allowbreak 命\allowbreak 中\allowbreak 施\allowbreak 工\allowbreak 特\allowbreak 定\allowbreak 词\allowbreak 即\allowbreak 便\allowbreak 出\allowbreak 现\allowbreak 工\allowbreak 业\allowbreak 公\allowbreak 司\allowbreak 也\allowbreak 放\allowbreak 行\allowbreak ；\allowbreak 工\allowbreak 业\allowbreak 公\allowbreak 司\allowbreak 且\allowbreak 施\allowbreak 工\allowbreak 信\allowbreak 号\allowbreak 少\allowbreak 于\allowbreak 三\allowbreak 个\allowbreak 则\allowbreak 否\allowbreak 决\allowbreak ；\allowbreak 非\allowbreak 工\allowbreak 业\allowbreak 公\allowbreak 司\allowbreak 且\allowbreak 施\allowbreak 工\allowbreak 信\allowbreak 号\allowbreak 达\allowbreak 到\allowbreak 两\allowbreak 个\allowbreak 以\allowbreak 上\allowbreak 则\allowbreak 放\allowbreak 行\allowbreak ；\allowbreak 仅\allowbreak 命\allowbreak 中\allowbreak 一\allowbreak 个\allowbreak 施\allowbreak 工\allowbreak 信\allowbreak 号\allowbreak 则\allowbreak 标\allowbreak 记\allowbreak 为\allowbreak 低\allowbreak 置\allowbreak 信\allowbreak 度\allowbreak 并\allowbreak 转\allowbreak 交\allowbreak 后\allowbreak 续\allowbreak 而\allowbreak 非\allowbreak 径\allowbreak 行\allowbreak 裁\allowbreak 决\allowbreak 。\allowbreak 分\allowbreak 类\allowbreak 器\allowbreak 返\allowbreak 回\allowbreak 带} \textit{high} {\cjkfont 或} \textit{low} {\cjkfont 置\allowbreak 信\allowbreak 度\allowbreak 的\allowbreak 二\allowbreak 值\allowbreak 判\allowbreak 断\allowbreak ；\allowbreak 仅\allowbreak 高\allowbreak 置\allowbreak 信\allowbreak 判\allowbreak 断\allowbreak 为\allowbreak 最\allowbreak 终\allowbreak 结\allowbreak 果\allowbreak 。}

\textbf{{\cjkfont 去\allowbreak 重\allowbreak 。}} {\cjkfont 同\allowbreak 一\allowbreak 起\allowbreak 事\allowbreak 故\allowbreak 在\allowbreak 中\allowbreak 国\allowbreak 新\allowbreak 闻\allowbreak 门\allowbreak 户\allowbreak 间\allowbreak 被\allowbreak 转\allowbreak 载\allowbreak 、\allowbreak 镜\allowbreak 像\allowbreak 报\allowbreak 道\allowbreak 十\allowbreak 分\allowbreak 常\allowbreak 见\allowbreak ，\allowbreak 因\allowbreak 此\allowbreak 记\allowbreak 录\allowbreak 的\allowbreak 来\allowbreak 源\allowbreak 并\allowbreak 不\allowbreak 构\allowbreak 成\allowbreak 其\allowbreak 唯\allowbreak 一\allowbreak 性\allowbreak 的\allowbreak 证\allowbreak 据\allowbreak 。\allowbreak 管\allowbreak 道\allowbreak 施\allowbreak 加\allowbreak 四\allowbreak 层\allowbreak 递\allowbreak 进\allowbreak 程\allowbreak 序\allowbreak ，\allowbreak 由\allowbreak 廉\allowbreak 价\allowbreak 且\allowbreak 精\allowbreak 确\allowbreak 者\allowbreak 向\allowbreak 昂\allowbreak 贵\allowbreak 且\allowbreak 近\allowbreak 似\allowbreak 者\allowbreak 排\allowbreak 序\allowbreak 。\allowbreak 第} 1 {\cjkfont 层\allowbreak 比\allowbreak 对} URL {\cjkfont 身\allowbreak 份\allowbreak ，\allowbreak 使\allowbreak 用\allowbreak 去\allowbreak 除\allowbreak 协\allowbreak 议\allowbreak 前\allowbreak 缀\allowbreak 、}`www`{\cjkfont 、\allowbreak 末\allowbreak 尾\allowbreak 斜\allowbreak 杠\allowbreak 与\allowbreak 片\allowbreak 段\allowbreak 后\allowbreak 的} URL {\cjkfont 哈\allowbreak 希\allowbreak 。\allowbreak 第} 2 {\cjkfont 层\allowbreak 比\allowbreak 对\allowbreak 结\allowbreak 构\allowbreak 化\allowbreak 指\allowbreak 纹\allowbreak ，\allowbreak 由\allowbreak 年\allowbreak 月\allowbreak 、\allowbreak 城\allowbreak 市\allowbreak 、\allowbreak 死\allowbreak 亡\allowbreak 人\allowbreak 数\allowbreak 、\allowbreak 受\allowbreak 伤\allowbreak 人\allowbreak 数\allowbreak 与\allowbreak 事\allowbreak 故\allowbreak 类\allowbreak 型\allowbreak 组\allowbreak 装\allowbreak 并\allowbreak 哈\allowbreak 希\allowbreak 为\allowbreak 十\allowbreak 二\allowbreak 位\allowbreak 十\allowbreak 六\allowbreak 进\allowbreak 制\allowbreak 字\allowbreak 符\allowbreak ；\allowbreak 由\allowbreak 于\allowbreak 指\allowbreak 纹\allowbreak 使\allowbreak 用\allowbreak 年\allowbreak 月\allowbreak 而\allowbreak 非\allowbreak 完\allowbreak 整\allowbreak 日\allowbreak 期\allowbreak ，\allowbreak 它\allowbreak 容\allowbreak 忍\allowbreak 门\allowbreak 户\allowbreak 转\allowbreak 载\allowbreak 调\allowbreak 查\allowbreak 报\allowbreak 告\allowbreak 时\allowbreak 出\allowbreak 现\allowbreak 的\allowbreak 数\allowbreak 日\allowbreak 发\allowbreak 布\allowbreak 日\allowbreak 差\allowbreak 异\allowbreak 。\allowbreak 第} 3 {\cjkfont 层\allowbreak 比\allowbreak 对\allowbreak 标\allowbreak 题\allowbreak 语\allowbreak 义\allowbreak ，\allowbreak 以\allowbreak 中\allowbreak 文\allowbreak 句\allowbreak 向\allowbreak 量\allowbreak 模\allowbreak 型\allowbreak 编\allowbreak 码\allowbreak 标\allowbreak 题\allowbreak 并\allowbreak 对\allowbreak 余\allowbreak 弦\allowbreak 相\allowbreak 似\allowbreak 度\allowbreak 达\allowbreak 到} 0.85 {\cjkfont 的\allowbreak 配\allowbreak 对\allowbreak 予\allowbreak 以\allowbreak 标\allowbreak 记\allowbreak ；\allowbreak 候\allowbreak 选\allowbreak 集\allowbreak 限\allowbreak 定\allowbreak 于\allowbreak 以\allowbreak 月\allowbreak 份\allowbreak 与\allowbreak 省\allowbreak 份\allowbreak 为\allowbreak 键\allowbreak 的\allowbreak 时\allowbreak 空\allowbreak 桶\allowbreak 内\allowbreak ，\allowbreak 从\allowbreak 而\allowbreak 使\allowbreak 比\allowbreak 对\allowbreak 量\allowbreak 不\allowbreak 随\allowbreak 语\allowbreak 料\allowbreak 规\allowbreak 模\allowbreak 增\allowbreak 长\allowbreak 。\allowbreak 当\allowbreak 句\allowbreak 向\allowbreak 量\allowbreak 模\allowbreak 型\allowbreak 不\allowbreak 可\allowbreak 用\allowbreak 时\allowbreak ，\allowbreak 改\allowbreak 以\allowbreak 阈\allowbreak 值} 0.80 {\cjkfont 的\allowbreak 编\allowbreak 辑\allowbreak 距\allowbreak 离\allowbreak 比\allowbreak 替\allowbreak 代\allowbreak 。\allowbreak 第} 4 {\cjkfont 层\allowbreak 独\allowbreak 立\allowbreak 于\allowbreak 标\allowbreak 题\allowbreak 措\allowbreak 辞\allowbreak 交\allowbreak 叉\allowbreak 校\allowbreak 验\allowbreak 四\allowbreak 项\allowbreak 事\allowbreak 实}------{\cjkfont 日\allowbreak 期\allowbreak 邻\allowbreak 近\allowbreak 性\allowbreak 、\allowbreak 城\allowbreak 市\allowbreak 同\allowbreak 一\allowbreak 性\allowbreak 、\allowbreak 伤\allowbreak 亡\allowbreak 吻\allowbreak 合\allowbreak 度\allowbreak 与\allowbreak 事\allowbreak 故\allowbreak 类\allowbreak 型\allowbreak 吻\allowbreak 合\allowbreak 度}------{\cjkfont 并\allowbreak 要\allowbreak 求\allowbreak 至\allowbreak 少\allowbreak 两\allowbreak 项\allowbreak 检\allowbreak 查\allowbreak 通\allowbreak 过\allowbreak 且\allowbreak 归\allowbreak 一\allowbreak 化\allowbreak 吻\allowbreak 合\allowbreak 度\allowbreak 高\allowbreak 于} 0.75{\cjkfont 。}

{\cjkfont 设\allowbreak 计\allowbreak 意\allowbreak 图\allowbreak 是\allowbreak 由\allowbreak 精\allowbreak 确\allowbreak 层\allowbreak 与\allowbreak 结\allowbreak 构\allowbreak 层\allowbreak 拦\allowbreak 截\allowbreak 绝\allowbreak 大\allowbreak 多\allowbreak 数\allowbreak 重\allowbreak 复\allowbreak ，\allowbreak 仅\allowbreak 真\allowbreak 正\allowbreak 歧\allowbreak 义\allowbreak 的\allowbreak 案\allowbreak 例\allowbreak 进\allowbreak 入\allowbreak 语\allowbreak 义\allowbreak 层\allowbreak 。\allowbreak 检\allowbreak 出\allowbreak 去\allowbreak 重\allowbreak 时\allowbreak ，\allowbreak 管\allowbreak 道\allowbreak 并\allowbreak 不\allowbreak 丢\allowbreak 弃\allowbreak 新\allowbreak 到\allowbreak 的\allowbreak 记\allowbreak 录\allowbreak ，\allowbreak 而\allowbreak 是}\textbf{{\cjkfont 合\allowbreak 并}}{\cjkfont ：\allowbreak 新\allowbreak 来\allowbreak 源} URL {\cjkfont 追\allowbreak 加\allowbreak 至\allowbreak 既\allowbreak 有\allowbreak 记\allowbreak 录\allowbreak ，\allowbreak 来\allowbreak 源\allowbreak 计\allowbreak 数\allowbreak 递\allowbreak 增\allowbreak ，\allowbreak 既\allowbreak 有\allowbreak 记\allowbreak 录\allowbreak 缺\allowbreak 失\allowbreak 的\allowbreak 字\allowbreak 段\allowbreak 由\allowbreak 新\allowbreak 记\allowbreak 录\allowbreak 补\allowbreak 全\allowbreak 。\allowbreak 数\allowbreak 据\allowbreak 库\allowbreak 中\allowbreak 有} 45 {\cjkfont 条\allowbreak 记\allowbreak 录\allowbreak 携\allowbreak 带\allowbreak 一\allowbreak 个\allowbreak 以\allowbreak 上\allowbreak 的\allowbreak 来\allowbreak 源} URL{\cjkfont ，\allowbreak 合\allowbreak 并\allowbreak 后\allowbreak 的\allowbreak 集\allowbreak 合\allowbreak 通\allowbreak 过\allowbreak 界\allowbreak 面\allowbreak 开\allowbreak 放\allowbreak ，\allowbreak 因\allowbreak 此\allowbreak 用\allowbreak 户\allowbreak 可\allowbreak 以\allowbreak 看\allowbreak 到\allowbreak 某\allowbreak 项\allowbreak 发\allowbreak 现\allowbreak 依\allowbreak 托\allowbreak 于\allowbreak 、\allowbreak 或\allowbreak 得\allowbreak 到\allowbreak 不\allowbreak 止\allowbreak 一\allowbreak 份\allowbreak 刊\allowbreak 物\allowbreak 的\allowbreak 印\allowbreak 证\allowbreak 。}

\textbf{{\cjkfont 整\allowbreak 合\allowbreak 。}} {\cjkfont 被\allowbreak 接\allowbreak 纳\allowbreak 的\allowbreak 记\allowbreak 录\allowbreak 依\allowbreak 据\allowbreak 两\allowbreak 条\allowbreak 幂\allowbreak 等\allowbreak 规\allowbreak 则\allowbreak 写\allowbreak 入\allowbreak 存\allowbreak 储\allowbreak 。\allowbreak 原\allowbreak 始\allowbreak 候\allowbreak 选\allowbreak 文\allowbreak 本\allowbreak 以\allowbreak 其\allowbreak 来\allowbreak 源} URL {\cjkfont 的\allowbreak 哈\allowbreak 希\allowbreak 为\allowbreak 键\allowbreak ，\allowbreak 因\allowbreak 此\allowbreak 对\allowbreak 同\allowbreak 一\allowbreak 来\allowbreak 源\allowbreak 重\allowbreak 复\allowbreak 采\allowbreak 集\allowbreak 不\allowbreak 会\allowbreak 新\allowbreak 增\allowbreak 任\allowbreak 何\allowbreak 内\allowbreak 容\allowbreak ，\allowbreak 管\allowbreak 道\allowbreak 可\allowbreak 在\allowbreak 中\allowbreak 断\allowbreak 后\allowbreak 续\allowbreak 跑\allowbreak 。\allowbreak 分\allowbreak 类}------{\cjkfont 单\allowbreak 条\allowbreak 记\allowbreak 录\allowbreak 成\allowbreak 本\allowbreak 最\allowbreak 高\allowbreak 的\allowbreak 一\allowbreak 步}------{\cjkfont 每\allowbreak 十\allowbreak 条\allowbreak 记\allowbreak 录\allowbreak 设\allowbreak 立\allowbreak 检\allowbreak 查\allowbreak 点\allowbreak ，\allowbreak 中\allowbreak 断\allowbreak 的\allowbreak 运\allowbreak 行\allowbreak 从\allowbreak 最\allowbreak 近\allowbreak 检\allowbreak 查\allowbreak 点\allowbreak 恢\allowbreak 复\allowbreak ，\allowbreak 而\allowbreak 非\allowbreak 重\allowbreak 新\allowbreak 发\allowbreak 出\allowbreak 模\allowbreak 型\allowbreak 调\allowbreak 用\allowbreak 。\allowbreak 对\allowbreak 于\allowbreak 由\allowbreak 小\allowbreak 团\allowbreak 队\allowbreak 构\allowbreak 建\allowbreak 的\allowbreak 此\allowbreak 等\allowbreak 规\allowbreak 模\allowbreak 语\allowbreak 料\allowbreak ，\allowbreak 这\allowbreak 两\allowbreak 条\allowbreak 性\allowbreak 质\allowbreak 都\allowbreak 很\allowbreak 重\allowbreak 要\allowbreak ：\allowbreak 它\allowbreak 们\allowbreak 使\allowbreak 采\allowbreak 集\allowbreak 与\allowbreak 分\allowbreak 类\allowbreak 可\allowbreak 在\allowbreak 无\allowbreak 需\allowbreak 人\allowbreak 工\allowbreak 簿\allowbreak 记\allowbreak 的\allowbreak 情\allowbreak 况\allowbreak 下\allowbreak 重\allowbreak 启\allowbreak 。}

\subsection{CAICS{\cjkfont ：\allowbreak 操\allowbreak 作\allowbreak 性\allowbreak 编\allowbreak 码\allowbreak 方\allowbreak 案}}
{\cjkfont 建\allowbreak 筑\allowbreak 施\allowbreak 工\allowbreak 事\allowbreak 故\allowbreak 综\allowbreak 合\allowbreak 致\allowbreak 因\allowbreak 系\allowbreak 统\allowbreak （}CAICS{\cjkfont ）\allowbreak 是\allowbreak 面\allowbreak 向\allowbreak 施\allowbreak 工\allowbreak 事\allowbreak 故\allowbreak 叙\allowbreak 述\allowbreak 的} 5 {\cjkfont 层} 19 {\cjkfont 类\allowbreak 编\allowbreak 码\allowbreak 方\allowbreak 案\allowbreak 。\allowbreak 其\allowbreak 层\allowbreak 级\allowbreak 遵\allowbreak 循\allowbreak 因\allowbreak 果\allowbreak 泄\allowbreak 漏\allowbreak 模\allowbreak 型\allowbreak ：\allowbreak 管\allowbreak 理\allowbreak 层\allowbreak （}M{\cjkfont ）\allowbreak 回\allowbreak 答\allowbreak 为\allowbreak 何\allowbreak 该\allowbreak 工\allowbreak 地\allowbreak 易\allowbreak 出\allowbreak 事\allowbreak 故\allowbreak ；\allowbreak 防\allowbreak 线\allowbreak 层\allowbreak （}F{\cjkfont ）\allowbreak 回\allowbreak 答\allowbreak 为\allowbreak 何\allowbreak 物\allowbreak 理\allowbreak 屏\allowbreak 障\allowbreak 失\allowbreak 效\allowbreak ；\allowbreak 触\allowbreak 发\allowbreak 层\allowbreak （}T{\cjkfont ）\allowbreak 回\allowbreak 答\allowbreak 何\allowbreak 种\allowbreak 物\allowbreak 质\allowbreak 条\allowbreak 件\allowbreak 引\allowbreak 发\allowbreak 事\allowbreak 件\allowbreak ；\allowbreak 人\allowbreak 的\allowbreak 行\allowbreak 为\allowbreak 层\allowbreak （}H{\cjkfont ）\allowbreak 回\allowbreak 答\allowbreak 何\allowbreak 种\allowbreak 不\allowbreak 安\allowbreak 全\allowbreak 行\allowbreak 为\allowbreak 在\allowbreak 前\allowbreak ；\allowbreak 后\allowbreak 果\allowbreak 层\allowbreak （}C{\cjkfont ）\allowbreak 回\allowbreak 答\allowbreak 为\allowbreak 何\allowbreak 后\allowbreak 果\allowbreak 超\allowbreak 出\allowbreak 预\allowbreak 期\allowbreak 。\allowbreak 表} 2 {\cjkfont 列\allowbreak 出\allowbreak 十\allowbreak 九\allowbreak 类\allowbreak 及\allowbreak 其\allowbreak 操\allowbreak 作\allowbreak 定\allowbreak 义\allowbreak 与\allowbreak 典\allowbreak 型\allowbreak 证\allowbreak 据\allowbreak 关\allowbreak 键\allowbreak 词\allowbreak 。}

\textbf{{\cjkfont 表} 2. CAICS {\cjkfont 编\allowbreak 码\allowbreak 方\allowbreak 案\allowbreak ：\allowbreak 操\allowbreak 作\allowbreak 定\allowbreak 义}}

\begingroup\setlength{\tabcolsep}{2pt}\small
\begin{longtable}{>{\raggedright\arraybackslash}p{0.2297\textwidth}>{\raggedright\arraybackslash}p{0.0446\textwidth}>{\raggedright\arraybackslash}p{0.1146\textwidth}>{\raggedright\arraybackslash}p{0.3152\textwidth}>{\raggedright\arraybackslash}p{0.2467\textwidth}}
\hline
{\cjkfont 层} & {\cjkfont 码} & {\cjkfont 类\allowbreak 别} & {\cjkfont 操\allowbreak 作\allowbreak 定\allowbreak 义} & {\cjkfont 典\allowbreak 型\allowbreak 证\allowbreak 据\allowbreak 关\allowbreak 键\allowbreak 词} \\
\hline\endfirsthead
\hline
{\cjkfont 层} & {\cjkfont 码} & {\cjkfont 类\allowbreak 别} & {\cjkfont 操\allowbreak 作\allowbreak 定\allowbreak 义} & {\cjkfont 典\allowbreak 型\allowbreak 证\allowbreak 据\allowbreak 关\allowbreak 键\allowbreak 词} \\
\hline\endhead
\textbf{M {\cjkfont 层} --- {\cjkfont 管\allowbreak 理\allowbreak 根\allowbreak 基}} & M1 & {\cjkfont 现\allowbreak 场\allowbreak 监\allowbreak 管\allowbreak 缺\allowbreak 失} & {\cjkfont 现\allowbreak 场\allowbreak 无\allowbreak 具\allowbreak 备\allowbreak 资\allowbreak 格\allowbreak 的\allowbreak 监\allowbreak 督\allowbreak 人\allowbreak 员\allowbreak ，\allowbreak 或\allowbreak 监\allowbreak 督\allowbreak 未\allowbreak 能\allowbreak 发\allowbreak 现\allowbreak 并\allowbreak 制\allowbreak 止\allowbreak 不\allowbreak 安\allowbreak 全\allowbreak 状\allowbreak 态} & {\cjkfont 无\allowbreak 人\allowbreak 监\allowbreak 督\allowbreak ；\allowbreak 监\allowbreak 理\allowbreak 未\allowbreak 到\allowbreak 场\allowbreak ；\allowbreak 未\allowbreak 予\allowbreak 制\allowbreak 止} \\
 & M2 & {\cjkfont 隐\allowbreak 患\allowbreak 排\allowbreak 查\allowbreak 失\allowbreak 效} & {\cjkfont 隐\allowbreak 患\allowbreak 未\allowbreak 排\allowbreak 查\allowbreak 、\allowbreak 未\allowbreak 整\allowbreak 改\allowbreak ，\allowbreak 或\allowbreak 长\allowbreak 期\allowbreak 存\allowbreak 在} & {\cjkfont 隐\allowbreak 患\allowbreak 未\allowbreak 整\allowbreak 改\allowbreak ；\allowbreak 检\allowbreak 查\allowbreak 走\allowbreak 过\allowbreak 场\allowbreak ；\allowbreak 反\allowbreak 复\allowbreak 上\allowbreak 报\allowbreak 未\allowbreak 处\allowbreak 置} \\
 & M3 & {\cjkfont 安\allowbreak 全\allowbreak 培\allowbreak 训\allowbreak 不\allowbreak 足} & {\cjkfont 未\allowbreak 进\allowbreak 行\allowbreak 入\allowbreak 场\allowbreak 培\allowbreak 训\allowbreak 或\allowbreak 安\allowbreak 全\allowbreak 交\allowbreak 底\allowbreak 流\allowbreak 于\allowbreak 形\allowbreak 式\allowbreak ；\allowbreak 作\allowbreak 业\allowbreak 人\allowbreak 员\allowbreak 缺\allowbreak 乏\allowbreak 必\allowbreak 要\allowbreak 安\allowbreak 全\allowbreak 知\allowbreak 识} & {\cjkfont 未\allowbreak 做\allowbreak 三\allowbreak 级\allowbreak 教\allowbreak 育\allowbreak ；\allowbreak 班\allowbreak 前\allowbreak 未\allowbreak 交\allowbreak 底\allowbreak ；\allowbreak 作\allowbreak 业\allowbreak 人\allowbreak 员\allowbreak 不\allowbreak 知\allowbreak 规\allowbreak 程} \\
 & M4 & {\cjkfont 制\allowbreak 度\allowbreak 规\allowbreak 程\allowbreak 漏\allowbreak 洞} & {\cjkfont 无\allowbreak 安\allowbreak 全\allowbreak 方\allowbreak 案\allowbreak 、\allowbreak 无\allowbreak 应\allowbreak 急\allowbreak 预\allowbreak 案\allowbreak ，\allowbreak 或\allowbreak 现\allowbreak 场\allowbreak 管\allowbreak 理\allowbreak 制\allowbreak 度\allowbreak 不\allowbreak 健\allowbreak 全} & {\cjkfont 无\allowbreak 专\allowbreak 项\allowbreak 施\allowbreak 工\allowbreak 方\allowbreak 案\allowbreak ；\allowbreak 应\allowbreak 急\allowbreak 预\allowbreak 案\allowbreak 缺\allowbreak 失\allowbreak ；\allowbreak 制\allowbreak 度\allowbreak 未\allowbreak 建\allowbreak 立} \\
 & M5 & {\cjkfont 市\allowbreak 场\allowbreak 行\allowbreak 为\allowbreak 违\allowbreak 规} & {\cjkfont 无\allowbreak 资\allowbreak 质\allowbreak 承\allowbreak 揽\allowbreak 、\allowbreak 违\allowbreak 法\allowbreak 分\allowbreak 包\allowbreak 、\allowbreak 转\allowbreak 包\allowbreak 挂\allowbreak 靠\allowbreak ，\allowbreak 或\allowbreak 违\allowbreak 规\allowbreak 压\allowbreak 缩\allowbreak 工\allowbreak 期} & {\cjkfont 施\allowbreak 工\allowbreak 单\allowbreak 位\allowbreak 无\allowbreak 资\allowbreak 质\allowbreak ；\allowbreak 违\allowbreak 法\allowbreak 分\allowbreak 包\allowbreak ；\allowbreak 压\allowbreak 缩\allowbreak 工\allowbreak 期} \\
\textbf{F {\cjkfont 层} --- {\cjkfont 防\allowbreak 线\allowbreak 失\allowbreak 效}} & F1 & {\cjkfont 临\allowbreak 边\allowbreak 洞\allowbreak 口\allowbreak 防\allowbreak 护\allowbreak 缺\allowbreak 失} & {\cjkfont 无\allowbreak 防\allowbreak 护\allowbreak 栏\allowbreak 杆\allowbreak 、\allowbreak 无\allowbreak 安\allowbreak 全\allowbreak 网\allowbreak 、\allowbreak 洞\allowbreak 口\allowbreak 未\allowbreak 盖\allowbreak 或\allowbreak 临\allowbreak 边\allowbreak 未\allowbreak 设\allowbreak 防} & {\cjkfont 井\allowbreak 口\allowbreak 无\allowbreak 护\allowbreak 栏\allowbreak ；\allowbreak 安\allowbreak 全\allowbreak 网\allowbreak 缺\allowbreak 失\allowbreak ；\allowbreak 洞\allowbreak 口\allowbreak 未\allowbreak 覆\allowbreak 盖} \\
 & F2 & {\cjkfont 支\allowbreak 撑\allowbreak 体\allowbreak 系\allowbreak 失\allowbreak 稳} & {\cjkfont 模\allowbreak 板\allowbreak 或\allowbreak 脚\allowbreak 手\allowbreak 架\allowbreak 失\allowbreak 稳\allowbreak 、\allowbreak 支\allowbreak 撑\allowbreak 不\allowbreak 足\allowbreak ，\allowbreak 或\allowbreak 临\allowbreak 时\allowbreak 结\allowbreak 构\allowbreak 失\allowbreak 效} & {\cjkfont 模\allowbreak 板\allowbreak 支\allowbreak 撑\allowbreak 坍\allowbreak 塌\allowbreak ；\allowbreak 脚\allowbreak 手\allowbreak 架\allowbreak 失\allowbreak 稳\allowbreak ；\allowbreak 支\allowbreak 撑\allowbreak 不\allowbreak 足} \\
 & F3 & {\cjkfont 设\allowbreak 备\allowbreak 安\allowbreak 全\allowbreak 装\allowbreak 置\allowbreak 失\allowbreak 效} & {\cjkfont 限\allowbreak 位\allowbreak 器\allowbreak 被\allowbreak 屏\allowbreak 蔽\allowbreak 、\allowbreak 防\allowbreak 坠\allowbreak 器\allowbreak 失\allowbreak 效\allowbreak ，\allowbreak 或\allowbreak 安\allowbreak 全\allowbreak 装\allowbreak 置\allowbreak 损\allowbreak 坏\allowbreak 、\allowbreak 未\allowbreak 安\allowbreak 装} & {\cjkfont 限\allowbreak 位\allowbreak 器\allowbreak 被\allowbreak 屏\allowbreak 蔽\allowbreak ；\allowbreak 防\allowbreak 坠\allowbreak 器\allowbreak 失\allowbreak 效\allowbreak ；\allowbreak 装\allowbreak 置\allowbreak 被\allowbreak 拆\allowbreak 除} \\
 & F4 & {\cjkfont 个\allowbreak 体\allowbreak 防\allowbreak 护\allowbreak 未\allowbreak 使\allowbreak 用} & {\cjkfont 应\allowbreak 按\allowbreak 要\allowbreak 求\allowbreak 佩\allowbreak 戴\allowbreak 的\allowbreak 安\allowbreak 全\allowbreak 带\allowbreak 、\allowbreak 安\allowbreak 全\allowbreak 帽\allowbreak 等\allowbreak 个\allowbreak 体\allowbreak 防\allowbreak 护\allowbreak 装\allowbreak 备\allowbreak 未\allowbreak 使\allowbreak 用} & {\cjkfont 未\allowbreak 系\allowbreak 安\allowbreak 全\allowbreak 带\allowbreak ；\allowbreak 未\allowbreak 戴\allowbreak 安\allowbreak 全\allowbreak 帽\allowbreak ；\allowbreak 无\allowbreak 个\allowbreak 体\allowbreak 防\allowbreak 护} \\
\textbf{T {\cjkfont 层} --- {\cjkfont 触\allowbreak 发\allowbreak 条\allowbreak 件}} & T1 & {\cjkfont 设\allowbreak 备\allowbreak 机\allowbreak 具\allowbreak 故\allowbreak 障} & {\cjkfont 设\allowbreak 备\allowbreak 故\allowbreak 障\allowbreak 、\allowbreak 钢\allowbreak 丝\allowbreak 绳\allowbreak 断\allowbreak 裂\allowbreak 、\allowbreak 机\allowbreak 械\allowbreak 失\allowbreak 灵\allowbreak 或\allowbreak 制\allowbreak 动\allowbreak 失\allowbreak 效} & {\cjkfont 提\allowbreak 升\allowbreak 机\allowbreak 故\allowbreak 障\allowbreak ；\allowbreak 钢\allowbreak 丝\allowbreak 绳\allowbreak 断\allowbreak 裂\allowbreak ；\allowbreak 制\allowbreak 动\allowbreak 失\allowbreak 灵} \\
 & T2 & {\cjkfont 物\allowbreak 料\allowbreak 状\allowbreak 态\allowbreak 异\allowbreak 常} & {\cjkfont 堆\allowbreak 载\allowbreak 超\allowbreak 出\allowbreak 设\allowbreak 计\allowbreak 承\allowbreak 载\allowbreak 、\allowbreak 材\allowbreak 料\allowbreak 缺\allowbreak 陷\allowbreak 或\allowbreak 物\allowbreak 料\allowbreak 失\allowbreak 稳} & {\cjkfont 堆\allowbreak 载\allowbreak 超\allowbreak 限\allowbreak ；\allowbreak 材\allowbreak 料\allowbreak 缺\allowbreak 陷\allowbreak ；\allowbreak 物\allowbreak 料\allowbreak 浸\allowbreak 水\allowbreak 增\allowbreak 重} \\
 & T3 & {\cjkfont 地\allowbreak 质\allowbreak 环\allowbreak 境\allowbreak 突\allowbreak 变} & {\cjkfont 涌\allowbreak 水\allowbreak 、\allowbreak 边\allowbreak 坡\allowbreak 失\allowbreak 稳\allowbreak 、\allowbreak 地\allowbreak 基\allowbreak 沉\allowbreak 降\allowbreak 或\allowbreak 流\allowbreak 沙} & {\cjkfont 突\allowbreak 发\allowbreak 涌\allowbreak 水\allowbreak ；\allowbreak 边\allowbreak 坡\allowbreak 失\allowbreak 稳\allowbreak ；\allowbreak 地\allowbreak 基\allowbreak 沉\allowbreak 降} \\
 & T4 & {\cjkfont 气\allowbreak 象\allowbreak 水\allowbreak 文\allowbreak 触\allowbreak 发} & {\cjkfont 暴\allowbreak 雨\allowbreak 、\allowbreak 大\allowbreak 风\allowbreak 、\allowbreak 冰\allowbreak 雪\allowbreak 、\allowbreak 雷\allowbreak 电\allowbreak 或\allowbreak 连\allowbreak 日\allowbreak 降\allowbreak 雨} & {\cjkfont 暴\allowbreak 雨\allowbreak ；\allowbreak 大\allowbreak 风\allowbreak ；\allowbreak 连\allowbreak 日\allowbreak 降\allowbreak 雨\allowbreak ；\allowbreak 雷\allowbreak 电} \\
\textbf{H {\cjkfont 层} --- {\cjkfont 人\allowbreak 的\allowbreak 行\allowbreak 为}} & H1 & {\cjkfont 违\allowbreak 反\allowbreak 操\allowbreak 作\allowbreak 规\allowbreak 程} & {\cjkfont 未\allowbreak 按\allowbreak 方\allowbreak 案\allowbreak 施\allowbreak 工\allowbreak 、\allowbreak 擅\allowbreak 自\allowbreak 变\allowbreak 更\allowbreak 方\allowbreak 案\allowbreak 或\allowbreak 未\allowbreak 按\allowbreak 规\allowbreak 范\allowbreak 搭\allowbreak 设} & {\cjkfont 未\allowbreak 按\allowbreak 方\allowbreak 案\allowbreak 施\allowbreak 工\allowbreak ；\allowbreak 擅\allowbreak 自\allowbreak 变\allowbreak 更\allowbreak ；\allowbreak 违\allowbreak 规\allowbreak 搭\allowbreak 设} \\
 & H2 & {\cjkfont 违\allowbreak 规\allowbreak 指\allowbreak 挥\allowbreak 决\allowbreak 策} & {\cjkfont 明\allowbreak 知\allowbreak 已\allowbreak 识\allowbreak 别\allowbreak 的\allowbreak 风\allowbreak 险\allowbreak 仍\allowbreak 下\allowbreak 令\allowbreak 或\allowbreak 继\allowbreak 续\allowbreak 作\allowbreak 业} & {\cjkfont 强\allowbreak 令\allowbreak 作\allowbreak 业\allowbreak ；\allowbreak 已\allowbreak 知\allowbreak 风\allowbreak 险\allowbreak 仍\allowbreak 继\allowbreak 续\allowbreak 施\allowbreak 工} \\
 & H3 & {\cjkfont 风\allowbreak 险\allowbreak 辨\allowbreak 识\allowbreak 失\allowbreak 败} & {\cjkfont 未\allowbreak 识\allowbreak 别\allowbreak 危\allowbreak 险\allowbreak ，\allowbreak 或\allowbreak 已\allowbreak 识\allowbreak 别\allowbreak 但\allowbreak 判\allowbreak 断\allowbreak 失\allowbreak 误} & {\cjkfont 未\allowbreak 察\allowbreak 觉\allowbreak 危\allowbreak 险\allowbreak ；\allowbreak 判\allowbreak 断\allowbreak 失\allowbreak 误\allowbreak ；\allowbreak 疏\allowbreak 忽\allowbreak 大\allowbreak 意} \\
 & H4 & {\cjkfont 应\allowbreak 急\allowbreak 行\allowbreak 为\allowbreak 失\allowbreak 当} & {\cjkfont 无\allowbreak 防\allowbreak 护\allowbreak 施\allowbreak 救\allowbreak 、\allowbreak 撤\allowbreak 离\allowbreak 不\allowbreak 及\allowbreak 时\allowbreak ，\allowbreak 或\allowbreak 救\allowbreak 援\allowbreak 方\allowbreak 法\allowbreak 不\allowbreak 当\allowbreak 致\allowbreak 后\allowbreak 果\allowbreak 扩\allowbreak 大} & {\cjkfont 盲\allowbreak 目\allowbreak 施\allowbreak 救\allowbreak ；\allowbreak 救\allowbreak 援\allowbreak 方\allowbreak 法\allowbreak 错\allowbreak 误\allowbreak ；\allowbreak 撤\allowbreak 离\allowbreak 不\allowbreak 及\allowbreak 时} \\
\textbf{C {\cjkfont 层} --- {\cjkfont 后\allowbreak 果\allowbreak 放\allowbreak 大}} & C1 & {\cjkfont 救\allowbreak 援\allowbreak 响\allowbreak 应\allowbreak 迟\allowbreak 缓} & {\cjkfont 救\allowbreak 援\allowbreak 启\allowbreak 动\allowbreak 滞\allowbreak 后\allowbreak 、\allowbreak 应\allowbreak 急\allowbreak 预\allowbreak 案\allowbreak 未\allowbreak 落\allowbreak 实\allowbreak ，\allowbreak 或\allowbreak 二\allowbreak 次\allowbreak 坍\allowbreak 塌\allowbreak 阻\allowbreak 碍\allowbreak 救\allowbreak 援} & {\cjkfont 救\allowbreak 援\allowbreak 不\allowbreak 及\allowbreak 时\allowbreak ；\allowbreak 二\allowbreak 次\allowbreak 坍\allowbreak 塌\allowbreak ；\allowbreak 无\allowbreak 应\allowbreak 急\allowbreak 预\allowbreak 案} \\
 & C2 & {\cjkfont 信\allowbreak 息\allowbreak 迟\allowbreak 报\allowbreak 瞒\allowbreak 报} & {\cjkfont 未\allowbreak 上\allowbreak 报\allowbreak 、\allowbreak 少\allowbreak 报\allowbreak 或\allowbreak 迟\allowbreak 报\allowbreak 事\allowbreak 故\allowbreak 信\allowbreak 息} & {\cjkfont 未\allowbreak 上\allowbreak 报\allowbreak ；\allowbreak 瞒\allowbreak 报\allowbreak 伤\allowbreak 亡\allowbreak ；\allowbreak 迟\allowbreak 报} \\
\hline
\end{longtable}\endgroup
\textit{{\cjkfont 关\allowbreak 键\allowbreak 词\allowbreak 为\allowbreak 语\allowbreak 料\allowbreak 中\allowbreak 中\allowbreak 文\allowbreak 表\allowbreak 述\allowbreak 的\allowbreak 概\allowbreak 括\allowbreak 。\allowbreak 不\allowbreak 依\allowbreak 赖\allowbreak 直\allowbreak 接\allowbreak 引\allowbreak 用\allowbreak 的\allowbreak 编\allowbreak 码\allowbreak 占\allowbreak 比\allowbreak 随\allowbreak 层\allowbreak 而\allowbreak 异}------{\cjkfont 管\allowbreak 理\allowbreak 层} 12.0\%{\cjkfont ，\allowbreak 触\allowbreak 发\allowbreak 层} 3.1\%------{\cjkfont 与} 3.1 {\cjkfont 节\allowbreak 的\allowbreak 三\allowbreak 级\allowbreak 证\allowbreak 据\allowbreak 方\allowbreak 案\allowbreak 一\allowbreak 致\allowbreak 。}}

CAICS {\cjkfont 与\allowbreak 源\allowbreak 自\allowbreak 航\allowbreak 空\allowbreak 领\allowbreak 域\allowbreak 的} HFACS {\cjkfont 改\allowbreak 编\allowbreak （}Shappell and Wiegmann, 2000; Liu et al., 2025a{\cjkfont ）\allowbreak 有\allowbreak 两\allowbreak 个\allowbreak 设\allowbreak 计\allowbreak 差\allowbreak 异\allowbreak ：}F {\cjkfont 层\allowbreak 刻\allowbreak 画\allowbreak 施\allowbreak 工\allowbreak 领\allowbreak 域\allowbreak 普\allowbreak 遍\allowbreak 的\allowbreak 物\allowbreak 理\allowbreak 防\allowbreak 护\allowbreak 失\allowbreak 效\allowbreak ，}C {\cjkfont 层\allowbreak 记\allowbreak 录\allowbreak 迟\allowbreak 报\allowbreak 瞒\allowbreak 报\allowbreak 等\allowbreak 后\allowbreak 果\allowbreak 放\allowbreak 大\allowbreak 因\allowbreak 素\allowbreak 。\allowbreak 方\allowbreak 案\allowbreak 经\allowbreak 泄\allowbreak 漏\allowbreak 模\allowbreak 型\allowbreak 指\allowbreak 导\allowbreak 下\allowbreak 的\allowbreak 语\allowbreak 料\allowbreak 驱\allowbreak 动\allowbreak 归\allowbreak 纳\allowbreak 得\allowbreak 出\allowbreak ：\allowbreak 从\allowbreak 语\allowbreak 料\allowbreak 抽\allowbreak 取\allowbreak 因\allowbreak 果\allowbreak 表\allowbreak 述\allowbreak 、\allowbreak 聚\allowbreak 类\allowbreak 为\allowbreak 候\allowbreak 选\allowbreak 类\allowbreak 别\allowbreak 、\allowbreak 整\allowbreak 合\allowbreak 为\allowbreak 五\allowbreak 层\allowbreak ，\allowbreak 频\allowbreak 次\allowbreak 低\allowbreak 于\allowbreak 下\allowbreak 限\allowbreak 的\allowbreak 类\allowbreak 别\allowbreak 合\allowbreak 并\allowbreak 或\allowbreak 删\allowbreak 除\allowbreak ；\allowbreak 保\allowbreak 留\allowbreak 类\allowbreak 别\allowbreak 均\allowbreak 配\allowbreak 有\allowbreak 操\allowbreak 作\allowbreak 定\allowbreak 义\allowbreak 与\allowbreak 语\allowbreak 料\allowbreak 归\allowbreak 纳\allowbreak 的\allowbreak 证\allowbreak 据\allowbreak 关\allowbreak 键\allowbreak 词\allowbreak 。\allowbreak 我\allowbreak 们\allowbreak 强\allowbreak 调\allowbreak 该\allowbreak 声\allowbreak 明\allowbreak 的\allowbreak 范\allowbreak 围\allowbreak ：\allowbreak 推\allowbreak 导\allowbreak 协\allowbreak 议\allowbreak 是\allowbreak 迭\allowbreak 代\allowbreak 式\allowbreak 工\allowbreak 程\allowbreak 流\allowbreak 程\allowbreak 而\allowbreak 非\allowbreak 正\allowbreak 式\allowbreak 的\allowbreak 内\allowbreak 容\allowbreak 效\allowbreak 度\allowbreak 研\allowbreak 究\allowbreak ，\allowbreak 且\allowbreak 尚\allowbreak 未\allowbreak 在\allowbreak 独\allowbreak 立\allowbreak 数\allowbreak 据\allowbreak 上\allowbreak 与} 24Model{\cjkfont （}Xuecai et al., 2024{\cjkfont ）\allowbreak 、}Accimap {\cjkfont 或} HFACS {\cjkfont 做\allowbreak 对\allowbreak 比\allowbreak 验\allowbreak 证\allowbreak 。\allowbreak 因\allowbreak 此\allowbreak 本\allowbreak 文\allowbreak 将} CAICS {\cjkfont 定\allowbreak 位\allowbreak 为\allowbreak 本\allowbreak 数\allowbreak 据\allowbreak 库\allowbreak 的\allowbreak 操\allowbreak 作\allowbreak 性\allowbreak 方\allowbreak 案}------{\cjkfont 面\allowbreak 向} LLM {\cjkfont 的\allowbreak 有\allowbreak 限\allowbreak 编\allowbreak 码\allowbreak 框\allowbreak 架\allowbreak ，\allowbreak 配\allowbreak 显\allowbreak 式\allowbreak 定\allowbreak 义}------{\cjkfont 而\allowbreak 非\allowbreak 已\allowbreak 经\allowbreak 确\allowbreak 立\allowbreak 的\allowbreak 致\allowbreak 因\allowbreak 理\allowbreak 论\allowbreak 。\allowbreak 其\allowbreak 可\allowbreak 用\allowbreak 性\allowbreak 证\allowbreak 据\allowbreak 包\allowbreak 括} 100 {\cjkfont 条\allowbreak 策\allowbreak 略\allowbreak 样\allowbreak 本\allowbreak 上\allowbreak 的\allowbreak 专\allowbreak 家\allowbreak 标\allowbreak 注\allowbreak 一\allowbreak 致\allowbreak 性\allowbreak （}Fleiss' kappa {\cjkfont 各\allowbreak 层} 0.859--1.000{\cjkfont ，\allowbreak 见} 4.1 {\cjkfont 节\allowbreak ）\allowbreak 。}

\subsection{{\cjkfont 管\allowbreak 道\allowbreak 架\allowbreak 构}}
{\cjkfont 图} 1 {\cjkfont 展\allowbreak 示\allowbreak 总\allowbreak 体\allowbreak 架\allowbreak 构\allowbreak 。\allowbreak 管\allowbreak 道\allowbreak 由\allowbreak 六\allowbreak 个\allowbreak 处\allowbreak 理\allowbreak 阶\allowbreak 段\allowbreak 与\allowbreak 随\allowbreak 后\allowbreak 的\allowbreak 存\allowbreak 储\allowbreak 、\allowbreak 访\allowbreak 问\allowbreak 层\allowbreak 构\allowbreak 成\allowbreak 。}

\begin{figure}[tbp]
\centering
\includegraphics[width=0.92\textwidth,height=0.80\textheight,keepaspectratio]{/home/suntianhao/construction_agent/.claude/worktrees/classifier-mod/docs/paper_v2/figures/fig1_architecture.png}
\end{figure}
\textbf{{\cjkfont 图} 1. CAIEMS {\cjkfont 系\allowbreak 统\allowbreak 架\allowbreak 构}}
\textbf{{\cjkfont 阶\allowbreak 段} 1{\cjkfont ：\allowbreak 多\allowbreak 源\allowbreak 采\allowbreak 集\allowbreak 。}} {\cjkfont 事\allowbreak 故\allowbreak 报\allowbreak 告\allowbreak 采\allowbreak 集\allowbreak 自\allowbreak 公\allowbreak 共\allowbreak 安\allowbreak 全\allowbreak 信\allowbreak 息\allowbreak 门\allowbreak 户\allowbreak （}safehoo.com{\cjkfont 、}esafety.cn{\cjkfont ）\allowbreak 、\allowbreak 省\allowbreak 级\allowbreak 与\allowbreak 国\allowbreak 家\allowbreak 级\allowbreak 应\allowbreak 急\allowbreak 管\allowbreak 理\allowbreak 部\allowbreak 门\allowbreak 网\allowbreak 站\allowbreak 及\allowbreak 公\allowbreak 开\allowbreak 事\allowbreak 故\allowbreak 汇\allowbreak 编\allowbreak ，\allowbreak 遵\allowbreak 循} 3.2 {\cjkfont 节\allowbreak 所\allowbreak 述\allowbreak 的\allowbreak 双\allowbreak 轨\allowbreak 方\allowbreak 案\allowbreak 与\allowbreak 采\allowbreak 集\allowbreak 、\allowbreak 清\allowbreak 洗\allowbreak 、\allowbreak 筛\allowbreak 选\allowbreak 流\allowbreak 程\allowbreak 。\allowbreak 阶\allowbreak 段} 1 {\cjkfont 产\allowbreak 出\allowbreak 表} 1 {\cjkfont 所\allowbreak 报\allowbreak 告\allowbreak 构\allowbreak 成\allowbreak 的\allowbreak 、\allowbreak 已\allowbreak 去\allowbreak 重\allowbreak 的\allowbreak 工\allowbreak 作\allowbreak 语\allowbreak 料} 2,623 {\cjkfont 条\allowbreak 。\allowbreak 中\allowbreak 国\allowbreak 承\allowbreak 包\allowbreak 商\allowbreak 海\allowbreak 外\allowbreak 工\allowbreak 地\allowbreak 事\allowbreak 故\allowbreak （}22 {\cjkfont 条\allowbreak 境\allowbreak 外\allowbreak 记\allowbreak 录\allowbreak ）\allowbreak 保\allowbreak 留\allowbreak 于\allowbreak 库\allowbreak 但\allowbreak 排\allowbreak 除\allowbreak 在\allowbreak 省\allowbreak 级\allowbreak 统\allowbreak 计\allowbreak 之\allowbreak 外\allowbreak 。}

\textbf{{\cjkfont 阶\allowbreak 段} 2{\cjkfont ：}LLM {\cjkfont 文\allowbreak 本\allowbreak 分\allowbreak 类\allowbreak 。}} {\cjkfont 候\allowbreak 选\allowbreak 页\allowbreak 面\allowbreak 异\allowbreak 质\allowbreak 性\allowbreak 显\allowbreak 著\allowbreak ，\allowbreak 且\allowbreak 提\allowbreak 及\allowbreak 事\allowbreak 故\allowbreak 的\allowbreak 页\allowbreak 面\allowbreak 未\allowbreak 必\allowbreak 是\allowbreak 某\allowbreak 起\allowbreak 事\allowbreak 故\allowbreak 的\allowbreak 报\allowbreak 告\allowbreak 。\allowbreak 轻\allowbreak 量} LLM {\cjkfont 分\allowbreak 类\allowbreak 器\allowbreak 以\allowbreak 五\allowbreak 个\allowbreak 标\allowbreak 注\allowbreak 示\allowbreak 例\allowbreak 的\allowbreak 少\allowbreak 样\allowbreak 本\allowbreak 提\allowbreak 示\allowbreak 为\allowbreak 每\allowbreak 条\allowbreak 记\allowbreak 录\allowbreak 分\allowbreak 配\allowbreak 六\allowbreak 种\allowbreak 文\allowbreak 本\allowbreak 类\allowbreak 型\allowbreak 之\allowbreak 一}------\textit{{\cjkfont 单\allowbreak 事\allowbreak 故\allowbreak 报\allowbreak 告}}{\cjkfont 、}\textit{{\cjkfont 多\allowbreak 事\allowbreak 故\allowbreak 汇\allowbreak 编}}{\cjkfont 、}\textit{{\cjkfont 监\allowbreak 管\allowbreak 文\allowbreak 件}}{\cjkfont 、}\textit{{\cjkfont 汇\allowbreak 总\allowbreak 统\allowbreak 计}}{\cjkfont 、}\textit{{\cjkfont 培\allowbreak 训\allowbreak 材\allowbreak 料}}{\cjkfont 或}\textit{{\cjkfont 其\allowbreak 他}}{\cjkfont 。\allowbreak 六\allowbreak 种\allowbreak 类\allowbreak 型\allowbreak 归\allowbreak 为\allowbreak 三\allowbreak 种\allowbreak 处\allowbreak 置\allowbreak ：\allowbreak 单\allowbreak 事\allowbreak 故\allowbreak 报\allowbreak 告\allowbreak 直\allowbreak 接\allowbreak 推\allowbreak 进\allowbreak ；\allowbreak 多\allowbreak 事\allowbreak 故\allowbreak 汇\allowbreak 编\allowbreak 拆\allowbreak 分\allowbreak 为\allowbreak 其\allowbreak 构\allowbreak 成\allowbreak 案\allowbreak 例\allowbreak ；\allowbreak 其\allowbreak 余\allowbreak 四\allowbreak 类\allowbreak 移\allowbreak 出\allowbreak 语\allowbreak 料\allowbreak 。\allowbreak 该\allowbreak 阶\allowbreak 段\allowbreak 剔\allowbreak 除\allowbreak 的\allowbreak 是} 3.2.4 {\cjkfont 节\allowbreak 规\allowbreak 则\allowbreak 闸\allowbreak 门\allowbreak 放\allowbreak 行\allowbreak 的\allowbreak 监\allowbreak 管\allowbreak 通\allowbreak 知\allowbreak 、\allowbreak 汇\allowbreak 总\allowbreak 统\allowbreak 计\allowbreak 、\allowbreak 培\allowbreak 训\allowbreak 材\allowbreak 料\allowbreak 以\allowbreak 及\allowbreak 广\allowbreak 告\allowbreak 与\allowbreak 导\allowbreak 航\allowbreak 文\allowbreak 本\allowbreak 。}3.2 {\cjkfont 节\allowbreak 所\allowbreak 述\allowbreak 的} 895 {\cjkfont 条\allowbreak 数\allowbreak 字\allowbreak 化\allowbreak 记\allowbreak 录\allowbreak 不\allowbreak 经\allowbreak 过\allowbreak 本\allowbreak 阶\allowbreak 段\allowbreak ；\allowbreak 其\allowbreak 向\allowbreak 独\allowbreak 立\allowbreak 案\allowbreak 例\allowbreak 的\allowbreak 切\allowbreak 分\allowbreak 由\allowbreak 该\allowbreak 节\allowbreak 所\allowbreak 述\allowbreak 的\allowbreak 文\allowbreak 献\allowbreak 级\allowbreak 切\allowbreak 分\allowbreak 路\allowbreak 径\allowbreak 完\allowbreak 成\allowbreak ，\allowbreak 在\allowbreak 记\allowbreak 录\allowbreak 进\allowbreak 入\allowbreak 语\allowbreak 料\allowbreak 之\allowbreak 前\allowbreak 即\allowbreak 实\allowbreak 现\allowbreak 了\allowbreak 同\allowbreak 样\allowbreak 的\allowbreak 区\allowbreak 分\allowbreak 。\allowbreak 需\allowbreak 说\allowbreak 明\allowbreak 的\allowbreak 是\allowbreak 分\allowbreak 类\allowbreak 器\allowbreak 并\allowbreak 不\allowbreak 完\allowbreak 美\allowbreak ：\allowbreak 一\allowbreak 致\allowbreak 性\allowbreak 研\allowbreak 究\allowbreak 的} 300 {\cjkfont 条\allowbreak 抽\allowbreak 样\allowbreak 中\allowbreak 仍\allowbreak 有} 11 {\cjkfont 条\allowbreak 在\allowbreak 标\allowbreak 注\allowbreak 时\allowbreak 因\allowbreak 含\allowbreak 多\allowbreak 起\allowbreak 事\allowbreak 故\allowbreak 被\allowbreak 排\allowbreak 除\allowbreak ，\allowbreak 说\allowbreak 明\allowbreak 文\allowbreak 本\allowbreak 分\allowbreak 类\allowbreak 存\allowbreak 在\allowbreak 残\allowbreak 留\allowbreak 错\allowbreak 误\allowbreak ，\allowbreak 当\allowbreak 前\allowbreak 管\allowbreak 道\allowbreak 尚\allowbreak 未\allowbreak 对\allowbreak 其\allowbreak 完\allowbreak 全\allowbreak 审\allowbreak 计\allowbreak ；\allowbreak 多\allowbreak 事\allowbreak 故\allowbreak 拆\allowbreak 分\allowbreak 规\allowbreak 则\allowbreak 仍\allowbreak 部\allowbreak 分\allowbreak 依\allowbreak 赖\allowbreak 启\allowbreak 发\allowbreak 式\allowbreak 。}

\textbf{{\cjkfont 阶\allowbreak 段} 3{\cjkfont ：\allowbreak 基\allowbreak 础\allowbreak 字\allowbreak 段\allowbreak 抽\allowbreak 取\allowbreak （}Step 1{\cjkfont ）\allowbreak 。}} {\cjkfont 对\allowbreak 每\allowbreak 条\allowbreak 通\allowbreak 过\allowbreak 分\allowbreak 类\allowbreak 的\allowbreak 记\allowbreak 录\allowbreak ，}LLM {\cjkfont 在\allowbreak 严\allowbreak 格} JSON {\cjkfont 工\allowbreak 具\allowbreak 模\allowbreak 式\allowbreak 下\allowbreak 抽\allowbreak 取} 33 {\cjkfont 字\allowbreak 段\allowbreak 基\allowbreak 础\allowbreak 信\allowbreak 息\allowbreak 向\allowbreak 量\allowbreak 。\allowbreak 评\allowbreak 测\allowbreak 运\allowbreak 行\allowbreak 中} 100\% {\cjkfont 记\allowbreak 录\allowbreak 获\allowbreak 得\allowbreak 符\allowbreak 合} schema {\cjkfont 的\allowbreak 输\allowbreak 出\allowbreak 。\allowbreak 抽\allowbreak 取\allowbreak 器\allowbreak 同\allowbreak 时\allowbreak 产\allowbreak 出\allowbreak 规\allowbreak 范\allowbreak 化\allowbreak 标\allowbreak 题\allowbreak 摘\allowbreak 要\allowbreak 用\allowbreak 于\allowbreak 搜\allowbreak 索\allowbreak 增\allowbreak 强\allowbreak 。}

\textbf{{\cjkfont 阶\allowbreak 段} 4{\cjkfont ：}CAICS {\cjkfont 致\allowbreak 因\allowbreak 编\allowbreak 码\allowbreak （}Step 2{\cjkfont ）\allowbreak 。}} {\cjkfont 编\allowbreak 码\allowbreak 器\allowbreak 应\allowbreak 用\allowbreak 两\allowbreak 步\allowbreak 判\allowbreak 断\allowbreak 流\allowbreak 程}------{\cjkfont 显\allowbreak 性\allowbreak 匹\allowbreak 配\allowbreak 后\allowbreak 授\allowbreak 权\allowbreak 推\allowbreak 断}------{\cjkfont 并\allowbreak 为\allowbreak 每\allowbreak 个\allowbreak 编\allowbreak 码\allowbreak 记\allowbreak 录\allowbreak 证\allowbreak 据\allowbreak 层\allowbreak 级\allowbreak 、\allowbreak 支\allowbreak 撑\allowbreak 文\allowbreak 本\allowbreak 与\allowbreak 推\allowbreak 理\allowbreak 链\allowbreak （}3.1 {\cjkfont 节\allowbreak ）\allowbreak 。\allowbreak 图} 2 {\cjkfont 以\allowbreak 一\allowbreak 条\allowbreak 真\allowbreak 实\allowbreak 记\allowbreak 录\allowbreak 的\allowbreak 示\allowbreak 例\allowbreak 展\allowbreak 示\allowbreak 五\allowbreak 层\allowbreak 结\allowbreak 构\allowbreak 。}

\begin{figure}[tbp]
\centering
\includegraphics[width=0.92\textwidth,height=0.80\textheight,keepaspectratio]{/home/suntianhao/construction_agent/.claude/worktrees/classifier-mod/docs/paper_v2/figures/fig2_caics_tree.png}
\end{figure}
\textbf{{\cjkfont 图} 2. CAICS {\cjkfont 五\allowbreak 层\allowbreak 结\allowbreak 构\allowbreak 与\allowbreak 示\allowbreak 例\allowbreak 证\allowbreak 据\allowbreak 链}}
\textbf{{\cjkfont 阶\allowbreak 段} 5{\cjkfont ：}LLM {\cjkfont 优\allowbreak 先\allowbreak 的} MECE {\cjkfont 枚\allowbreak 举\allowbreak 标\allowbreak 准\allowbreak 化\allowbreak 。}} {\cjkfont 五\allowbreak 个\allowbreak 自\allowbreak 由\allowbreak 文\allowbreak 本\allowbreak 字\allowbreak 段\allowbreak （\allowbreak 工\allowbreak 序\allowbreak 、\allowbreak 设\allowbreak 备\allowbreak 、\allowbreak 材\allowbreak 料\allowbreak 、\allowbreak 工\allowbreak 种\allowbreak 、\allowbreak 建\allowbreak 筑\allowbreak 结\allowbreak 构\allowbreak ）\allowbreak 经} LLM {\cjkfont 优\allowbreak 先\allowbreak 流\allowbreak 程\allowbreak 规\allowbreak 范\allowbreak 化\allowbreak ：}LLM {\cjkfont 审\allowbreak 视\allowbreak 字\allowbreak 段\allowbreak 取\allowbreak 值\allowbreak 分\allowbreak 布\allowbreak 、\allowbreak 设\allowbreak 计} MECE {\cjkfont 框\allowbreak 架\allowbreak 、\allowbreak 映\allowbreak 射\allowbreak 全\allowbreak 部\allowbreak 观\allowbreak 察\allowbreak 值\allowbreak 。\allowbreak 表} 3 {\cjkfont 报\allowbreak 告\allowbreak 所\allowbreak 得\allowbreak 枚\allowbreak 举\allowbreak ，\allowbreak 并\allowbreak 显\allowbreak 式\allowbreak 给\allowbreak 出\allowbreak 非\allowbreak 空\allowbreak 分\allowbreak 母\allowbreak 。}

\textbf{{\cjkfont 表} 3. LLM {\cjkfont 优\allowbreak 先\allowbreak 流\allowbreak 程\allowbreak 产\allowbreak 出\allowbreak 的} MECE {\cjkfont 枚\allowbreak 举}}

\begingroup\setlength{\tabcolsep}{2pt}\small
\begin{longtable}{>{\raggedright\arraybackslash}p{0.1553\textwidth}>{\raggedright\arraybackslash}p{0.0681\textwidth}>{\raggedright\arraybackslash}p{0.2126\textwidth}>{\raggedright\arraybackslash}p{0.3346\textwidth}>{\raggedright\arraybackslash}p{0.1802\textwidth}}
\hline
{\cjkfont 字\allowbreak 段} & {\cjkfont 类\allowbreak 别\allowbreak 数} & {\cjkfont 非\allowbreak 空\allowbreak 记\allowbreak 录\allowbreak （\allowbreak 占} 2,623{\cjkfont ）} & {\cjkfont 最\allowbreak 大\allowbreak 类\allowbreak 别\allowbreak （\allowbreak 占\allowbreak 非\allowbreak 空\allowbreak ）} & {\cjkfont 残\allowbreak 余}"{\cjkfont 其\allowbreak 他}"{\cjkfont （\allowbreak 占\allowbreak 非\allowbreak 空\allowbreak ）} \\
\hline\endfirsthead
\hline
{\cjkfont 字\allowbreak 段} & {\cjkfont 类\allowbreak 别\allowbreak 数} & {\cjkfont 非\allowbreak 空\allowbreak 记\allowbreak 录\allowbreak （\allowbreak 占} 2,623{\cjkfont ）} & {\cjkfont 最\allowbreak 大\allowbreak 类\allowbreak 别\allowbreak （\allowbreak 占\allowbreak 非\allowbreak 空\allowbreak ）} & {\cjkfont 残\allowbreak 余}"{\cjkfont 其\allowbreak 他}"{\cjkfont （\allowbreak 占\allowbreak 非\allowbreak 空\allowbreak ）} \\
\hline\endhead
{\cjkfont 工\allowbreak 序\allowbreak （}WT{\cjkfont ）} & 11 & 2,202{\cjkfont （}83.9\%{\cjkfont ）} & WT02 {\cjkfont 主\allowbreak 体\allowbreak 结\allowbreak 构\allowbreak 工\allowbreak 程\allowbreak （}397{\cjkfont ，}18.0\%{\cjkfont ）} & 2.8\% \\
{\cjkfont 涉\allowbreak 事\allowbreak 设\allowbreak 备\allowbreak （}EQP{\cjkfont ）} & 13 & 1,373{\cjkfont （}52.3\%{\cjkfont ）} & EQP01 {\cjkfont 起\allowbreak 重\allowbreak 机\allowbreak 械\allowbreak （}335{\cjkfont ，}24.4\%{\cjkfont ）} & 1.9\% \\
{\cjkfont 涉\allowbreak 事\allowbreak 材\allowbreak 料\allowbreak （}MAT{\cjkfont ）} & 7 & 1,465{\cjkfont （}55.9\%{\cjkfont ）} & MAT01 {\cjkfont 混\allowbreak 凝\allowbreak 土\allowbreak 及\allowbreak 水\allowbreak 泥\allowbreak 制\allowbreak 品\allowbreak （}423{\cjkfont ，}28.9\%{\cjkfont ）} & 11.1\% \\
{\cjkfont 工\allowbreak 种\allowbreak （}WKR{\cjkfont ）} & 13 & 1,905{\cjkfont （}72.6\%{\cjkfont ）} & WKR01 {\cjkfont 通\allowbreak 用\allowbreak 劳\allowbreak 务\allowbreak 工\allowbreak （}928{\cjkfont ，}48.7\%{\cjkfont ）} & 6.3\% \\
{\cjkfont 建\allowbreak 筑\allowbreak 结\allowbreak 构\allowbreak （}BLD{\cjkfont ）} & 12 & 359{\cjkfont （}13.7\%{\cjkfont ）} & BLD01 {\cjkfont 钢\allowbreak 结\allowbreak 构\allowbreak （}127{\cjkfont ，}35.4\%{\cjkfont ）} & 3.6\% \\
\hline
\end{longtable}\endgroup
\textbf{{\cjkfont 阶\allowbreak 段} 6{\cjkfont ：\allowbreak 存\allowbreak 储\allowbreak 与\allowbreak 访\allowbreak 问\allowbreak 。}} {\cjkfont 全\allowbreak 部\allowbreak 产\allowbreak 物\allowbreak 存\allowbreak 储\allowbreak 于\allowbreak 含\allowbreak 九\allowbreak 张\allowbreak 表\allowbreak 的} SQLite {\cjkfont 数\allowbreak 据\allowbreak 库\allowbreak 。\allowbreak 交\allowbreak 互\allowbreak 式} Web {\cjkfont 平\allowbreak 台\allowbreak 与} REST API {\cjkfont 开\allowbreak 放\allowbreak 访\allowbreak 问\allowbreak ：\allowbreak 用\allowbreak 户\allowbreak 可\allowbreak 按\allowbreak 事\allowbreak 故\allowbreak 类\allowbreak 型\allowbreak 、\allowbreak 等\allowbreak 级\allowbreak 、\allowbreak 地\allowbreak 区\allowbreak 、\allowbreak 工\allowbreak 序\allowbreak 与\allowbreak 致\allowbreak 因\allowbreak 编\allowbreak 码\allowbreak 筛\allowbreak 选\allowbreak ，\allowbreak 并\allowbreak 查\allowbreak 看\allowbreak 每\allowbreak 条\allowbreak 标\allowbreak 注\allowbreak 背\allowbreak 后\allowbreak 的\allowbreak 证\allowbreak 据\allowbreak （\allowbreak 证\allowbreak 据\allowbreak 层\allowbreak 级\allowbreak 分\allowbreak 开\allowbreak 展\allowbreak 示\allowbreak ）\allowbreak 。\allowbreak 数\allowbreak 据\allowbreak 库\allowbreak 与\allowbreak 代\allowbreak 码\allowbreak 于\allowbreak 论\allowbreak 文\allowbreak 发\allowbreak 表\allowbreak 后\allowbreak 开\allowbreak 放\allowbreak 检\allowbreak 查\allowbreak ；\allowbreak 此\allowbreak 前\allowbreak 可\allowbreak 向\allowbreak 作\allowbreak 者\allowbreak 索\allowbreak 取\allowbreak 。}

\subsection{{\cjkfont 多\allowbreak 模\allowbreak 态\allowbreak 扩\allowbreak 展}}
{\cjkfont 文\allowbreak 本\allowbreak 之\allowbreak 外\allowbreak ，\allowbreak 管\allowbreak 道\allowbreak 通\allowbreak 过\allowbreak 三\allowbreak 个\allowbreak 展\allowbreak 示\allowbreak 型\allowbreak 模\allowbreak 块\allowbreak 增\allowbreak 强\allowbreak 记\allowbreak 录\allowbreak ：\allowbreak 聚\allowbreak 合\allowbreak 搜\allowbreak 索} API {\cjkfont 的\allowbreak 搜\allowbreak 索\allowbreak 增\allowbreak 强\allowbreak （}1,636 {\cjkfont 条\allowbreak ，}62.4\%{\cjkfont ）\allowbreak 、\allowbreak 零\allowbreak 样\allowbreak 本} CLIP {\cjkfont 分\allowbreak 类\allowbreak 器\allowbreak 过\allowbreak 滤\allowbreak 的\allowbreak 图\allowbreak 片\allowbreak 采\allowbreak 集\allowbreak （}226 {\cjkfont 张\allowbreak 高\allowbreak 于} 0.5 {\cjkfont 相\allowbreak 似\allowbreak 度\allowbreak 阈\allowbreak 值\allowbreak ）\allowbreak 、\allowbreak 以\allowbreak 及} GLM-5V-Turbo {\cjkfont 的\allowbreak 两\allowbreak 阶\allowbreak 段} VLM {\cjkfont 描\allowbreak 述\allowbreak （}282 {\cjkfont 条\allowbreak 分\allowbreak 析\allowbreak 覆\allowbreak 盖} 176 {\cjkfont 起\allowbreak 事\allowbreak 故\allowbreak ）\allowbreak 。\allowbreak 这\allowbreak 些\allowbreak 模\allowbreak 块\allowbreak 是\allowbreak 扩\allowbreak 展\allowbreak 而\allowbreak 非\allowbreak 经\allowbreak 评\allowbreak 估\allowbreak 的\allowbreak 贡\allowbreak 献\allowbreak ：\allowbreak 图\allowbreak 片\allowbreak 源\allowbreak 头\allowbreak 质\allowbreak 量\allowbreak 约\allowbreak 束} VLM {\cjkfont 分\allowbreak 析\allowbreak 的\allowbreak 贡\allowbreak 献\allowbreak ，\allowbreak 描\allowbreak 述\allowbreak 作\allowbreak 为\allowbreak 补\allowbreak 充\allowbreak 材\allowbreak 料\allowbreak 呈\allowbreak 现\allowbreak ，\allowbreak 而\allowbreak 非\allowbreak 经\allowbreak 验\allowbreak 证\allowbreak 标\allowbreak 注\allowbreak 。}

\subsection{{\cjkfont 数\allowbreak 据\allowbreak 库\allowbreak 统\allowbreak 计}}
{\cjkfont 表} 4 {\cjkfont 汇\allowbreak 总\allowbreak 数\allowbreak 据\allowbreak 库\allowbreak ，\allowbreak 全\allowbreak 部\allowbreak 计\allowbreak 数\allowbreak 与\allowbreak 总\allowbreak 量\allowbreak 闭\allowbreak 合\allowbreak 。}

\textbf{{\cjkfont 表} 4. {\cjkfont 数\allowbreak 据\allowbreak 库\allowbreak 统\allowbreak 计}}

\begingroup\setlength{\tabcolsep}{2pt}\small
\begin{longtable}{>{\raggedright\arraybackslash}p{0.1284\textwidth}>{\raggedright\arraybackslash}p{0.8514\textwidth}}
\hline
{\cjkfont 维\allowbreak 度} & {\cjkfont 数\allowbreak 值} \\
\hline\endfirsthead
\hline
{\cjkfont 维\allowbreak 度} & {\cjkfont 数\allowbreak 值} \\
\hline\endhead
{\cjkfont 事\allowbreak 故\allowbreak 记\allowbreak 录} & 2,623 \\
{\cjkfont 时\allowbreak 间\allowbreak 跨\allowbreak 度} & 1978-08 {\cjkfont 至} 2026-06{\cjkfont （}1991--2004 {\cjkfont 无\allowbreak 记\allowbreak 录\allowbreak ）} \\
{\cjkfont 事\allowbreak 故\allowbreak 类\allowbreak 型} & 12 {\cjkfont 类\allowbreak ，}2,617 {\cjkfont 条\allowbreak 有\allowbreak 标\allowbreak 注\allowbreak （\allowbreak 坍\allowbreak 塌} 1,092{\cjkfont ；\allowbreak 高\allowbreak 处\allowbreak 坠\allowbreak 落} 570{\cjkfont ；\allowbreak 起\allowbreak 重\allowbreak 伤\allowbreak 害} 225{\cjkfont ；\allowbreak 土\allowbreak 方\allowbreak 坍\allowbreak 塌} 222{\cjkfont ；\allowbreak 触\allowbreak 电} 160{\cjkfont ；\allowbreak 物\allowbreak 体\allowbreak 打\allowbreak 击} 119{\cjkfont ；\allowbreak 其\allowbreak 他} 50{\cjkfont ；\allowbreak 火\allowbreak 灾\allowbreak 爆\allowbreak 炸} 47{\cjkfont ；\allowbreak 机\allowbreak 械\allowbreak 伤\allowbreak 害} 46{\cjkfont ；\allowbreak 中\allowbreak 毒\allowbreak 窒\allowbreak 息} 41{\cjkfont ；\allowbreak 车\allowbreak 辆\allowbreak 伤\allowbreak 害} 38{\cjkfont ；\allowbreak 淹\allowbreak 溺} 7{\cjkfont ）\allowbreak ；}6 {\cjkfont 条\allowbreak 未\allowbreak 标\allowbreak 注} \\
{\cjkfont 事\allowbreak 故\allowbreak 等\allowbreak 级} & {\cjkfont 一\allowbreak 般} 1,889{\cjkfont ；\allowbreak 较\allowbreak 大} 672{\cjkfont ；\allowbreak 重\allowbreak 大} 53{\cjkfont ；\allowbreak 特\allowbreak 别\allowbreak 重\allowbreak 大} 8{\cjkfont ；}1 {\cjkfont 条\allowbreak 未\allowbreak 标\allowbreak 注} \\
{\cjkfont 伤\allowbreak 亡} & {\cjkfont 死\allowbreak 亡} 5,894{\cjkfont ；\allowbreak 受\allowbreak 伤} 5,796 \\
{\cjkfont 致\allowbreak 因\allowbreak 标\allowbreak 注} & 16,323 {\cjkfont 条\allowbreak （}explicit 14,784 / inferred 1,117 / prior 422{\cjkfont ）} \\
{\cjkfont 搜\allowbreak 索\allowbreak 增\allowbreak 强\allowbreak 记\allowbreak 录} & 1,636 \\
CLIP {\cjkfont 过\allowbreak 滤\allowbreak 图\allowbreak 片} & 226 \\
VLM {\cjkfont 分\allowbreak 析} & 282 {\cjkfont 条} / 176 {\cjkfont 起\allowbreak 事\allowbreak 故} \\
\hline
\end{longtable}\endgroup
{\cjkfont 图} 3 {\cjkfont 展\allowbreak 示\allowbreak 十\allowbreak 九\allowbreak 个} CAICS {\cjkfont 类\allowbreak 别\allowbreak 的\allowbreak 标\allowbreak 注\allowbreak 分\allowbreak 布\allowbreak ；\allowbreak 高\allowbreak 频\allowbreak 编\allowbreak 码\allowbreak （}M1 {\cjkfont 监\allowbreak 管\allowbreak 缺\allowbreak 失\allowbreak 、}H1 {\cjkfont 违\allowbreak 章\allowbreak 操\allowbreak 作\allowbreak 、}H3 {\cjkfont 风\allowbreak 险\allowbreak 辨\allowbreak 识\allowbreak 失\allowbreak 败\allowbreak ）\allowbreak 部\allowbreak 分\allowbreak 反\allowbreak 映\allowbreak 生\allowbreak 产\allowbreak 提\allowbreak 示\allowbreak 授\allowbreak 权\allowbreak 的\allowbreak 先\allowbreak 验\allowbreak 推\allowbreak 断\allowbreak ，\allowbreak 这\allowbreak 也\allowbreak 是\allowbreak 总\allowbreak 量\allowbreak 旁\allowbreak 并\allowbreak 列\allowbreak 分\allowbreak 层\allowbreak 细\allowbreak 目\allowbreak 的\allowbreak 原\allowbreak 因\allowbreak 。}

\begin{figure}[tbp]
\centering
\includegraphics[width=0.92\textwidth,height=0.80\textheight,keepaspectratio]{/home/suntianhao/construction_agent/.claude/worktrees/classifier-mod/docs/paper_v2/figures/fig3_code_dist.png}
\end{figure}
\textbf{{\cjkfont 图} 3. {\cjkfont 十\allowbreak 九\allowbreak 个} CAICS {\cjkfont 类\allowbreak 别\allowbreak 的\allowbreak 致\allowbreak 因\allowbreak 标\allowbreak 注\allowbreak 分\allowbreak 布}}
\section{{\cjkfont 实\allowbreak 验\allowbreak 评\allowbreak 估}}
{\cjkfont 评\allowbreak 估\allowbreak 以\allowbreak 明\allowbreak 确\allowbreak 声\allowbreak 明\allowbreak 口\allowbreak 径\allowbreak 回\allowbreak 答\allowbreak 三\allowbreak 个\allowbreak 问\allowbreak 题\allowbreak ：\allowbreak 管\allowbreak 道\allowbreak 在\allowbreak 其\allowbreak 实\allowbreak 际\allowbreak 成\allowbreak 本\allowbreak 约\allowbreak 束\allowbreak 下\allowbreak 应\allowbreak 采\allowbreak 用\allowbreak 哪\allowbreak 个\allowbreak 骨\allowbreak 干\allowbreak 模\allowbreak 型\allowbreak （}4.3 {\cjkfont 节\allowbreak ）\allowbreak ；\allowbreak 致\allowbreak 因\allowbreak 编\allowbreak 码\allowbreak 提\allowbreak 示\allowbreak 应\allowbreak 如\allowbreak 何\allowbreak 措\allowbreak 辞\allowbreak （}4.4 {\cjkfont 节\allowbreak ）\allowbreak ；}LLM {\cjkfont 的\allowbreak 基\allowbreak 础\allowbreak 字\allowbreak 段\allowbreak 标\allowbreak 注\allowbreak 是\allowbreak 否\allowbreak 与\allowbreak 受\allowbreak 训\allowbreak 人\allowbreak 工\allowbreak 标\allowbreak 注\allowbreak 者\allowbreak 一\allowbreak 致\allowbreak （}4.5 {\cjkfont 节\allowbreak ）\allowbreak 。}4.1 {\cjkfont 与} 4.2 {\cjkfont 节\allowbreak 定\allowbreak 义\allowbreak 数\allowbreak 据\allowbreak 集\allowbreak 与\allowbreak 指\allowbreak 标\allowbreak 。}

\subsection{{\cjkfont 数\allowbreak 据\allowbreak 集\allowbreak 与\allowbreak 抽\allowbreak 样}}
\textbf{{\cjkfont 全\allowbreak 量\allowbreak 语\allowbreak 料\allowbreak 。}} {\cjkfont 管\allowbreak 道\allowbreak 产\allowbreak 出} 2,623 {\cjkfont 条\allowbreak 单\allowbreak 事\allowbreak 故\allowbreak 记\allowbreak 录\allowbreak ，\allowbreak 跨} 1978--2026 {\cjkfont 年\allowbreak ，\allowbreak 含} 33 {\cjkfont 字\allowbreak 段\allowbreak 与} 16,323 {\cjkfont 条\allowbreak 分\allowbreak 层\allowbreak 致\allowbreak 因\allowbreak 标\allowbreak 注\allowbreak 。\allowbreak 语\allowbreak 料\allowbreak 由\allowbreak 网\allowbreak 络\allowbreak 采\allowbreak 集\allowbreak 轨\allowbreak （}1,728 {\cjkfont 条\allowbreak ）\allowbreak 与\allowbreak 文\allowbreak 献\allowbreak 数\allowbreak 字\allowbreak 化\allowbreak 轨\allowbreak （}895 {\cjkfont 条\allowbreak ）\allowbreak 构\allowbreak 成\allowbreak ；}3.2 {\cjkfont 节\allowbreak 报\allowbreak 告\allowbreak 其\allowbreak 构\allowbreak 成\allowbreak 与\allowbreak 构\allowbreak 建\allowbreak 方\allowbreak 式\allowbreak ，\allowbreak 包\allowbreak 括} 1991--2004 {\cjkfont 年\allowbreak 的\allowbreak 断\allowbreak 代\allowbreak 。}

\textbf{{\cjkfont 银\allowbreak 标\allowbreak 准\allowbreak 。}} 972 {\cjkfont 条\allowbreak 银\allowbreak 标\allowbreak 准\allowbreak 由\allowbreak 管\allowbreak 道\allowbreak 版\allowbreak 本\allowbreak 一\allowbreak 致\allowbreak 构\allowbreak 建\allowbreak ，\allowbreak 而\allowbreak 非\allowbreak 独\allowbreak 立\allowbreak 人\allowbreak 工\allowbreak 重\allowbreak 标\allowbreak 注\allowbreak ：\allowbreak 基\allowbreak 础\allowbreak 字\allowbreak 段\allowbreak 在\allowbreak 两\allowbreak 版\allowbreak 本\allowbreak 一\allowbreak 致\allowbreak 时\allowbreak 直\allowbreak 接\allowbreak 采\allowbreak 纳\allowbreak 、\allowbreak 不\allowbreak 一\allowbreak 致\allowbreak 时\allowbreak 规\allowbreak 则\allowbreak 裁\allowbreak 决\allowbreak ；\allowbreak 致\allowbreak 因\allowbreak 维\allowbreak 度\allowbreak 标\allowbreak 签\allowbreak 取\allowbreak 两\allowbreak 版\allowbreak 本\allowbreak 输\allowbreak 出\allowbreak 的\allowbreak 合\allowbreak 并\allowbreak 去\allowbreak 重\allowbreak 。\allowbreak 我\allowbreak 们\allowbreak 显\allowbreak 式\allowbreak 陈\allowbreak 述\allowbreak 这\allowbreak 一\allowbreak 构\allowbreak 建\allowbreak 方\allowbreak 式\allowbreak ，\allowbreak 因\allowbreak 为\allowbreak 它\allowbreak 界\allowbreak 定\allowbreak 了\allowbreak 银\allowbreak 标\allowbreak 准\allowbreak 能\allowbreak 认\allowbreak 证\allowbreak 什\allowbreak 么\allowbreak ：\allowbreak 其\allowbreak 可\allowbreak 靠\allowbreak 程\allowbreak 度\allowbreak 取\allowbreak 决\allowbreak 于\allowbreak 两\allowbreak 个\allowbreak 来\allowbreak 源\allowbreak 版\allowbreak 本\allowbreak 的\allowbreak 可\allowbreak 靠\allowbreak 程\allowbreak 度\allowbreak ，\allowbreak 它\allowbreak 不\allowbreak 是\allowbreak 独\allowbreak 立\allowbreak 人\allowbreak 工\allowbreak 金\allowbreak 标\allowbreak 准\allowbreak 。\allowbreak 银\allowbreak 标\allowbreak 准\allowbreak 中\allowbreak 致\allowbreak 因\allowbreak 维\allowbreak 度\allowbreak 为\allowbreak 空\allowbreak 表\allowbreak 示\allowbreak 该\allowbreak 维\allowbreak 度\allowbreak 在\allowbreak 原\allowbreak 文\allowbreak 未\allowbreak 提\allowbreak 及\allowbreak 。\allowbreak 一\allowbreak 个} 100 {\cjkfont 条\allowbreak 子\allowbreak 集\allowbreak 获\allowbreak 得\allowbreak 了\allowbreak 完\allowbreak 全\allowbreak 独\allowbreak 立\allowbreak 的\allowbreak 人\allowbreak 工\allowbreak 验\allowbreak 证\allowbreak ：\allowbreak 三\allowbreak 位\allowbreak 施\allowbreak 工\allowbreak 安\allowbreak 全\allowbreak 领\allowbreak 域\allowbreak 专\allowbreak 家\allowbreak 在} LabelStudio {\cjkfont 平\allowbreak 台\allowbreak 上\allowbreak 标\allowbreak 注} CAICS {\cjkfont 编\allowbreak 码\allowbreak ，}Fleiss' kappa {\cjkfont 在\allowbreak 五\allowbreak 层\allowbreak 上\allowbreak 介\allowbreak 于} 0.859 {\cjkfont 与} 1.000{\cjkfont ；\allowbreak 该\allowbreak 子\allowbreak 集\allowbreak 真\allowbreak 值\allowbreak 为\allowbreak 三\allowbreak 专\allowbreak 家\allowbreak 多\allowbreak 数\allowbreak 投\allowbreak 票\allowbreak ，\allowbreak 平\allowbreak 均\allowbreak 每\allowbreak 条} 5.1 {\cjkfont 个\allowbreak 编\allowbreak 码\allowbreak 。}

\textbf{{\cjkfont 评\allowbreak 测\allowbreak 样\allowbreak 本\allowbreak 。}}

1. \textit{{\cjkfont 模\allowbreak 型\allowbreak 比\allowbreak 较\allowbreak 样\allowbreak 本}}{\cjkfont （}n=200{\cjkfont ）\allowbreak 。\allowbreak 按\allowbreak 事\allowbreak 故\allowbreak 类\allowbreak 型\allowbreak 分\allowbreak 层\allowbreak （}11 {\cjkfont 类\allowbreak ，\allowbreak 每\allowbreak 类\allowbreak 至\allowbreak 少} 3 {\cjkfont 条\allowbreak ）\allowbreak 从\allowbreak 银\allowbreak 标\allowbreak 准\allowbreak 抽\allowbreak 取\allowbreak 。} 2. \textit{{\cjkfont 策\allowbreak 略\allowbreak 比\allowbreak 较\allowbreak 样\allowbreak 本}}{\cjkfont （}n=100{\cjkfont ）\allowbreak 。\allowbreak 从\allowbreak 银\allowbreak 标\allowbreak 准\allowbreak 抽\allowbreak 取\allowbreak ，\allowbreak 覆\allowbreak 盖\allowbreak 全\allowbreak 部} 12 {\cjkfont 种\allowbreak 事\allowbreak 故\allowbreak 类\allowbreak 型\allowbreak 与} 19 {\cjkfont 个} CAICS {\cjkfont 类\allowbreak 别\allowbreak ；\allowbreak 致\allowbreak 因\allowbreak 真\allowbreak 值\allowbreak 为\allowbreak 上\allowbreak 述\allowbreak 三\allowbreak 专\allowbreak 家\allowbreak 多\allowbreak 数\allowbreak 投\allowbreak 票\allowbreak 。}\textbf{{\cjkfont 口\allowbreak 径\allowbreak 声\allowbreak 明}}{\cjkfont ：\allowbreak 该\allowbreak 样\allowbreak 本\allowbreak 承\allowbreak 担\allowbreak 了\allowbreak 提\allowbreak 示\allowbreak 开\allowbreak 发}------{\cjkfont 发\allowbreak 现} v1 {\cjkfont 召\allowbreak 回\allowbreak 问\allowbreak 题\allowbreak 、\allowbreak 分\allowbreak 析\allowbreak 错\allowbreak 误\allowbreak 类\allowbreak 别\allowbreak 、\allowbreak 设\allowbreak 计} v2{\cjkfont 、\allowbreak 运\allowbreak 行\allowbreak 消\allowbreak 融\allowbreak 、\allowbreak 选\allowbreak 定\allowbreak 生\allowbreak 产\allowbreak 配\allowbreak 置\allowbreak 。\allowbreak 开\allowbreak 发\allowbreak 与\allowbreak 最\allowbreak 终\allowbreak 报\allowbreak 告\allowbreak 共\allowbreak 用\allowbreak 同\allowbreak 一\allowbreak 样\allowbreak 本\allowbreak ，\allowbreak 因\allowbreak 此\allowbreak 策\allowbreak 略\allowbreak 结\allowbreak 果\allowbreak 带\allowbreak 有\allowbreak 开\allowbreak 发\allowbreak 集\allowbreak 复\allowbreak 用\allowbreak 风\allowbreak 险\allowbreak ；\allowbreak 三\allowbreak 个\allowbreak 少\allowbreak 样\allowbreak 本\allowbreak 示\allowbreak 例\allowbreak 取\allowbreak 自\allowbreak 同\allowbreak 一\allowbreak 银\allowbreak 标\allowbreak 准\allowbreak ，\allowbreak 未\allowbreak 按\allowbreak 显\allowbreak 式\allowbreak 流\allowbreak 程\allowbreak 与\allowbreak 测\allowbreak 试\allowbreak 记\allowbreak 录\allowbreak 隔\allowbreak 离\allowbreak 。\allowbreak 所\allowbreak 报\allowbreak 数\allowbreak 字\allowbreak 最\allowbreak 好\allowbreak 解\allowbreak 读\allowbreak 为\allowbreak 开\allowbreak 发\allowbreak 评\allowbreak 估\allowbreak ，\allowbreak 而\allowbreak 非\allowbreak 冻\allowbreak 结\allowbreak 测\allowbreak 试\allowbreak 集\allowbreak 认\allowbreak 证\allowbreak 。} 3. \textit{{\cjkfont 一\allowbreak 致\allowbreak 性\allowbreak 样\allowbreak 本}}{\cjkfont （}n=289{\cjkfont ）\allowbreak 。\allowbreak 两\allowbreak 阶\allowbreak 段\allowbreak 分\allowbreak 层\allowbreak 抽\allowbreak 样\allowbreak （\allowbreak 事\allowbreak 故\allowbreak 类\allowbreak 型} $\times$ {\cjkfont 等\allowbreak 级\allowbreak ，\allowbreak 比\allowbreak 例\allowbreak 分\allowbreak 配\allowbreak 、\allowbreak 每\allowbreak 格\allowbreak 保\allowbreak 底} 2 {\cjkfont 条\allowbreak ）\allowbreak 从} 2,623 {\cjkfont 条\allowbreak 全\allowbreak 量\allowbreak 语\allowbreak 料\allowbreak 抽\allowbreak 取}------{\cjkfont 而\allowbreak 非\allowbreak 银\allowbreak 标\allowbreak 准\allowbreak 。\allowbreak 一\allowbreak 名\allowbreak 受\allowbreak 训\allowbreak 标\allowbreak 注\allowbreak 者\allowbreak （\allowbreak 安\allowbreak 全\allowbreak 工\allowbreak 程\allowbreak 研\allowbreak 究\allowbreak 生\allowbreak ）\allowbreak 经} LabelStudio {\cjkfont 标\allowbreak 注\allowbreak 六\allowbreak 项\allowbreak 基\allowbreak 础\allowbreak 特\allowbreak 征\allowbreak （\allowbreak 事\allowbreak 故\allowbreak 类\allowbreak 型\allowbreak 、\allowbreak 等\allowbreak 级\allowbreak 、\allowbreak 省\allowbreak 份\allowbreak 、\allowbreak 死\allowbreak 亡\allowbreak 人\allowbreak 数\allowbreak 、\allowbreak 受\allowbreak 伤\allowbreak 人\allowbreak 数\allowbreak 、\allowbreak 日\allowbreak 期\allowbreak ）\allowbreak ，\allowbreak 标\allowbreak 注\allowbreak 过\allowbreak 程\allowbreak 未\allowbreak 接\allowbreak 触\allowbreak 管\allowbreak 道\allowbreak 的\allowbreak 抽\allowbreak 取\allowbreak 值\allowbreak 。\allowbreak 原} 300 {\cjkfont 条\allowbreak 中} 11 {\cjkfont 条\allowbreak 因\allowbreak 多\allowbreak 事\allowbreak 故\allowbreak 文\allowbreak 本\allowbreak 被\allowbreak 排\allowbreak 除\allowbreak 。}

{\cjkfont 图} 4 {\cjkfont 汇\allowbreak 总\allowbreak 样\allowbreak 本\allowbreak 层\allowbreak 级\allowbreak 。}

\begin{figure}[tbp]
\centering
\includegraphics[width=0.92\textwidth,height=0.80\textheight,keepaspectratio]{/home/suntianhao/construction_agent/.claude/worktrees/classifier-mod/docs/paper_v2/figures/fig4_sampling.png}
\end{figure}
\textbf{{\cjkfont 图} 4. {\cjkfont 数\allowbreak 据\allowbreak 集\allowbreak 层\allowbreak 级\allowbreak 与\allowbreak 三\allowbreak 个\allowbreak 评\allowbreak 测\allowbreak 样\allowbreak 本}}
\subsection{{\cjkfont 评\allowbreak 测\allowbreak 指\allowbreak 标}}
{\cjkfont 基\allowbreak 础\allowbreak 字\allowbreak 段\allowbreak 抽\allowbreak 取\allowbreak 采\allowbreak 用\allowbreak 逐\allowbreak 字\allowbreak 段\allowbreak 精\allowbreak 确\allowbreak 匹\allowbreak 配\allowbreak 准\allowbreak 确\allowbreak 率\allowbreak ，\allowbreak 地\allowbreak 点\allowbreak 字\allowbreak 段\allowbreak 做\allowbreak 行\allowbreak 政\allowbreak 后\allowbreak 缀\allowbreak 归\allowbreak 一\allowbreak 化\allowbreak 。\allowbreak 致\allowbreak 因\allowbreak 编\allowbreak 码\allowbreak 按\allowbreak 维\allowbreak 度\allowbreak （\allowbreak 管\allowbreak 理\allowbreak 、\allowbreak 防\allowbreak 线\allowbreak 、\allowbreak 触\allowbreak 发\allowbreak 、\allowbreak 人\allowbreak 的\allowbreak 行\allowbreak 为\allowbreak ）\allowbreak 计\allowbreak 算\allowbreak 精\allowbreak 确\allowbreak 率\allowbreak 、\allowbreak 召\allowbreak 回\allowbreak 率\allowbreak 与} F1{\cjkfont ，\allowbreak 采\allowbreak 用} Liu et al. (2025a) {\cjkfont 的\allowbreak 维\allowbreak 度\allowbreak 级\allowbreak 匹\allowbreak 配\allowbreak 定\allowbreak 义\allowbreak ；\allowbreak 六\allowbreak 策\allowbreak 略\allowbreak 对\allowbreak 比\allowbreak 另\allowbreak 报\allowbreak 百\allowbreak 分\allowbreak 位} bootstrap 95\% {\cjkfont 置\allowbreak 信\allowbreak 区\allowbreak 间\allowbreak （}1,000 {\cjkfont 次\allowbreak 重\allowbreak 采\allowbreak 样\allowbreak ）\allowbreak 。\allowbreak 标\allowbreak 注\allowbreak 者\allowbreak 一\allowbreak 致\allowbreak 性\allowbreak 方\allowbreak 面\allowbreak ：\allowbreak 分\allowbreak 类\allowbreak 字\allowbreak 段\allowbreak 用\allowbreak 带} 95\% {\cjkfont 置\allowbreak 信\allowbreak 区\allowbreak 间\allowbreak 的} Cohen's kappa{\cjkfont ，\allowbreak 计\allowbreak 数\allowbreak 字\allowbreak 段\allowbreak 用} ICC(2,1){\cjkfont ，\allowbreak 日\allowbreak 期\allowbreak 用\allowbreak 年}-{\cjkfont 月\allowbreak 匹\allowbreak 配\allowbreak 率}------{\cjkfont 三\allowbreak 种\allowbreak 统\allowbreak 计\allowbreak 量\allowbreak 按\allowbreak 各\allowbreak 自\allowbreak 名\allowbreak 称\allowbreak 报\allowbreak 告\allowbreak ，\allowbreak 不\allowbreak 以\allowbreak 单\allowbreak 一\allowbreak 标\allowbreak 签\allowbreak 统\allowbreak 称\allowbreak 。}

\subsection{{\cjkfont 实\allowbreak 验\allowbreak 一\allowbreak ：\allowbreak 显\allowbreak 式\allowbreak 约\allowbreak 束\allowbreak 下\allowbreak 的\allowbreak 骨\allowbreak 干\allowbreak 选\allowbreak 型}}
{\cjkfont 四\allowbreak 个\allowbreak 候\allowbreak 选\allowbreak 模\allowbreak 型\allowbreak 在\allowbreak 模\allowbreak 型\allowbreak 比\allowbreak 较\allowbreak 样\allowbreak 本\allowbreak 上\allowbreak 评\allowbreak 测\allowbreak ：}DeepSeek V4{\cjkfont 、}Qwen 3.7{\cjkfont 、}GLM 5.2{\cjkfont 、}Kimi K2.6{\cjkfont ，\allowbreak 使\allowbreak 用\allowbreak 相\allowbreak 同\allowbreak 系\allowbreak 统\allowbreak 提\allowbreak 示\allowbreak 与\allowbreak 三\allowbreak 个\allowbreak 少\allowbreak 样\allowbreak 本\allowbreak 示\allowbreak 例\allowbreak 。\allowbreak 各\allowbreak 模\allowbreak 型\allowbreak 运\allowbreak 行\allowbreak 条\allowbreak 件\allowbreak 并\allowbreak 非\allowbreak 完\allowbreak 全\allowbreak 对\allowbreak 等\allowbreak ：}Kimi K2.6 {\cjkfont 按} API {\cjkfont 规\allowbreak 范\allowbreak 强\allowbreak 制} temperature 1.0{\cjkfont ；}GLM 5.2 {\cjkfont 因\allowbreak 推\allowbreak 理\allowbreak 痕\allowbreak 迹\allowbreak 需} 16,384 token {\cjkfont 输\allowbreak 出\allowbreak 预\allowbreak 算\allowbreak ；}DeepSeek V4 {\cjkfont 经\allowbreak 提\allowbreak 供\allowbreak 商\allowbreak 自\allowbreak 动\allowbreak 档\allowbreak 位\allowbreak 路\allowbreak 由\allowbreak 访\allowbreak 问\allowbreak ，\allowbreak 运\allowbreak 行\allowbreak 日\allowbreak 志\allowbreak 中\allowbreak 的\allowbreak 请\allowbreak 求\allowbreak 标\allowbreak 识\allowbreak 覆\allowbreak 盖\allowbreak 其\allowbreak 档\allowbreak 位\allowbreak 家\allowbreak 族\allowbreak 。\allowbreak 由\allowbreak 于\allowbreak 这\allowbreak 些\allowbreak 差\allowbreak 异\allowbreak ，\allowbreak 且\allowbreak 每\allowbreak 种\allowbreak 配\allowbreak 置\allowbreak 仅\allowbreak 运\allowbreak 行\allowbreak 一\allowbreak 次\allowbreak ，\allowbreak 本\allowbreak 实\allowbreak 验\allowbreak 定\allowbreak 位\allowbreak 为\allowbreak 支\allowbreak 持\allowbreak 选\allowbreak 型\allowbreak 决\allowbreak 策\allowbreak 的\allowbreak 约\allowbreak 束\allowbreak 性\allowbreak 系\allowbreak 统\allowbreak 比\allowbreak 较\allowbreak ，\allowbreak 而\allowbreak 非\allowbreak 一\allowbreak 般\allowbreak 性\allowbreak 模\allowbreak 型\allowbreak 能\allowbreak 力\allowbreak 排\allowbreak 名\allowbreak 。}

{\cjkfont 表} 5 {\cjkfont 报\allowbreak 告\allowbreak 逐\allowbreak 字\allowbreak 段\allowbreak 准\allowbreak 确\allowbreak 率\allowbreak 。\allowbreak 三\allowbreak 个\allowbreak 最\allowbreak 强\allowbreak 模\allowbreak 型\allowbreak 总\allowbreak 体\allowbreak 达} 91.2--91.5\%{\cjkfont ，\allowbreak 事\allowbreak 故\allowbreak 等\allowbreak 级\allowbreak 、\allowbreak 死\allowbreak 亡\allowbreak 人\allowbreak 数\allowbreak 、\allowbreak 省\allowbreak 级\allowbreak 地\allowbreak 点\allowbreak 普\allowbreak 遍\allowbreak 高\allowbreak 准\allowbreak 确\allowbreak 。\allowbreak 城\allowbreak 市\allowbreak 提\allowbreak 取\allowbreak 最\allowbreak 难\allowbreak （}81.5--88.0\%{\cjkfont ）\allowbreak ：\allowbreak 中\allowbreak 文\allowbreak 报\allowbreak 道\allowbreak 常\allowbreak 点\allowbreak 名\allowbreak 县\allowbreak 级\allowbreak 市\allowbreak 而\allowbreak 参\allowbreak 考\allowbreak 记\allowbreak 录\allowbreak 为\allowbreak 上\allowbreak 级\allowbreak 地\allowbreak 级\allowbreak 市\allowbreak 。}Kimi K2.6 {\cjkfont 大\allowbreak 幅\allowbreak 落\allowbreak 后\allowbreak （}63.9\%{\cjkfont ）\allowbreak ，\allowbreak 与\allowbreak 其\allowbreak 强\allowbreak 制\allowbreak 温\allowbreak 度\allowbreak 带\allowbreak 来\allowbreak 的\allowbreak 输\allowbreak 出\allowbreak 不\allowbreak 稳\allowbreak 定\allowbreak 一\allowbreak 致\allowbreak 。}

\textbf{{\cjkfont 表} 5. {\cjkfont 基\allowbreak 础\allowbreak 字\allowbreak 段\allowbreak 提\allowbreak 取\allowbreak 准\allowbreak 确\allowbreak 率\allowbreak （}\%{\cjkfont ）\allowbreak ，\allowbreak 实\allowbreak 验\allowbreak 一}}

\begingroup\setlength{\tabcolsep}{2pt}\small
\begin{longtable}{>{\raggedright\arraybackslash}p{0.1294\textwidth}>{\raggedright\arraybackslash}p{0.2162\textwidth}>{\raggedright\arraybackslash}p{0.2017\textwidth}>{\raggedright\arraybackslash}p{0.2009\textwidth}>{\raggedright\arraybackslash}p{0.2026\textwidth}}
\hline
{\cjkfont 字\allowbreak 段} & DeepSeek V4 & Qwen 3.7 & GLM 5.2 & Kimi K2.6 \\
\hline\endfirsthead
\hline
{\cjkfont 字\allowbreak 段} & DeepSeek V4 & Qwen 3.7 & GLM 5.2 & Kimi K2.6 \\
\hline\endhead
{\cjkfont 事\allowbreak 故\allowbreak 类\allowbreak 型} & 86.0 & 85.0 & 84.5 & 61.5 \\
{\cjkfont 事\allowbreak 故\allowbreak 等\allowbreak 级} & 98.0 & 90.5 & 97.5 & 69.5 \\
{\cjkfont 死\allowbreak 亡\allowbreak 人\allowbreak 数} & 96.0 & 96.0 & 97.0 & 68.5 \\
{\cjkfont 受\allowbreak 伤\allowbreak 人\allowbreak 数} & 90.0 & 91.0 & 91.0 & 58.5 \\
{\cjkfont 省\allowbreak 份} & 96.0 & 98.5 & 96.0 & 67.0 \\
{\cjkfont 城\allowbreak 市} & 81.5 & 88.0 & 83.0 & 58.5 \\
\textbf{{\cjkfont 总\allowbreak 体}} & \textbf{91.2} & \textbf{91.5} & \textbf{91.5} & \textbf{63.9} \\
\hline
\end{longtable}\endgroup
{\cjkfont 表} 6 {\cjkfont 报\allowbreak 告\allowbreak 按\allowbreak 维\allowbreak 度\allowbreak 的\allowbreak 致\allowbreak 因\allowbreak 编\allowbreak 码\allowbreak 表\allowbreak 现\allowbreak 。}DeepSeek V4 {\cjkfont 平\allowbreak 均\allowbreak 精\allowbreak 确\allowbreak 率\allowbreak 最\allowbreak 高\allowbreak （}89.5\%{\cjkfont ）\allowbreak 、\allowbreak 强\allowbreak 模\allowbreak 型\allowbreak 间\allowbreak 召\allowbreak 回\allowbreak 率\allowbreak 最\allowbreak 低\allowbreak （}52.7\%{\cjkfont ）\allowbreak ；}Qwen 3.7 {\cjkfont 呈\allowbreak 镜\allowbreak 像\allowbreak 取\allowbreak 舍\allowbreak （\allowbreak 精\allowbreak 确} 78.0\%{\cjkfont 、\allowbreak 召\allowbreak 回} 80.5\%{\cjkfont ）\allowbreak 。\allowbreak 成\allowbreak 本\allowbreak 与\allowbreak 延\allowbreak 迟\allowbreak 补\allowbreak 全\allowbreak 选\allowbreak 型\allowbreak 图\allowbreak 景\allowbreak ：}200 {\cjkfont 条\allowbreak 记\allowbreak 录\allowbreak 成\allowbreak 本} 0.13 {\cjkfont 美\allowbreak 元\allowbreak （}DeepSeek V4{\cjkfont ）\allowbreak 、}0.43{\cjkfont （}Qwen 3.7{\cjkfont ）\allowbreak 、}2.91{\cjkfont （}GLM 5.2{\cjkfont ）\allowbreak 、}2.84{\cjkfont （}Kimi K2.6{\cjkfont ）\allowbreak ；\allowbreak 单\allowbreak 条\allowbreak 延\allowbreak 迟} 9.8 {\cjkfont 秒\allowbreak 、}8.6 {\cjkfont 秒\allowbreak 、}27.0 {\cjkfont 秒\allowbreak 、}74.9 {\cjkfont 秒\allowbreak 。\allowbreak 管\allowbreak 道\allowbreak 采\allowbreak 用} DeepSeek V4{\cjkfont 。}

\textbf{{\cjkfont 表} 6. {\cjkfont 致\allowbreak 因\allowbreak 编\allowbreak 码\allowbreak 精\allowbreak 确\allowbreak 率} / {\cjkfont 召\allowbreak 回\allowbreak 率} / F1{\cjkfont （}\%{\cjkfont ）\allowbreak ，\allowbreak 按\allowbreak 维\allowbreak 度\allowbreak ，\allowbreak 实\allowbreak 验\allowbreak 一}}

\begingroup\setlength{\tabcolsep}{2pt}\small
\begin{longtable}{>{\raggedright\arraybackslash}p{0.0784\textwidth}>{\raggedright\arraybackslash}p{0.2308\textwidth}>{\raggedright\arraybackslash}p{0.2139\textwidth}>{\raggedright\arraybackslash}p{0.2134\textwidth}>{\raggedright\arraybackslash}p{0.2143\textwidth}}
\hline
{\cjkfont 维\allowbreak 度} & DeepSeek V4 & Qwen 3.7 & GLM 5.2 & Kimi K2.6 \\
\hline\endfirsthead
\hline
{\cjkfont 维\allowbreak 度} & DeepSeek V4 & Qwen 3.7 & GLM 5.2 & Kimi K2.6 \\
\hline\endhead
{\cjkfont 管\allowbreak 理} & 89.9 / 54.9 / 68.1 & 67.5 / 95.6 / 79.1 & 92.3 / 63.7 / 75.4 & 97.6 / 36.3 / 52.9 \\
{\cjkfont 防\allowbreak 线} & 91.9 / 41.9 / 57.6 & 88.8 / 64.0 / 74.4 & 87.6 / 67.6 / 76.3 & 96.4 / 39.0 / 55.5 \\
{\cjkfont 触\allowbreak 发} & 81.0 / 69.4 / 74.7 & 72.4 / 85.7 / 78.5 & 73.9 / 69.4 / 71.6 & 81.5 / 44.9 / 57.9 \\
{\cjkfont 人\allowbreak 的\allowbreak 行\allowbreak 为} & 95.0 / 44.7 / 60.8 & 83.3 / 76.5 / 79.8 & 83.9 / 61.2 / 70.7 & 88.2 / 35.3 / 50.4 \\
\textbf{{\cjkfont 平\allowbreak 均}} & \textbf{89.5 / 52.7 / 65.3} & \textbf{78.0 / 80.5 / 78.0} & \textbf{84.4 / 65.5 / 73.5} & \textbf{90.9 / 38.9 / 54.2} \\
\hline
\end{longtable}\endgroup
\subsection{{\cjkfont 实\allowbreak 验\allowbreak 二\allowbreak ：\allowbreak 提\allowbreak 示\allowbreak 策\allowbreak 略\allowbreak 对\allowbreak 比\allowbreak 与\allowbreak 词\allowbreak 汇\allowbreak 门\allowbreak 控\allowbreak 观\allowbreak 察}}
\subsubsection{{\cjkfont 实\allowbreak 验\allowbreak 设\allowbreak 计}}
{\cjkfont 六\allowbreak 种\allowbreak 提\allowbreak 示\allowbreak 策\allowbreak 略\allowbreak 在\allowbreak 策\allowbreak 略\allowbreak 样\allowbreak 本\allowbreak （}n=100{\cjkfont ，\allowbreak 三\allowbreak 专\allowbreak 家\allowbreak 多\allowbreak 数\allowbreak 投\allowbreak 票\allowbreak 真\allowbreak 值\allowbreak ）\allowbreak 上\allowbreak 以} DeepSeek V4{\cjkfont （}temperature 0.0{\cjkfont ）\allowbreak 对\allowbreak 比\allowbreak 评\allowbreak 测\allowbreak ：}S0 {\cjkfont 基\allowbreak 础\allowbreak 零\allowbreak 样\allowbreak 本\allowbreak 、}S1 {\cjkfont 基\allowbreak 础\allowbreak 少\allowbreak 样\allowbreak 本\allowbreak （\allowbreak 三\allowbreak 个\allowbreak 专\allowbreak 家\allowbreak 标\allowbreak 注\allowbreak 示\allowbreak 例\allowbreak ）\allowbreak 、}S2 Auto-CoT{\cjkfont 、}S3 {\cjkfont 计\allowbreak 划\allowbreak 与\allowbreak 求\allowbreak 解\allowbreak 、}S4 CAICS-CoT{\cjkfont （\allowbreak 两\allowbreak 步\allowbreak 判\allowbreak 断\allowbreak 领\allowbreak 域\allowbreak 知\allowbreak 识} CoT{\cjkfont ）\allowbreak 、}S5 CAICS-CoT+{\cjkfont （\allowbreak 带\allowbreak 层\allowbreak 间\allowbreak 摘\allowbreak 要\allowbreak 注\allowbreak 入\allowbreak 的\allowbreak 串\allowbreak 行\allowbreak 逐\allowbreak 层\allowbreak 调\allowbreak 用\allowbreak ）\allowbreak 。\allowbreak 附\allowbreak 录} A {\cjkfont 以\allowbreak 压\allowbreak 缩\allowbreak 形\allowbreak 式\allowbreak 收\allowbreak 录\allowbreak 各\allowbreak 提\allowbreak 示\allowbreak 。}

\textbf{{\cjkfont 表} 7. {\cjkfont 六\allowbreak 种\allowbreak 提\allowbreak 示\allowbreak 策\allowbreak 略\allowbreak 性\allowbreak 能\allowbreak （\allowbreak 多\allowbreak 数\allowbreak 投\allowbreak 票\allowbreak 真\allowbreak 值\allowbreak ，}n=100{\cjkfont ）}}

\begingroup\setlength{\tabcolsep}{2pt}\small
\begin{longtable}{>{\raggedright\arraybackslash}p{0.1605\textwidth}>{\raggedright\arraybackslash}p{0.1800\textwidth}>{\raggedright\arraybackslash}p{0.1522\textwidth}>{\raggedright\arraybackslash}p{0.0871\textwidth}>{\raggedright\arraybackslash}p{0.0901\textwidth}>{\raggedright\arraybackslash}p{0.0901\textwidth}>{\raggedright\arraybackslash}p{0.0761\textwidth}>{\raggedright\arraybackslash}p{0.0856\textwidth}}
\hline
{\cjkfont 策\allowbreak 略} & {\cjkfont 平\allowbreak 均} F1 [95\% CI] & {\cjkfont 平\allowbreak 均\allowbreak 召\allowbreak 回} & {\cjkfont 平\allowbreak 均\allowbreak 精\allowbreak 确} & {\cjkfont 类\allowbreak 型\allowbreak 准\allowbreak 确\allowbreak 率} & {\cjkfont 等\allowbreak 级\allowbreak 准\allowbreak 确\allowbreak 率} & {\cjkfont 有\allowbreak 效\allowbreak 输\allowbreak 出\allowbreak 率} & {\cjkfont 耗\allowbreak 时}(s) \\
\hline\endfirsthead
\hline
{\cjkfont 策\allowbreak 略} & {\cjkfont 平\allowbreak 均} F1 [95\% CI] & {\cjkfont 平\allowbreak 均\allowbreak 召\allowbreak 回} & {\cjkfont 平\allowbreak 均\allowbreak 精\allowbreak 确} & {\cjkfont 类\allowbreak 型\allowbreak 准\allowbreak 确\allowbreak 率} & {\cjkfont 等\allowbreak 级\allowbreak 准\allowbreak 确\allowbreak 率} & {\cjkfont 有\allowbreak 效\allowbreak 输\allowbreak 出\allowbreak 率} & {\cjkfont 耗\allowbreak 时}(s) \\
\hline\endhead
S0 {\cjkfont 零\allowbreak 样\allowbreak 本} & 0.600 [0.520, 0.644] & 0.515 & 0.808 & 0.630 & 0.960 & 96\% & 1,382 \\
S1 {\cjkfont 少\allowbreak 样\allowbreak 本} & \textbf{0.669 [0.615, 0.722]} & 0.634 & 0.753 & 0.840 & 0.970 & 100\% & 1,302 \\
S2 Auto-CoT & 0.634 [0.564, 0.691] & 0.564 & 0.772 & 0.860 & 0.980 & 85\% & 2,401 \\
S3 {\cjkfont 计\allowbreak 划\allowbreak 与\allowbreak 求\allowbreak 解} & 0.598 [0.529, 0.652] & 0.540 & 0.741 & 0.850 & 0.990 & 92\% & 1,735 \\
S4 CAICS-CoT & 0.656 [0.603, 0.695] & \textbf{0.872} & 0.565 & 0.830 & 0.990 & 99\% & 1,843 \\
S5 CAICS-CoT+ & 0.665 [0.606, 0.719] & 0.750 & 0.618 & 0.830 & 0.980 & 100\% & 5,430 \\
\hline
\end{longtable}\endgroup
\textit{\textit{{\cjkfont 注\allowbreak 。}} {\cjkfont 真\allowbreak 值\allowbreak 为\allowbreak 三\allowbreak 位\allowbreak 独\allowbreak 立\allowbreak 专\allowbreak 家\allowbreak 标\allowbreak 注\allowbreak 者\allowbreak 的\allowbreak 多\allowbreak 数\allowbreak 投\allowbreak 票\allowbreak 。\allowbreak 置\allowbreak 信\allowbreak 区\allowbreak 间\allowbreak 为\allowbreak 百\allowbreak 分\allowbreak 位} bootstrap{\cjkfont （}1,000 {\cjkfont 次\allowbreak 重\allowbreak 采\allowbreak 样\allowbreak ）\allowbreak 。}S2 {\cjkfont 的\allowbreak 编\allowbreak 码\allowbreak 指\allowbreak 标\allowbreak 与\allowbreak 类\allowbreak 型}/{\cjkfont 等\allowbreak 级\allowbreak 准\allowbreak 确\allowbreak 率\allowbreak 均\allowbreak 经\allowbreak 后\allowbreak 处\allowbreak 理\allowbreak 恢\allowbreak 复\allowbreak 计\allowbreak 算\allowbreak （}96/100 {\cjkfont 输\allowbreak 出\allowbreak 发\allowbreak 生\allowbreak 层\allowbreak 级\allowbreak 格\allowbreak 式\allowbreak 漂\allowbreak 移\allowbreak ）\allowbreak ；\allowbreak 其\allowbreak 他\allowbreak 策\allowbreak 略\allowbreak 未\allowbreak 做\allowbreak 后\allowbreak 处\allowbreak 理\allowbreak ，\allowbreak 因\allowbreak 此} S2 {\cjkfont 的\allowbreak 得\allowbreak 分\allowbreak 描\allowbreak 述\allowbreak 的\allowbreak 是\allowbreak 后\allowbreak 处\allowbreak 理\allowbreak 管\allowbreak 道\allowbreak 而\allowbreak 非\allowbreak 原\allowbreak 始\allowbreak 输\allowbreak 出\allowbreak 。\allowbreak 全\allowbreak 部\allowbreak 配\allowbreak 置\allowbreak 在} temperature 0.0 {\cjkfont 下\allowbreak 各\allowbreak 运\allowbreak 行\allowbreak 一\allowbreak 次\allowbreak 。}S1{\cjkfont 、}S4{\cjkfont 、}S5 {\cjkfont 的} F1 {\cjkfont 差\allowbreak 异\allowbreak 处\allowbreak 于\allowbreak 重\allowbreak 叠\allowbreak 的\allowbreak 置\allowbreak 信\allowbreak 区\allowbreak 间\allowbreak 内\allowbreak ，\allowbreak 故\allowbreak 将} S1 {\cjkfont 描\allowbreak 述\allowbreak 为\allowbreak 最\allowbreak 高\allowbreak 点\allowbreak 估\allowbreak 计\allowbreak ，\allowbreak 而\allowbreak 非\allowbreak 经\allowbreak 统\allowbreak 计\allowbreak 确\allowbreak 认\allowbreak 的\allowbreak 胜\allowbreak 者\allowbreak 。}}\\
\subsubsection{{\cjkfont 少\allowbreak 样\allowbreak 本\allowbreak 在\allowbreak 本\allowbreak 模\allowbreak 型}-{\cjkfont 任\allowbreak 务\allowbreak 组\allowbreak 合\allowbreak 上\allowbreak 优\allowbreak 于\allowbreak 领\allowbreak 域} CoT}
{\cjkfont 表} 7 {\cjkfont 汇\allowbreak 总\allowbreak 了\allowbreak 该\allowbreak 对\allowbreak 比\allowbreak 。}S1 {\cjkfont 的\allowbreak 示\allowbreak 例\allowbreak 实\allowbreak 现\allowbreak 格\allowbreak 式}-{\cjkfont 内\allowbreak 容\allowbreak 联\allowbreak 合\allowbreak 校\allowbreak 准\allowbreak ：\allowbreak 示\allowbreak 范\allowbreak 输\allowbreak 出} schema{\cjkfont 、\allowbreak 期\allowbreak 望\allowbreak 编\allowbreak 码\allowbreak 密\allowbreak 度\allowbreak 、\allowbreak 编\allowbreak 码\allowbreak 组\allowbreak 合\allowbreak 模\allowbreak 式\allowbreak 与\allowbreak 推\allowbreak 理\allowbreak 深\allowbreak 度\allowbreak 。}S1 {\cjkfont 的\allowbreak 平\allowbreak 均\allowbreak 编\allowbreak 码\allowbreak 密\allowbreak 度\allowbreak 在\allowbreak 六\allowbreak 策\allowbreak 略\allowbreak 中\allowbreak 最\allowbreak 接\allowbreak 近\allowbreak 真\allowbreak 值\allowbreak （}4.2 {\cjkfont 对} 5.1{\cjkfont ）\allowbreak ，}S0 {\cjkfont 偏\allowbreak 少\allowbreak （}2.8{\cjkfont ）\allowbreak 、}S4 {\cjkfont 偏\allowbreak 多\allowbreak （}9.2{\cjkfont ）\allowbreak 。\allowbreak 与} Liu et al. (2025a) {\cjkfont 的\allowbreak 比\allowbreak 较\allowbreak 有\allowbreak 信\allowbreak 息\allowbreak 量\allowbreak 但\allowbreak 存\allowbreak 在\allowbreak 混\allowbreak 淆\allowbreak ：\allowbreak 两\allowbreak 项\allowbreak 研\allowbreak 究\allowbreak 在\allowbreak 模\allowbreak 型\allowbreak （}GPT-4o {\cjkfont 对} DeepSeek V4{\cjkfont ）\allowbreak 、\allowbreak 任\allowbreak 务\allowbreak （\allowbreak 航\allowbreak 空} HFACS {\cjkfont 对\allowbreak 施\allowbreak 工} CAICS{\cjkfont ）\allowbreak 、\allowbreak 语\allowbreak 料\allowbreak 与\allowbreak 评\allowbreak 测\allowbreak 条\allowbreak 件\allowbreak 上\allowbreak 同\allowbreak 时\allowbreak 不\allowbreak 同\allowbreak 。\allowbreak 因\allowbreak 此\allowbreak 观\allowbreak 测\allowbreak 到\allowbreak 的\allowbreak 排\allowbreak 名\allowbreak 差\allowbreak 异\allowbreak 与\allowbreak 模\allowbreak 型\allowbreak 依\allowbreak 赖\allowbreak 、\allowbreak 任\allowbreak 务\allowbreak 依\allowbreak 赖\allowbreak 或\allowbreak 两\allowbreak 者\allowbreak 兼\allowbreak 有\allowbreak 一\allowbreak 致\allowbreak ；\allowbreak 我\allowbreak 们\allowbreak 不\allowbreak 声\allowbreak 称\allowbreak 分\allowbreak 离\allowbreak 出\allowbreak 了\allowbreak 模\allowbreak 型\allowbreak 依\allowbreak 赖\allowbreak ，\allowbreak 实\allowbreak 践\allowbreak 解\allowbreak 读\allowbreak 仅\allowbreak 为\allowbreak ：\allowbreak 新\allowbreak 模\allowbreak 型}-{\cjkfont 任\allowbreak 务\allowbreak 组\allowbreak 合\allowbreak 的\allowbreak 最\allowbreak 优\allowbreak 策\allowbreak 略\allowbreak 不\allowbreak 能\allowbreak 从\allowbreak 单\allowbreak 一\allowbreak 已\allowbreak 发\allowbreak 表\allowbreak 排\allowbreak 名\allowbreak 推\allowbreak 定\allowbreak 。}

\subsubsection{{\cjkfont 词\allowbreak 汇\allowbreak 门\allowbreak 控\allowbreak 观\allowbreak 察\allowbreak 及\allowbreak 其\allowbreak 补\allowbreak 救}}
S4 {\cjkfont 初\allowbreak 版\allowbreak 提\allowbreak 示\allowbreak （}v1{\cjkfont ）\allowbreak 将\allowbreak 领\allowbreak 域\allowbreak 知\allowbreak 识\allowbreak 嵌\allowbreak 入\allowbreak 为\allowbreak 条\allowbreak 件\allowbreak 检\allowbreak 查\allowbreak ：\allowbreak 每\allowbreak 个\allowbreak 编\allowbreak 码\allowbreak 配\allowbreak 一\allowbreak 个\allowbreak 疑\allowbreak 问\allowbreak 式\allowbreak 门\allowbreak 控\allowbreak （}"{\cjkfont 文\allowbreak 本\allowbreak 提\allowbreak 到}'{\cjkfont 无\allowbreak 人\allowbreak 监\allowbreak 管}''{\cjkfont 监\allowbreak 理\allowbreak 失\allowbreak 职}'{\cjkfont 吗\allowbreak ？}"{\cjkfont ）\allowbreak 加\allowbreak 保\allowbreak 守\allowbreak 指\allowbreak 令\allowbreak 。\allowbreak 在\allowbreak 策\allowbreak 略\allowbreak 样\allowbreak 本\allowbreak 上\allowbreak ，}v1 {\cjkfont 的} F1=0.503{\cjkfont 、\allowbreak 召\allowbreak 回}=0.423{\cjkfont 、}21\% {\cjkfont 空\allowbreak 输\allowbreak 出\allowbreak 、}M1 {\cjkfont 召\allowbreak 回} 0.141------{\cjkfont 而} M1 {\cjkfont 是\allowbreak 真\allowbreak 值\allowbreak 中\allowbreak 最\allowbreak 高\allowbreak 频\allowbreak 编\allowbreak 码\allowbreak （}100 {\cjkfont 条\allowbreak 中\allowbreak 出\allowbreak 现} 85 {\cjkfont 次\allowbreak ）\allowbreak 。\allowbreak 多\allowbreak 数\allowbreak 事\allowbreak 故\allowbreak 叙\allowbreak 述\allowbreak 从\allowbreak 不\allowbreak 使\allowbreak 用\allowbreak 这\allowbreak 些\allowbreak 精\allowbreak 确\allowbreak 措\allowbreak 辞\allowbreak ，\allowbreak 而\allowbreak 是\allowbreak 写}"{\cjkfont 安\allowbreak 全\allowbreak 管\allowbreak 理\allowbreak 不\allowbreak 到\allowbreak 位}""{\cjkfont 未\allowbreak 及\allowbreak 时\allowbreak 发\allowbreak 现\allowbreak 制\allowbreak 止}"{\cjkfont 。\allowbreak 在\allowbreak 条\allowbreak 件\allowbreak 检\allowbreak 查\allowbreak 措\allowbreak 辞\allowbreak 下\allowbreak ，\allowbreak 模\allowbreak 型\allowbreak 似\allowbreak 乎\allowbreak 将\allowbreak 所\allowbreak 列\allowbreak 关\allowbreak 键\allowbreak 词\allowbreak 当\allowbreak 作\allowbreak 强\allowbreak 制\allowbreak 性\allowbreak 门\allowbreak 槛\allowbreak 而\allowbreak 非\allowbreak 示\allowbreak 例\allowbreak 线\allowbreak 索\allowbreak 。}

{\cjkfont 我\allowbreak 们\allowbreak 将\allowbreak 这\allowbreak 一\allowbreak 观\allowbreak 测\allowbreak 模\allowbreak 式\allowbreak 命\allowbreak 名\allowbreak 为}\textbf{{\cjkfont 词\allowbreak 汇\allowbreak 门\allowbreak 控\allowbreak 效\allowbreak 应}}{\cjkfont ：\allowbreak 疑\allowbreak 问\allowbreak 式\allowbreak 条\allowbreak 件\allowbreak 检\allowbreak 查} CoT {\cjkfont 措\allowbreak 辞\allowbreak 抑\allowbreak 制\allowbreak 隐\allowbreak 性\allowbreak 类\allowbreak 别\allowbreak 的\allowbreak 同\allowbreak 义\allowbreak 证\allowbreak 据\allowbreak ，\allowbreak 召\allowbreak 回\allowbreak 损\allowbreak 失\allowbreak 集\allowbreak 中\allowbreak 于\allowbreak 最\allowbreak 高\allowbreak 频\allowbreak 的\allowbreak 隐\allowbreak 性\allowbreak 编\allowbreak 码\allowbreak 。\allowbreak 该\allowbreak 术\allowbreak 语\allowbreak 描\allowbreak 述\allowbreak 的\allowbreak 是\allowbreak 观\allowbreak 察\allowbreak ；\allowbreak 它\allowbreak 本\allowbreak 身\allowbreak 并\allowbreak 不\allowbreak 识\allowbreak 别\allowbreak 机\allowbreak 制\allowbreak 。}

{\cjkfont 补\allowbreak 救\allowbreak 以\allowbreak 两\allowbreak 步\allowbreak 判\allowbreak 断\allowbreak 取\allowbreak 代\allowbreak 门\allowbreak 控\allowbreak ：}(1) {\cjkfont 显\allowbreak 性\allowbreak 匹\allowbreak 配}------{\cjkfont 文\allowbreak 本\allowbreak 是\allowbreak 否\allowbreak 以\allowbreak 任\allowbreak 何\allowbreak 同\allowbreak 义\allowbreak 表\allowbreak 述\allowbreak 直\allowbreak 接\allowbreak 支\allowbreak 持\allowbreak 该\allowbreak 编\allowbreak 码\allowbreak ？}(2) {\cjkfont 授\allowbreak 权\allowbreak 推\allowbreak 断}------{\cjkfont 根\allowbreak 据\allowbreak 事\allowbreak 故\allowbreak 性\allowbreak 质\allowbreak 与\allowbreak 施\allowbreak 工\allowbreak 安\allowbreak 全\allowbreak 常\allowbreak 识\allowbreak ，\allowbreak 该\allowbreak 致\allowbreak 因\allowbreak 即\allowbreak 使\allowbreak 没\allowbreak 有\allowbreak 显\allowbreak 式\allowbreak 措\allowbreak 辞\allowbreak 是\allowbreak 否\allowbreak 必\allowbreak 然\allowbreak 存\allowbreak 在\allowbreak ？\allowbreak 保\allowbreak 守\allowbreak 指\allowbreak 令\allowbreak 亦\allowbreak 替\allowbreak 换\allowbreak 为\allowbreak 全\allowbreak 覆\allowbreak 盖\allowbreak 偏\allowbreak 好\allowbreak 。}v1$\rightarrow$v2 {\cjkfont 的\allowbreak 组\allowbreak 合\allowbreak 变\allowbreak 化\allowbreak 将} F1 {\cjkfont 从} 0.503 {\cjkfont 提\allowbreak 至} 0.656{\cjkfont ，\allowbreak 召\allowbreak 回\allowbreak 从} 0.423 {\cjkfont 提\allowbreak 至} 0.872{\cjkfont ，}M1 {\cjkfont 召\allowbreak 回\allowbreak 从} 0.141 {\cjkfont 提\allowbreak 至} 0.988{\cjkfont ，\allowbreak 空\allowbreak 输\allowbreak 出\allowbreak 从} 21\% {\cjkfont 降\allowbreak 至} 0{\cjkfont 。\allowbreak 图} 5 {\cjkfont 可\allowbreak 视\allowbreak 化\allowbreak 该\allowbreak 变\allowbreak 化\allowbreak 。}

\begin{figure}[tbp]
\centering
\includegraphics[width=0.92\textwidth,height=0.80\textheight,keepaspectratio]{/home/suntianhao/construction_agent/.claude/worktrees/classifier-mod/docs/paper_v2/figures/fig5_gatekeeping.png}
\end{figure}
\textbf{{\cjkfont 图} 5. {\cjkfont 词\allowbreak 汇\allowbreak 门\allowbreak 控\allowbreak 观\allowbreak 察\allowbreak 及\allowbreak 其\allowbreak 补\allowbreak 救}}
\subsubsection{{\cjkfont 措\allowbreak 辞\allowbreak 变\allowbreak 化\allowbreak 的\allowbreak 约\allowbreak 束\allowbreak 性\allowbreak 四\allowbreak 格\allowbreak 比\allowbreak 较}}
v1$\rightarrow$v2 {\cjkfont 的\allowbreak 转\allowbreak 变\allowbreak 同\allowbreak 时\allowbreak 改\allowbreak 变\allowbreak 了\allowbreak 多\allowbreak 个\allowbreak 提\allowbreak 示\allowbreak 属\allowbreak 性}------{\cjkfont 疑\allowbreak 问\allowbreak 句\allowbreak 式\allowbreak 、\allowbreak 两\allowbreak 步\allowbreak 判\allowbreak 断\allowbreak 的\allowbreak 加\allowbreak 入\allowbreak 、\allowbreak 推\allowbreak 断\allowbreak 的\allowbreak 显\allowbreak 式\allowbreak 授\allowbreak 权\allowbreak 、\allowbreak 输\allowbreak 出\allowbreak 政\allowbreak 策\allowbreak 指\allowbreak 令\allowbreak 。\allowbreak 因\allowbreak 此\allowbreak 我\allowbreak 们\allowbreak 在\allowbreak 同\allowbreak 一\allowbreak 样\allowbreak 本\allowbreak 上\allowbreak 运\allowbreak 行\allowbreak 了\allowbreak 四\allowbreak 格\allowbreak 比\allowbreak 较\allowbreak ：}v1{\cjkfont （\allowbreak 条\allowbreak 件\allowbreak 检\allowbreak 查}+{\cjkfont 保\allowbreak 守\allowbreak ）\allowbreak 、}v2a{\cjkfont （\allowbreak 两\allowbreak 步\allowbreak 判\allowbreak 断}+{\cjkfont 保\allowbreak 守\allowbreak ）\allowbreak 、}v2b{\cjkfont （\allowbreak 条\allowbreak 件\allowbreak 检\allowbreak 查}+{\cjkfont 全\allowbreak 覆\allowbreak 盖\allowbreak ）\allowbreak 、}v2{\cjkfont （\allowbreak 两\allowbreak 步\allowbreak 判\allowbreak 断}+{\cjkfont 全\allowbreak 覆\allowbreak 盖\allowbreak ）\allowbreak 。\allowbreak 表} 8 {\cjkfont 报\allowbreak 告\allowbreak 结\allowbreak 果\allowbreak 。\allowbreak 保\allowbreak 持\allowbreak 保\allowbreak 守\allowbreak 指\allowbreak 令\allowbreak 不\allowbreak 变\allowbreak ，\allowbreak 仅\allowbreak 措\allowbreak 辞\allowbreak 改\allowbreak 变\allowbreak （}v1 $\rightarrow$ v2a{\cjkfont ）\allowbreak 将} M1 {\cjkfont 召\allowbreak 回\allowbreak 从} 0.141 {\cjkfont 提\allowbreak 至} 0.941------M1 {\cjkfont 总\allowbreak 恢\allowbreak 复\allowbreak 量} 0.847 {\cjkfont 中\allowbreak 的} 0.800------{\cjkfont 而\allowbreak 保\allowbreak 持\allowbreak 措\allowbreak 辞\allowbreak 不\allowbreak 变\allowbreak ，\allowbreak 仅\allowbreak 政\allowbreak 策\allowbreak 改\allowbreak 变\allowbreak （}v1 $\rightarrow$ v2b{\cjkfont ）\allowbreak 只\allowbreak 提\allowbreak 至} 0.329{\cjkfont 。\allowbreak 聚\allowbreak 合\allowbreak 层\allowbreak 面\allowbreak 两\allowbreak 种\allowbreak 操\allowbreak 纵\allowbreak 贡\allowbreak 献\allowbreak 相\allowbreak 当\allowbreak （}F1 +0.161 {\cjkfont 对} +0.150{\cjkfont ）\allowbreak 。}

{\cjkfont 两\allowbreak 条\allowbreak 范\allowbreak 围\allowbreak 边\allowbreak 界\allowbreak 随\allowbreak 之\allowbreak 而\allowbreak 来\allowbreak 。\allowbreak 其\allowbreak 一\allowbreak ，\allowbreak 这\allowbreak 些\allowbreak 是\allowbreak 简\allowbreak 单\allowbreak 逐\allowbreak 对\allowbreak 差\allowbreak 值\allowbreak ，\allowbreak 不\allowbreak 是\allowbreak 带\allowbreak 交\allowbreak 互\allowbreak 项\allowbreak 估\allowbreak 计\allowbreak 的\allowbreak 正\allowbreak 式\allowbreak 因\allowbreak 子\allowbreak 分\allowbreak 析\allowbreak ，\allowbreak 且\allowbreak 四\allowbreak 个\allowbreak 格\allowbreak 子\allowbreak 在\allowbreak 共\allowbreak 享\allowbreak 开\allowbreak 发\allowbreak 样\allowbreak 本\allowbreak 上\allowbreak 各\allowbreak 运\allowbreak 行\allowbreak 一\allowbreak 次\allowbreak 。\allowbreak 其\allowbreak 二\allowbreak ，\allowbreak 结\allowbreak 果\allowbreak 描\allowbreak 述\allowbreak 的\allowbreak 是\allowbreak 输\allowbreak 出\allowbreak 行\allowbreak 为\allowbreak 而\allowbreak 非\allowbreak 模\allowbreak 型\allowbreak 内\allowbreak 部\allowbreak 机\allowbreak 制\allowbreak ：\allowbreak 疑\allowbreak 问\allowbreak 措\allowbreak 辞\allowbreak 究\allowbreak 竟\allowbreak 影\allowbreak 响\allowbreak 注\allowbreak 意\allowbreak 力\allowbreak 、\allowbreak 表\allowbreak 征\allowbreak 还\allowbreak 是\allowbreak 生\allowbreak 成\allowbreak ，\allowbreak 超\allowbreak 出\allowbreak 本\allowbreak 行\allowbreak 为\allowbreak 实\allowbreak 验\allowbreak 可\allowbreak 确\allowbreak 立\allowbreak 的\allowbreak 范\allowbreak 围\allowbreak 。\allowbreak 可\allowbreak 辩\allowbreak 护\allowbreak 的\allowbreak 解\allowbreak 读\allowbreak 是\allowbreak ：\allowbreak 在\allowbreak 当\allowbreak 前\allowbreak 模\allowbreak 型\allowbreak 与} CAICS {\cjkfont 任\allowbreak 务\allowbreak 上\allowbreak ，\allowbreak 提\allowbreak 示\allowbreak 措\allowbreak 辞\allowbreak 实\allowbreak 质\allowbreak 性\allowbreak 地\allowbreak 改\allowbreak 变\allowbreak 隐\allowbreak 性\allowbreak 类\allowbreak 别\allowbreak 的\allowbreak 召\allowbreak 回\allowbreak ，\allowbreak 两\allowbreak 步\allowbreak 判\allowbreak 断\allowbreak 是\allowbreak 实\allowbreak 用\allowbreak 补\allowbreak 救\allowbreak 。\allowbreak 我\allowbreak 们\allowbreak 不\allowbreak 声\allowbreak 称\allowbreak 已\allowbreak 分\allowbreak 离\allowbreak 因\allowbreak 果\allowbreak 机\allowbreak 制\allowbreak 。}

\textbf{{\cjkfont 表} 8. {\cjkfont 措\allowbreak 辞\allowbreak 与\allowbreak 输\allowbreak 出\allowbreak 政\allowbreak 策\allowbreak 的\allowbreak 四\allowbreak 格\allowbreak 比\allowbreak 较\allowbreak （}n=100{\cjkfont ，\allowbreak 多\allowbreak 数\allowbreak 投\allowbreak 票\allowbreak 真\allowbreak 值\allowbreak ，}DeepSeek V4{\cjkfont ，}temperature 0.0{\cjkfont ）}}

\begingroup\setlength{\tabcolsep}{2pt}\small
\begin{longtable}{>{\raggedright\arraybackslash}p{0.0525\textwidth}>{\raggedright\arraybackslash}p{0.0773\textwidth}>{\raggedright\arraybackslash}p{0.0517\textwidth}>{\raggedright\arraybackslash}p{0.1477\textwidth}>{\raggedright\arraybackslash}p{0.1480\textwidth}>{\raggedright\arraybackslash}p{0.0828\textwidth}>{\raggedright\arraybackslash}p{0.0572\textwidth}>{\raggedright\arraybackslash}p{0.0950\textwidth}>{\raggedright\arraybackslash}p{0.0950\textwidth}>{\raggedright\arraybackslash}p{0.0950\textwidth}}
\hline
{\cjkfont 单\allowbreak 元} & {\cjkfont 措\allowbreak 辞} & {\cjkfont 输\allowbreak 出\allowbreak 政\allowbreak 策} & F1 & {\cjkfont 召\allowbreak 回} & {\cjkfont 精\allowbreak 确} & {\cjkfont 空\allowbreak 输\allowbreak 出} & M1 {\cjkfont 召\allowbreak 回} & M3 {\cjkfont 召\allowbreak 回} & H3 {\cjkfont 召\allowbreak 回} \\
\hline\endfirsthead
\hline
{\cjkfont 单\allowbreak 元} & {\cjkfont 措\allowbreak 辞} & {\cjkfont 输\allowbreak 出\allowbreak 政\allowbreak 策} & F1 & {\cjkfont 召\allowbreak 回} & {\cjkfont 精\allowbreak 确} & {\cjkfont 空\allowbreak 输\allowbreak 出} & M1 {\cjkfont 召\allowbreak 回} & M3 {\cjkfont 召\allowbreak 回} & H3 {\cjkfont 召\allowbreak 回} \\
\hline\endhead
v1 & {\cjkfont 条\allowbreak 件\allowbreak 检\allowbreak 查} & {\cjkfont 保\allowbreak 守} & 0.503 & 0.423 & 0.856 & 21\% & 0.141 & 0.250 & 0.111 \\
v2a & {\cjkfont 两\allowbreak 步\allowbreak 判\allowbreak 断} & {\cjkfont 保\allowbreak 守} & 0.665 & 0.677 & 0.709 & 2\% & 0.941 & 0.667 & 0.511 \\
v2b & {\cjkfont 条\allowbreak 件\allowbreak 检\allowbreak 查} & {\cjkfont 全\allowbreak 覆\allowbreak 盖\allowbreak 偏\allowbreak 好} & 0.653 & 0.641 & 0.724 & 5\% & 0.329 & 0.444 & 0.533 \\
v2 & {\cjkfont 两\allowbreak 步\allowbreak 判\allowbreak 断} & {\cjkfont 全\allowbreak 覆\allowbreak 盖\allowbreak 偏\allowbreak 好} & \textbf{0.656} & \textbf{0.872} & 0.565 & 0\% & 0.988 & 0.972 & 1.000 \\
 & $\Delta$({\cjkfont 措\allowbreak 辞}) = v2a − v1 &  & +0.161 & +0.254 & −0.147 & −19 pp & +0.800 & +0.417 & +0.400 \\
 & $\Delta$({\cjkfont 政\allowbreak 策}) = v2b − v1 &  & +0.150 & +0.218 & −0.132 & −16 pp & +0.188 & +0.194 & +0.422 \\
\hline
\end{longtable}\endgroup
\subsubsection{{\cjkfont 生\allowbreak 产\allowbreak 管\allowbreak 道\allowbreak 的\allowbreak 策\allowbreak 略\allowbreak 选\allowbreak 择}}
{\cjkfont 生\allowbreak 产\allowbreak 编\allowbreak 码\allowbreak 器\allowbreak 采\allowbreak 用} S4 {\cjkfont 两\allowbreak 步\allowbreak 判\allowbreak 断\allowbreak 设\allowbreak 计\allowbreak 而\allowbreak 非} S1{\cjkfont 。\allowbreak 该\allowbreak 选\allowbreak 择\allowbreak 偏\allowbreak 向\allowbreak 召\allowbreak 回\allowbreak （}0.872 {\cjkfont 对} 0.634{\cjkfont ）\allowbreak 与\allowbreak 数\allowbreak 据\allowbreak 库\allowbreak 所\allowbreak 需\allowbreak 的\allowbreak 分\allowbreak 层\allowbreak 证\allowbreak 据\allowbreak 输\allowbreak 出\allowbreak ，\allowbreak 代\allowbreak 价\allowbreak 是\allowbreak 精\allowbreak 确\allowbreak 率\allowbreak （}0.565{\cjkfont ）\allowbreak ，\allowbreak 与} S1 {\cjkfont 的} F1 {\cjkfont 差\allowbreak 距\allowbreak 小\allowbreak 且\allowbreak 置\allowbreak 信\allowbreak 区\allowbreak 间\allowbreak 重\allowbreak 叠\allowbreak 。\allowbreak 三\allowbreak 层\allowbreak 证\allowbreak 据\allowbreak 分\allowbreak 级\allowbreak 方\allowbreak 案\allowbreak （}3.1 {\cjkfont 节\allowbreak ）\allowbreak 即\allowbreak 补\allowbreak 偿\allowbreak 设\allowbreak 计\allowbreak ：\allowbreak 低\allowbreak 精\allowbreak 确\allowbreak 编\allowbreak 码\allowbreak 按\allowbreak 层\allowbreak 级\allowbreak 标\allowbreak 注\allowbreak ，\allowbreak 下\allowbreak 游\allowbreak 用\allowbreak 户\allowbreak 得\allowbreak 以\allowbreak 区\allowbreak 分\allowbreak 逐\allowbreak 字\allowbreak 证\allowbreak 据\allowbreak 编\allowbreak 码\allowbreak 与\allowbreak 先\allowbreak 验\allowbreak 推\allowbreak 断\allowbreak 编\allowbreak 码\allowbreak ，\allowbreak 而\allowbreak 非\allowbreak 将\allowbreak 每\allowbreak 条\allowbreak 标\allowbreak 注\allowbreak 视\allowbreak 为\allowbreak 同\allowbreak 等\allowbreak 支\allowbreak 撑\allowbreak 。}

\subsection{{\cjkfont 实\allowbreak 验\allowbreak 三\allowbreak ：}LLM {\cjkfont 与\allowbreak 一\allowbreak 名\allowbreak 受\allowbreak 训\allowbreak 标\allowbreak 注\allowbreak 者\allowbreak 的\allowbreak 一\allowbreak 致\allowbreak 性}}
{\cjkfont 第\allowbreak 三\allowbreak 个\allowbreak 实\allowbreak 验\allowbreak 测\allowbreak 量} LLM {\cjkfont 的\allowbreak 基\allowbreak 础\allowbreak 字\allowbreak 段\allowbreak 标\allowbreak 注\allowbreak 是\allowbreak 否\allowbreak 与\allowbreak 受\allowbreak 训\allowbreak 人\allowbreak 工\allowbreak 标\allowbreak 注\allowbreak 者\allowbreak 一\allowbreak 致\allowbreak 。\allowbreak 口\allowbreak 径\allowbreak 精\allowbreak 确\allowbreak 表\allowbreak 述\allowbreak ：\allowbreak 这\allowbreak 是\allowbreak 一\allowbreak 项\allowbreak 单\allowbreak 一} LLM {\cjkfont 配\allowbreak 置\allowbreak 与\allowbreak 一\allowbreak 名\allowbreak 标\allowbreak 注\allowbreak 者\allowbreak 、\allowbreak 六\allowbreak 项\allowbreak 字\allowbreak 段\allowbreak （}33 {\cjkfont 项\allowbreak 中\allowbreak ）\allowbreak 的\allowbreak 一\allowbreak 致\allowbreak 性\allowbreak 研\allowbreak 究\allowbreak ，\allowbreak 不\allowbreak 是\allowbreak 专\allowbreak 家\allowbreak 级\allowbreak 准\allowbreak 确\allowbreak 性\allowbreak 的\allowbreak 认\allowbreak 证\allowbreak ，\allowbreak 也\allowbreak 不\allowbreak 涉\allowbreak 及\allowbreak 致\allowbreak 因\allowbreak 编\allowbreak 码\allowbreak 。}

{\cjkfont 在\allowbreak 一\allowbreak 致\allowbreak 性\allowbreak 样\allowbreak 本\allowbreak （}n=289{\cjkfont ）\allowbreak 上\allowbreak ，}LLM {\cjkfont 与\allowbreak 标\allowbreak 注\allowbreak 者\allowbreak 独\allowbreak 立\allowbreak 标\allowbreak 注\allowbreak 六\allowbreak 项\allowbreak 字\allowbreak 段\allowbreak 。\allowbreak 表} 9 {\cjkfont 报\allowbreak 告\allowbreak 结\allowbreak 果\allowbreak 。\allowbreak 三\allowbreak 个\allowbreak 分\allowbreak 类\allowbreak 字\allowbreak 段\allowbreak 的} Cohen's kappa {\cjkfont 为} 0.9136{\cjkfont （\allowbreak 事\allowbreak 故\allowbreak 类\allowbreak 型\allowbreak ）\allowbreak 、}0.9477{\cjkfont （\allowbreak 等\allowbreak 级\allowbreak ）\allowbreak 、}0.9746{\cjkfont （\allowbreak 省\allowbreak 份\allowbreak ）\allowbreak ；\allowbreak 两\allowbreak 个\allowbreak 计\allowbreak 数\allowbreak 字\allowbreak 段\allowbreak 的} ICC(2,1) {\cjkfont 为} 0.9849{\cjkfont （\allowbreak 死\allowbreak 亡\allowbreak ）\allowbreak 与} 0.9853{\cjkfont （\allowbreak 受\allowbreak 伤\allowbreak ）\allowbreak ，\allowbreak 平\allowbreak 均\allowbreak 绝\allowbreak 对\allowbreak 误\allowbreak 差\allowbreak 每\allowbreak 条} 0.118 {\cjkfont 与} 0.087{\cjkfont ；\allowbreak 日\allowbreak 期\allowbreak 年}-{\cjkfont 月\allowbreak 匹\allowbreak 配\allowbreak 率} 99.7\%{\cjkfont 。\allowbreak 三\allowbreak 种\allowbreak 统\allowbreak 计\allowbreak 量\allowbreak 针\allowbreak 对\allowbreak 三\allowbreak 种\allowbreak 字\allowbreak 段\allowbreak ，\allowbreak 我\allowbreak 们\allowbreak 按\allowbreak 各\allowbreak 自\allowbreak 名\allowbreak 称\allowbreak 报\allowbreak 告\allowbreak ，\allowbreak 不\allowbreak 以\allowbreak 单\allowbreak 一} kappa {\cjkfont 标\allowbreak 签\allowbreak 统\allowbreak 称\allowbreak 。}

\textbf{{\cjkfont 表} 9. LLM {\cjkfont 与\allowbreak 一\allowbreak 名\allowbreak 受\allowbreak 训\allowbreak 标\allowbreak 注\allowbreak 者\allowbreak 在\allowbreak 六\allowbreak 项\allowbreak 基\allowbreak 础\allowbreak 字\allowbreak 段\allowbreak 上\allowbreak 的\allowbreak 一\allowbreak 致\allowbreak 性\allowbreak （}n=289{\cjkfont ）}}

\begingroup\setlength{\tabcolsep}{2pt}\small
\begin{longtable}{>{\raggedright\arraybackslash}p{0.1738\textwidth}>{\raggedright\arraybackslash}p{0.2189\textwidth}>{\raggedright\arraybackslash}p{0.3939\textwidth}>{\raggedright\arraybackslash}p{0.1738\textwidth}}
\hline
{\cjkfont 特\allowbreak 征} & {\cjkfont 指\allowbreak 标} & {\cjkfont 数\allowbreak 值} [95\% CI] & {\cjkfont 一\allowbreak 致\allowbreak 率} \\
\hline\endfirsthead
\hline
{\cjkfont 特\allowbreak 征} & {\cjkfont 指\allowbreak 标} & {\cjkfont 数\allowbreak 值} [95\% CI] & {\cjkfont 一\allowbreak 致\allowbreak 率} \\
\hline\endhead
{\cjkfont 事\allowbreak 故\allowbreak 类\allowbreak 型} & Cohen's $\kappa$ & 0.9136 [0.8747, 0.9466] & 93.4\% \\
{\cjkfont 事\allowbreak 故\allowbreak 等\allowbreak 级} & Cohen's $\kappa$ & 0.9477 [0.9024, 0.9841] & 97.6\% \\
{\cjkfont 省\allowbreak 份} & Cohen's $\kappa$ & 0.9746 [0.9563, 0.9927] & 97.6\% \\
{\cjkfont 死\allowbreak 亡\allowbreak 人\allowbreak 数} & ICC(2,1) & 0.9849 [0.9292, 0.9991] & MAE 0.118 \\
{\cjkfont 受\allowbreak 伤\allowbreak 人\allowbreak 数} & ICC(2,1) & 0.9853 [0.9426, 0.9995] & MAE 0.087 \\
{\cjkfont 日\allowbreak 期\allowbreak （\allowbreak 年}-{\cjkfont 月\allowbreak ）} & {\cjkfont 匹\allowbreak 配\allowbreak 率} & 0.9965 & 99.7\% \\
\hline
\end{longtable}\endgroup
{\cjkfont 另\allowbreak 两\allowbreak 条\allowbreak 边\allowbreak 界\allowbreak 。\allowbreak 其\allowbreak 一\allowbreak ，\allowbreak 一\allowbreak 致\allowbreak 性\allowbreak 不\allowbreak 等\allowbreak 于\allowbreak 准\allowbreak 确\allowbreak 性\allowbreak ：\allowbreak 双\allowbreak 方\allowbreak 可\allowbreak 能\allowbreak 共\allowbreak 同\allowbreak 犯\allowbreak 错\allowbreak ，\allowbreak 该\allowbreak 样\allowbreak 本\allowbreak 无\allowbreak 独\allowbreak 立\allowbreak 仲\allowbreak 裁\allowbreak 标\allowbreak 签\allowbreak 。\allowbreak 其\allowbreak 二\allowbreak ，\allowbreak 比\allowbreak 较\allowbreak 仅\allowbreak 覆\allowbreak 盖\allowbreak 六\allowbreak 项\allowbreak 基\allowbreak 础\allowbreak 字\allowbreak 段\allowbreak ；\allowbreak 论\allowbreak 文\allowbreak 的\allowbreak 核\allowbreak 心\allowbreak 标\allowbreak 注\allowbreak 对\allowbreak 象\allowbreak 致\allowbreak 因\allowbreak 编\allowbreak 码\allowbreak 没\allowbreak 有\allowbreak 同\allowbreak 规\allowbreak 模\allowbreak 的\allowbreak 一\allowbreak 致\allowbreak 性\allowbreak 研\allowbreak 究\allowbreak ，\allowbreak 三\allowbreak 专\allowbreak 家\allowbreak 标\allowbreak 注\allowbreak 仅\allowbreak 覆\allowbreak 盖} 100 {\cjkfont 条\allowbreak 策\allowbreak 略\allowbreak 子\allowbreak 集\allowbreak 。\allowbreak 早\allowbreak 期\allowbreak 草\allowbreak 稿\allowbreak 曾\allowbreak 引\allowbreak 用\allowbreak 另\allowbreak 一\allowbreak 批} 100 {\cjkfont 条\allowbreak 三\allowbreak 专\allowbreak 家\allowbreak 样\allowbreak 本\allowbreak 的\allowbreak 专\allowbreak 家\allowbreak 间\allowbreak 区\allowbreak 间\allowbreak 作\allowbreak 为\allowbreak 等\allowbreak 效\allowbreak 界\allowbreak 限\allowbreak ；\allowbreak 因\allowbreak 该\allowbreak 样\allowbreak 本\allowbreak 与} 289 {\cjkfont 条\allowbreak 样\allowbreak 本\allowbreak 在\allowbreak 构\allowbreak 成\allowbreak 与\allowbreak 难\allowbreak 度\allowbreak 上\allowbreak 不\allowbreak 同\allowbreak ，\allowbreak 本\allowbreak 版\allowbreak 不\allowbreak 再\allowbreak 将\allowbreak 其\allowbreak 作\allowbreak 为\allowbreak 界\allowbreak 限\allowbreak 呈\allowbreak 现\allowbreak ，\allowbreak 仅\allowbreak 按\allowbreak 自\allowbreak 身\allowbreak 数\allowbreak 值\allowbreak 报\allowbreak 告\allowbreak 一\allowbreak 致\allowbreak 性\allowbreak 。\allowbreak 图} 6 {\cjkfont 绘\allowbreak 制\allowbreak 各} kappa {\cjkfont 值\allowbreak 及\allowbreak 其\allowbreak 置\allowbreak 信\allowbreak 区\allowbreak 间\allowbreak 。}

\begin{figure}[tbp]
\centering
\includegraphics[width=0.92\textwidth,height=0.80\textheight,keepaspectratio]{/home/suntianhao/construction_agent/.claude/worktrees/classifier-mod/docs/paper_v2/figures/fig6_kappa.png}
\end{figure}
\textbf{{\cjkfont 图} 6. LLM {\cjkfont 与\allowbreak 一\allowbreak 名\allowbreak 标\allowbreak 注\allowbreak 者\allowbreak 在\allowbreak 六\allowbreak 项\allowbreak 字\allowbreak 段\allowbreak 上\allowbreak 的\allowbreak 一\allowbreak 致\allowbreak 性}}
\section{{\cjkfont 讨\allowbreak 论}}
\subsection{{\cjkfont 意\allowbreak 义}}
{\cjkfont 对\allowbreak 数\allowbreak 据\allowbreak 工\allowbreak 程\allowbreak 实\allowbreak 践\allowbreak 有\allowbreak 三\allowbreak 重\allowbreak 推\allowbreak 论\allowbreak ，\allowbreak 均\allowbreak 在\allowbreak 各\allowbreak 自\allowbreak 范\allowbreak 围\allowbreak 内\allowbreak 陈\allowbreak 述\allowbreak 。\allowbreak 其\allowbreak 一\allowbreak ，\allowbreak 三\allowbreak 层\allowbreak 证\allowbreak 据\allowbreak 分\allowbreak 级\allowbreak 方\allowbreak 案\allowbreak 为}"{\cjkfont 每\allowbreak 个\allowbreak 编\allowbreak 码\allowbreak 必\allowbreak 须\allowbreak 带\allowbreak 逐\allowbreak 字\allowbreak 支\allowbreak 撑}"{\cjkfont 的\allowbreak 全\allowbreak 有\allowbreak 或\allowbreak 全\allowbreak 无\allowbreak 声\allowbreak 明\allowbreak 提\allowbreak 供\allowbreak 了\allowbreak 可\allowbreak 用\allowbreak 的\allowbreak 替\allowbreak 代\allowbreak ：\allowbreak 先\allowbreak 验\allowbreak 推\allowbreak 断\allowbreak 编\allowbreak 码\allowbreak 对\allowbreak 下\allowbreak 游\allowbreak 分\allowbreak 析\allowbreak 有\allowbreak 用\allowbreak ，\allowbreak 前\allowbreak 提\allowbreak 是\allowbreak 分\allowbreak 字\allowbreak 段\allowbreak 存\allowbreak 储\allowbreak 、\allowbreak 展\allowbreak 示\allowbreak 与\allowbreak 报\allowbreak 告}------{\cjkfont 管\allowbreak 道\allowbreak 正\allowbreak 是\allowbreak 如\allowbreak 此\allowbreak 。\allowbreak 其\allowbreak 二\allowbreak ，\allowbreak 措\allowbreak 辞\allowbreak 观\allowbreak 察\allowbreak 概\allowbreak 括\allowbreak 为\allowbreak 一\allowbreak 条\allowbreak 设\allowbreak 计\allowbreak 告\allowbreak 诫\allowbreak ：\allowbreak 在\allowbreak 本\allowbreak 模\allowbreak 型\allowbreak 与\allowbreak 本\allowbreak 任\allowbreak 务\allowbreak 上\allowbreak ，\allowbreak 疑\allowbreak 问\allowbreak 式\allowbreak 条\allowbreak 件\allowbreak 检\allowbreak 查\allowbreak 措\allowbreak 辞\allowbreak 抑\allowbreak 制\allowbreak 了\allowbreak 隐\allowbreak 性\allowbreak 类\allowbreak 别\allowbreak 的\allowbreak 同\allowbreak 义\allowbreak 证\allowbreak 据\allowbreak ，\allowbreak 两\allowbreak 步\allowbreak 判\allowbreak 断\allowbreak 措\allowbreak 辞\allowbreak 恢\allowbreak 复\allowbreak 了\allowbreak 其\allowbreak 中\allowbreak 大\allowbreak 部\allowbreak 分\allowbreak （}M1 {\cjkfont 恢\allowbreak 复\allowbreak 总\allowbreak 量} 0.847 {\cjkfont 中\allowbreak 的} 0.800{\cjkfont ）\allowbreak 。\allowbreak 设\allowbreak 计\allowbreak 召\allowbreak 回\allowbreak 关\allowbreak 键\allowbreak 抽\allowbreak 取\allowbreak 提\allowbreak 示\allowbreak 的\allowbreak 实\allowbreak 践\allowbreak 者\allowbreak 应\allowbreak 测\allowbreak 试\allowbreak 措\allowbreak 辞\allowbreak 变\allowbreak 体\allowbreak ，\allowbreak 而\allowbreak 非\allowbreak 假\allowbreak 定\allowbreak 所\allowbreak 列\allowbreak 关\allowbreak 键\allowbreak 词\allowbreak 天\allowbreak 然\allowbreak 作\allowbreak 为\allowbreak 示\allowbreak 例\allowbreak 线\allowbreak 索\allowbreak 起\allowbreak 作\allowbreak 用\allowbreak 。\allowbreak 其\allowbreak 三\allowbreak ，\allowbreak 一\allowbreak 致\allowbreak 性\allowbreak 结\allowbreak 果\allowbreak 支\allowbreak 持\allowbreak 一\allowbreak 项\allowbreak 有\allowbreak 限\allowbreak 的\allowbreak 声\allowbreak 明\allowbreak ：\allowbreak 在\allowbreak 六\allowbreak 项\allowbreak 基\allowbreak 础\allowbreak 字\allowbreak 段\allowbreak 上\allowbreak ，\allowbreak 单\allowbreak 一} LLM {\cjkfont 配\allowbreak 置\allowbreak 与\allowbreak 一\allowbreak 名\allowbreak 受\allowbreak 训\allowbreak 标\allowbreak 注\allowbreak 者\allowbreak 达\allowbreak 到\allowbreak 高\allowbreak 协\allowbreak 议\allowbreak 一\allowbreak 致\allowbreak 性}------{\cjkfont 这\allowbreak 支\allowbreak 持\allowbreak 以} LLM {\cjkfont 在\allowbreak 人\allowbreak 工\allowbreak 监\allowbreak 督\allowbreak 下}\textbf{{\cjkfont 减\allowbreak 轻}}{\cjkfont （\allowbreak 而\allowbreak 非\allowbreak 取\allowbreak 代\allowbreak ）\allowbreak 这\allowbreak 些\allowbreak 字\allowbreak 段\allowbreak 的\allowbreak 标\allowbreak 注\allowbreak 工\allowbreak 作\allowbreak 量\allowbreak 。}

\subsection{{\cjkfont 与\allowbreak 既\allowbreak 有\allowbreak 工\allowbreak 作\allowbreak 的\allowbreak 对\allowbreak 比}}
{\cjkfont 表} 10 {\cjkfont 沿\allowbreak 可\allowbreak 据\allowbreak 已\allowbreak 发\allowbreak 表\allowbreak 报\allowbreak 告\allowbreak 比\allowbreak 较\allowbreak 的\allowbreak 维\allowbreak 度}------{\cjkfont 数\allowbreak 据\allowbreak 来\allowbreak 源\allowbreak 、\allowbreak 记\allowbreak 录\allowbreak 规\allowbreak 模\allowbreak 、\allowbreak 人\allowbreak 工\allowbreak 验\allowbreak 证\allowbreak 、\allowbreak 证\allowbreak 据\allowbreak 绑\allowbreak 定\allowbreak 、\allowbreak 系\allowbreak 统\allowbreak 可\allowbreak 访\allowbreak 问\allowbreak 性}------{\cjkfont 将} CAIEMS {\cjkfont 与\allowbreak 代\allowbreak 表\allowbreak 性\allowbreak 既\allowbreak 有\allowbreak 系\allowbreak 统\allowbreak 并\allowbreak 置\allowbreak 。\allowbreak 我\allowbreak 们\allowbreak 将\allowbreak 其\allowbreak 呈\allowbreak 现\allowbreak 为\allowbreak 描\allowbreak 述\allowbreak 性\allowbreak 对\allowbreak 比\allowbreak ，\allowbreak 而\allowbreak 非\allowbreak 排\allowbreak 他\allowbreak 性\allowbreak 新\allowbreak 颖\allowbreak 性\allowbreak 的\allowbreak 证\allowbreak 据\allowbreak ；\allowbreak 若\allowbreak 干\allowbreak 既\allowbreak 有\allowbreak 系\allowbreak 统\allowbreak 已\allowbreak 组\allowbreak 合\allowbreak 了\allowbreak 采\allowbreak 集\allowbreak 、}LLM {\cjkfont 抽\allowbreak 取\allowbreak 与\allowbreak 图\allowbreak 谱}/{\cjkfont 数\allowbreak 据\allowbreak 库\allowbreak 存\allowbreak 储\allowbreak ，\allowbreak 本\allowbreak 工\allowbreak 作\allowbreak 的\allowbreak 侧\allowbreak 重\allowbreak 在\allowbreak 分\allowbreak 层\allowbreak 证\allowbreak 据\allowbreak 标\allowbreak 注\allowbreak 与\allowbreak 可\allowbreak 更\allowbreak 新\allowbreak 查\allowbreak 询\allowbreak 接\allowbreak 口\allowbreak 。}

\textbf{{\cjkfont 表} 10. {\cjkfont 与\allowbreak 代\allowbreak 表\allowbreak 性\allowbreak 既\allowbreak 有\allowbreak 系\allowbreak 统\allowbreak 的\allowbreak 描\allowbreak 述\allowbreak 性\allowbreak 对\allowbreak 比}}

\begingroup\setlength{\tabcolsep}{2pt}\small
\begin{longtable}{>{\raggedright\arraybackslash}p{0.1932\textwidth}>{\raggedright\arraybackslash}p{0.0835\textwidth}>{\raggedright\arraybackslash}p{0.1941\textwidth}>{\raggedright\arraybackslash}p{0.1641\textwidth}>{\raggedright\arraybackslash}p{0.0623\textwidth}>{\raggedright\arraybackslash}p{0.1164\textwidth}>{\raggedright\arraybackslash}p{0.1179\textwidth}}
\hline
{\cjkfont 系\allowbreak 统} & {\cjkfont 领\allowbreak 域} & LLM & {\cjkfont 编\allowbreak 码\allowbreak 方\allowbreak 案} & {\cjkfont 证\allowbreak 据\allowbreak 绑\allowbreak 定} & {\cjkfont 人\allowbreak 工\allowbreak 验\allowbreak 证} & {\cjkfont 交\allowbreak 互\allowbreak 访\allowbreak 问} \\
\hline\endfirsthead
\hline
{\cjkfont 系\allowbreak 统} & {\cjkfont 领\allowbreak 域} & LLM & {\cjkfont 编\allowbreak 码\allowbreak 方\allowbreak 案} & {\cjkfont 证\allowbreak 据\allowbreak 绑\allowbreak 定} & {\cjkfont 人\allowbreak 工\allowbreak 验\allowbreak 证} & {\cjkfont 交\allowbreak 互\allowbreak 访\allowbreak 问} \\
\hline\endhead
Liu et al. (2025a) & {\cjkfont 航\allowbreak 空} & GPT-4o & HFACS & {\cjkfont 无} & {\cjkfont 专\allowbreak 家\allowbreak 标\allowbreak 注\allowbreak 数\allowbreak 据\allowbreak 集} & {\cjkfont 无} \\
Kim and Han (2024) & {\cjkfont 施\allowbreak 工\allowbreak （\allowbreak 韩\allowbreak ）} & Claude 3.5 & HFACS & {\cjkfont 无} & {\cjkfont 无} & {\cjkfont 无} \\
Ahmadi et al. (2024) & {\cjkfont 施\allowbreak 工\allowbreak （\allowbreak 美\allowbreak ）} & GPT-4.0 & {\cjkfont 无} & {\cjkfont 无} & {\cjkfont 无} & {\cjkfont 无} \\
Liu and Luo (2025) & {\cjkfont 施\allowbreak 工} & GPT-4o & Accimap & {\cjkfont 无} & {\cjkfont 无} & {\cjkfont 无} \\
Jiang et al. (2026) & {\cjkfont 煤\allowbreak 矿} & LLM & {\cjkfont 自\allowbreak 定\allowbreak 本\allowbreak 体} & {\cjkfont 无} & {\cjkfont 无} & {\cjkfont 仅\allowbreak 可\allowbreak 视\allowbreak 化} \\
\textbf{CAIEMS{\cjkfont （\allowbreak 本\allowbreak 文\allowbreak ）}} & \textbf{{\cjkfont 施\allowbreak 工\allowbreak （\allowbreak 中\allowbreak ）}} & \textbf{DeepSeek V4} & \textbf{CAICS{\cjkfont （\allowbreak 操\allowbreak 作\allowbreak 性\allowbreak ）}} & \textbf{{\cjkfont 分\allowbreak 层}} & \textbf{{\cjkfont 部\allowbreak 分\allowbreak （}100 {\cjkfont 条\allowbreak 子\allowbreak 集\allowbreak ）}} & \textbf{Web + API} \\
\hline
\end{longtable}\endgroup
\subsection{{\cjkfont 局\allowbreak 限\allowbreak 性}}
{\cjkfont 以\allowbreak 下\allowbreak 按\allowbreak 严\allowbreak 重\allowbreak 程\allowbreak 度\allowbreak 列\allowbreak 出\allowbreak 界\allowbreak 定\allowbreak 主\allowbreak 要\allowbreak 结\allowbreak 论\allowbreak 的\allowbreak 局\allowbreak 限\allowbreak 性\allowbreak 。}

(1) \textit{{\cjkfont 参\allowbreak 考\allowbreak 标\allowbreak 准\allowbreak 独\allowbreak 立\allowbreak 性\allowbreak 。}} 972 {\cjkfont 条\allowbreak 银\allowbreak 标\allowbreak 准\allowbreak 部\allowbreak 分\allowbreak 由\allowbreak 模\allowbreak 型\allowbreak 产\allowbreak 出\allowbreak （\allowbreak 管\allowbreak 道\allowbreak 版\allowbreak 本\allowbreak 一\allowbreak 致\allowbreak 加\allowbreak 规\allowbreak 则\allowbreak 裁\allowbreak 决\allowbreak ）\allowbreak ；\allowbreak 仅} 100 {\cjkfont 条\allowbreak 致\allowbreak 因\allowbreak 子\allowbreak 集\allowbreak 承\allowbreak 载\allowbreak 独\allowbreak 立\allowbreak 三\allowbreak 专\allowbreak 家\allowbreak 验\allowbreak 证\allowbreak 。\allowbreak 基\allowbreak 于\allowbreak 银\allowbreak 标\allowbreak 准\allowbreak 得\allowbreak 出\allowbreak 的\allowbreak 结\allowbreak 论\allowbreak 继\allowbreak 承\allowbreak 该\allowbreak 边\allowbreak 界\allowbreak ；\allowbreak 完\allowbreak 全\allowbreak 独\allowbreak 立\allowbreak 的\allowbreak 人\allowbreak 工\allowbreak 金\allowbreak 标\allowbreak 准}------{\cjkfont 含\allowbreak 公\allowbreak 开\allowbreak 标\allowbreak 注\allowbreak 指\allowbreak 南\allowbreak 、\allowbreak 标\allowbreak 注\allowbreak 者\allowbreak 背\allowbreak 景\allowbreak 、\allowbreak 盲\allowbreak 法\allowbreak 、\allowbreak 仲\allowbreak 裁\allowbreak 规\allowbreak 则\allowbreak 与\allowbreak 逐\allowbreak 类\allowbreak 样\allowbreak 本\allowbreak 量}------{\cjkfont 仍\allowbreak 是\allowbreak 未\allowbreak 来\allowbreak 工\allowbreak 作\allowbreak 。}

(2) \textit{{\cjkfont 开\allowbreak 发}-{\cjkfont 测\allowbreak 试\allowbreak 重\allowbreak 叠\allowbreak 。}} 100 {\cjkfont 条\allowbreak 策\allowbreak 略\allowbreak 样\allowbreak 本\allowbreak 同\allowbreak 时\allowbreak 承\allowbreak 担\allowbreak 提\allowbreak 示\allowbreak 开\allowbreak 发\allowbreak （}v1 {\cjkfont 诊\allowbreak 断\allowbreak 、}v2 {\cjkfont 设\allowbreak 计\allowbreak 、\allowbreak 消\allowbreak 融\allowbreak ）\allowbreak 与\allowbreak 最\allowbreak 终\allowbreak 报\allowbreak 告\allowbreak ，\allowbreak 三\allowbreak 个\allowbreak 少\allowbreak 样\allowbreak 本\allowbreak 示\allowbreak 例\allowbreak 未\allowbreak 按\allowbreak 显\allowbreak 式\allowbreak 流\allowbreak 程\allowbreak 与\allowbreak 测\allowbreak 试\allowbreak 记\allowbreak 录\allowbreak 隔\allowbreak 离\allowbreak 。\allowbreak 策\allowbreak 略\allowbreak 结\allowbreak 果\allowbreak 因\allowbreak 此\allowbreak 带\allowbreak 有\allowbreak 开\allowbreak 发\allowbreak 集\allowbreak 复\allowbreak 用\allowbreak 风\allowbreak 险\allowbreak ，\allowbreak 应\allowbreak 解\allowbreak 读\allowbreak 为\allowbreak 开\allowbreak 发\allowbreak 评\allowbreak 估\allowbreak ；\allowbreak 与\allowbreak 全\allowbreak 部\allowbreak 提\allowbreak 示\allowbreak 设\allowbreak 计\allowbreak 决\allowbreak 策\allowbreak 隔\allowbreak 离\allowbreak 的\allowbreak 冻\allowbreak 结\allowbreak 测\allowbreak 试\allowbreak 集\allowbreak 计\allowbreak 划\allowbreak 于\allowbreak 下\allowbreak 一\allowbreak 迭\allowbreak 代\allowbreak 建\allowbreak 立\allowbreak 。}

(3) \textit{{\cjkfont 提\allowbreak 示\allowbreak 措\allowbreak 辞\allowbreak 推\allowbreak 断\allowbreak 。}} {\cjkfont 四\allowbreak 格\allowbreak 比\allowbreak 较\allowbreak 在\allowbreak 单\allowbreak 一\allowbreak 样\allowbreak 本\allowbreak 上\allowbreak 报\allowbreak 告\allowbreak 简\allowbreak 单\allowbreak 逐\allowbreak 对\allowbreak 差\allowbreak 值\allowbreak 且\allowbreak 各\allowbreak 运\allowbreak 行\allowbreak 一\allowbreak 次\allowbreak ；\allowbreak 未\allowbreak 估\allowbreak 计\allowbreak 正\allowbreak 式\allowbreak 交\allowbreak 互\allowbreak 效\allowbreak 应\allowbreak ，\allowbreak 也\allowbreak 不\allowbreak 支\allowbreak 持\allowbreak 模\allowbreak 型\allowbreak 内\allowbreak 部\allowbreak 机\allowbreak 制\allowbreak 层\allowbreak 面\allowbreak 的\allowbreak 声\allowbreak 明\allowbreak 。\allowbreak 词\allowbreak 汇\allowbreak 门\allowbreak 控\allowbreak 效\allowbreak 应\allowbreak 作\allowbreak 为\allowbreak 带\allowbreak 实\allowbreak 用\allowbreak 补\allowbreak 救\allowbreak 的\allowbreak 观\allowbreak 察\allowbreak 报\allowbreak 告\allowbreak ，\allowbreak 而\allowbreak 非\allowbreak 已\allowbreak 分\allowbreak 离\allowbreak 的\allowbreak 因\allowbreak 果\allowbreak 机\allowbreak 制\allowbreak 。}

(4) \textit{{\cjkfont 模\allowbreak 型\allowbreak 比\allowbreak 较\allowbreak 条\allowbreak 件\allowbreak 。}} {\cjkfont 四\allowbreak 模\allowbreak 型\allowbreak 运\allowbreak 行\allowbreak 条\allowbreak 件\allowbreak 不\allowbreak 对\allowbreak 等\allowbreak （\allowbreak 强\allowbreak 制\allowbreak 温\allowbreak 度\allowbreak 、\allowbreak 输\allowbreak 出\allowbreak 预\allowbreak 算\allowbreak 、\allowbreak 档\allowbreak 位\allowbreak 路\allowbreak 由\allowbreak ）\allowbreak ，\allowbreak 各\allowbreak 配\allowbreak 置\allowbreak 仅\allowbreak 运\allowbreak 行\allowbreak 一\allowbreak 次\allowbreak ，\allowbreak 模\allowbreak 型\allowbreak 标\allowbreak 识\allowbreak 记\allowbreak 录\allowbreak 在\allowbreak 家\allowbreak 族\allowbreak 粒\allowbreak 度\allowbreak 。\allowbreak 该\allowbreak 比\allowbreak 较\allowbreak 支\allowbreak 持\allowbreak 此\allowbreak 约\allowbreak 束\allowbreak 下\allowbreak 的\allowbreak 选\allowbreak 型\allowbreak 决\allowbreak 策\allowbreak ，\allowbreak 而\allowbreak 非\allowbreak 一\allowbreak 般\allowbreak 性\allowbreak 能\allowbreak 力\allowbreak 排\allowbreak 名\allowbreak 。}

(5) \textit{{\cjkfont 一\allowbreak 致\allowbreak 性\allowbreak 范\allowbreak 围\allowbreak 。}} {\cjkfont 一\allowbreak 致\allowbreak 性\allowbreak 研\allowbreak 究\allowbreak 覆\allowbreak 盖\allowbreak 一\allowbreak 名\allowbreak 标\allowbreak 注\allowbreak 者\allowbreak 、\allowbreak 六\allowbreak 项\allowbreak 字\allowbreak 段\allowbreak （}33 {\cjkfont 项\allowbreak 中\allowbreak ）\allowbreak 、\allowbreak 无\allowbreak 独\allowbreak 立\allowbreak 仲\allowbreak 裁\allowbreak ；\allowbreak 致\allowbreak 因\allowbreak 编\allowbreak 码\allowbreak 缺\allowbreak 大\allowbreak 规\allowbreak 模\allowbreak 一\allowbreak 致\allowbreak 性\allowbreak 研\allowbreak 究\allowbreak 。\allowbreak 结\allowbreak 果\allowbreak 支\allowbreak 持\allowbreak 人\allowbreak 工\allowbreak 监\allowbreak 督\allowbreak 下\allowbreak 基\allowbreak 础\allowbreak 字\allowbreak 段\allowbreak 的\allowbreak 工\allowbreak 作\allowbreak 量\allowbreak 减\allowbreak 轻\allowbreak ，\allowbreak 而\allowbreak 非\allowbreak 专\allowbreak 家\allowbreak 替\allowbreak 代\allowbreak 。}

(6) \textit{{\cjkfont 数\allowbreak 据\allowbreak 库\allowbreak 审\allowbreak 计\allowbreak 。}} {\cjkfont 语\allowbreak 料\allowbreak 存\allowbreak 在\allowbreak 已\allowbreak 知\allowbreak 质\allowbreak 量\allowbreak 缺\allowbreak 口\allowbreak ：\allowbreak 分\allowbreak 类\allowbreak 器\allowbreak 的\allowbreak 多\allowbreak 事\allowbreak 故\allowbreak 残\allowbreak 留\allowbreak 错\allowbreak 误\allowbreak 未\allowbreak 完\allowbreak 全\allowbreak 审\allowbreak 计\allowbreak ，}3.2.4 {\cjkfont 节\allowbreak 所\allowbreak 述\allowbreak 镜\allowbreak 像\allowbreak 报\allowbreak 道\allowbreak 的\allowbreak 四\allowbreak 层\allowbreak 去\allowbreak 重\allowbreak 之\allowbreak 后\allowbreak 未\allowbreak 有\allowbreak 对\allowbreak 残\allowbreak 留\allowbreak 重\allowbreak 复\allowbreak 率\allowbreak 的\allowbreak 独\allowbreak 立\allowbreak 估\allowbreak 计\allowbreak ，\allowbreak 字\allowbreak 段\allowbreak 缺\allowbreak 失\allowbreak 率\allowbreak 跨\allowbreak 字\allowbreak 段\allowbreak 差\allowbreak 异\allowbreak 大\allowbreak （\allowbreak 五\allowbreak 个\allowbreak 枚\allowbreak 举\allowbreak 字\allowbreak 段\allowbreak 非\allowbreak 空\allowbreak 率} 13.7\%--83.9\%{\cjkfont ）\allowbreak ，\allowbreak 且\allowbreak 因\allowbreak 来\allowbreak 源\allowbreak 为\allowbreak 公\allowbreak 共\allowbreak 新\allowbreak 闻\allowbreak 门\allowbreak 户\allowbreak 与\allowbreak 官\allowbreak 方\allowbreak 调\allowbreak 查\allowbreak 报\allowbreak 告\allowbreak ，\allowbreak 语\allowbreak 料\allowbreak 过\allowbreak 度\allowbreak 代\allowbreak 表\allowbreak 严\allowbreak 重\allowbreak 事\allowbreak 故\allowbreak 。\allowbreak 数\allowbreak 据\allowbreak 库\allowbreak 覆\allowbreak 盖\allowbreak 公\allowbreak 开\allowbreak 可\allowbreak 检\allowbreak 索\allowbreak 的\allowbreak 事\allowbreak 故\allowbreak 报\allowbreak 道\allowbreak ；\allowbreak 它\allowbreak 不\allowbreak 是\allowbreak 全\allowbreak 部\allowbreak 施\allowbreak 工\allowbreak 事\allowbreak 故\allowbreak 的\allowbreak 代\allowbreak 表\allowbreak 性\allowbreak 样\allowbreak 本\allowbreak 。}

(7) \textit{{\cjkfont 数\allowbreak 据\allowbreak 治\allowbreak 理\allowbreak 。}} {\cjkfont 管\allowbreak 道\allowbreak 采\allowbreak 集\allowbreak 并\allowbreak 计\allowbreak 划\allowbreak 再\allowbreak 分\allowbreak 发\allowbreak 新\allowbreak 闻\allowbreak 门\allowbreak 户\allowbreak 与\allowbreak 政\allowbreak 府\allowbreak 网\allowbreak 站\allowbreak 的\allowbreak 文\allowbreak 本\allowbreak 与\allowbreak 图\allowbreak 片\allowbreak 。\allowbreak 公\allowbreak 开\allowbreak 可\allowbreak 访\allowbreak 问\allowbreak 本\allowbreak 身\allowbreak 不\allowbreak 足\allowbreak 以\allowbreak 解\allowbreak 决\allowbreak 再\allowbreak 发\allowbreak 布\allowbreak 权\allowbreak 利\allowbreak 问\allowbreak 题\allowbreak ；\allowbreak 逐\allowbreak 源\allowbreak 许\allowbreak 可\allowbreak 、\allowbreak 撤\allowbreak 回\allowbreak 机\allowbreak 制\allowbreak 与\allowbreak 创\allowbreak 伤\allowbreak 影\allowbreak 像\allowbreak 访\allowbreak 问\allowbreak 控\allowbreak 制\allowbreak 见\allowbreak 第} 7 {\cjkfont 节\allowbreak ，\allowbreak 须\allowbreak 在\allowbreak 公\allowbreak 开\allowbreak 前\allowbreak 完\allowbreak 成\allowbreak 。}

(8) \textit{{\cjkfont 可\allowbreak 复\allowbreak 现\allowbreak 性\allowbreak 。}} {\cjkfont 数\allowbreak 据\allowbreak 、\allowbreak 代\allowbreak 码\allowbreak 与\allowbreak 完\allowbreak 整\allowbreak 提\allowbreak 示\allowbreak 承\allowbreak 诺\allowbreak 发\allowbreak 表\allowbreak 后\allowbreak 公\allowbreak 开\allowbreak 而\allowbreak 非\allowbreak 投\allowbreak 稿\allowbreak 时\allowbreak 公\allowbreak 开\allowbreak ；\allowbreak 单\allowbreak 次\allowbreak 运\allowbreak 行\allowbreak 配\allowbreak 置\allowbreak 与\allowbreak 家\allowbreak 族\allowbreak 粒\allowbreak 度\allowbreak 模\allowbreak 型\allowbreak 标\allowbreak 识\allowbreak 限\allowbreak 制\allowbreak 精\allowbreak 确\allowbreak 复\allowbreak 现\allowbreak 。\allowbreak 我\allowbreak 们\allowbreak 承\allowbreak 诺\allowbreak 在\allowbreak 重\allowbreak 投\allowbreak 前\allowbreak 提\allowbreak 供\allowbreak 匿\allowbreak 名\allowbreak 只\allowbreak 读\allowbreak 数\allowbreak 据\allowbreak 库\allowbreak 、\allowbreak 逐\allowbreak 字\allowbreak 提\allowbreak 示\allowbreak 、\allowbreak 最\allowbreak 小\allowbreak 复\allowbreak 现\allowbreak 包\allowbreak 与\allowbreak 系\allowbreak 统\allowbreak 演\allowbreak 示\allowbreak 。}

\section{{\cjkfont 结\allowbreak 论}}
{\cjkfont 本\allowbreak 文\allowbreak 报\allowbreak 告\allowbreak 了} CAIEMS------{\cjkfont 一\allowbreak 条\allowbreak 探\allowbreak 索\allowbreak 性} LLM {\cjkfont 管\allowbreak 道\allowbreak ，\allowbreak 将\allowbreak 公\allowbreak 开\allowbreak 可\allowbreak 获\allowbreak 取\allowbreak 的\allowbreak 中\allowbreak 文\allowbreak 施\allowbreak 工\allowbreak 事\allowbreak 故\allowbreak 叙\allowbreak 述\allowbreak 转\allowbreak 化\allowbreak 为\allowbreak 可\allowbreak 查\allowbreak 询\allowbreak 的\allowbreak 数\allowbreak 据\allowbreak 库\allowbreak 原\allowbreak 型\allowbreak ：}2,623 {\cjkfont 条\allowbreak 记\allowbreak 录\allowbreak ，\allowbreak 取\allowbreak 自\allowbreak 网\allowbreak 络\allowbreak 门\allowbreak 户\allowbreak 与\allowbreak 公\allowbreak 开\allowbreak 事\allowbreak 故\allowbreak 汇\allowbreak 编\allowbreak ，\allowbreak 含} 33 {\cjkfont 字\allowbreak 段\allowbreak 与} 16,323 {\cjkfont 条\allowbreak 三\allowbreak 层\allowbreak 分\allowbreak 级\allowbreak 致\allowbreak 因\allowbreak 标\allowbreak 注\allowbreak ，\allowbreak 经} Web {\cjkfont 界\allowbreak 面\allowbreak 与} REST API {\cjkfont 服\allowbreak 务\allowbreak 。\allowbreak 三\allowbreak 项\allowbreak 各\allowbreak 带\allowbreak 口\allowbreak 径\allowbreak 声\allowbreak 明\allowbreak 的\allowbreak 支\allowbreak 撑\allowbreak 实\allowbreak 验\allowbreak 锚\allowbreak 定\allowbreak 可\allowbreak 靠\allowbreak 性\allowbreak 讨\allowbreak 论\allowbreak ：\allowbreak 支\allowbreak 持\allowbreak 采\allowbreak 用} DeepSeek V4 {\cjkfont 的\allowbreak 约\allowbreak 束\allowbreak 性\allowbreak 四\allowbreak 模\allowbreak 型\allowbreak 比\allowbreak 较\allowbreak （}91.2\% {\cjkfont 基\allowbreak 础\allowbreak 字\allowbreak 段\allowbreak 准\allowbreak 确\allowbreak 率\allowbreak 、\allowbreak 最\allowbreak 低\allowbreak 成\allowbreak 本\allowbreak ）\allowbreak ；\allowbreak 发\allowbreak 现\allowbreak 词\allowbreak 汇\allowbreak 门\allowbreak 控\allowbreak 观\allowbreak 察\allowbreak 及\allowbreak 其\allowbreak 两\allowbreak 步\allowbreak 判\allowbreak 断\allowbreak 补\allowbreak 救\allowbreak 的\allowbreak 六\allowbreak 策\allowbreak 略\allowbreak 提\allowbreak 示\allowbreak 对\allowbreak 比\allowbreak ；\allowbreak 以\allowbreak 及\allowbreak 显\allowbreak 示} LLM {\cjkfont 与\allowbreak 一\allowbreak 名\allowbreak 受\allowbreak 训\allowbreak 标\allowbreak 注\allowbreak 者\allowbreak 在\allowbreak 六\allowbreak 项\allowbreak 基\allowbreak 础\allowbreak 字\allowbreak 段\allowbreak 上\allowbreak 高\allowbreak 一\allowbreak 致\allowbreak 性\allowbreak 的\allowbreak 协\allowbreak 议\allowbreak 研\allowbreak 究\allowbreak 。\allowbreak 全\allowbreak 文\allowbreak 对\allowbreak 结\allowbreak 果\allowbreak 边\allowbreak 界\allowbreak 作\allowbreak 显\allowbreak 式\allowbreak 声\allowbreak 明\allowbreak ：\allowbreak 银\allowbreak 标\allowbreak 准\allowbreak 部\allowbreak 分\allowbreak 由\allowbreak 模\allowbreak 型\allowbreak 产\allowbreak 出\allowbreak 、\allowbreak 开\allowbreak 发\allowbreak 与\allowbreak 测\allowbreak 试\allowbreak 样\allowbreak 本\allowbreak 重\allowbreak 叠\allowbreak 、\allowbreak 不\allowbreak 声\allowbreak 称\allowbreak 因\allowbreak 果\allowbreak 机\allowbreak 制\allowbreak 或\allowbreak 专\allowbreak 家\allowbreak 替\allowbreak 代\allowbreak 。\allowbreak 本\allowbreak 工\allowbreak 作\allowbreak 的\allowbreak 下\allowbreak 一\allowbreak 迭\allowbreak 代\allowbreak 遵\allowbreak 循\allowbreak 明\allowbreak 确\allowbreak 顺\allowbreak 序\allowbreak ：\allowbreak 致\allowbreak 因\allowbreak 编\allowbreak 码\allowbreak 的\allowbreak 独\allowbreak 立\allowbreak 人\allowbreak 工\allowbreak 金\allowbreak 标\allowbreak 准\allowbreak 、\allowbreak 与\allowbreak 提\allowbreak 示\allowbreak 开\allowbreak 发\allowbreak 隔\allowbreak 离\allowbreak 的\allowbreak 冻\allowbreak 结\allowbreak 测\allowbreak 试\allowbreak 集\allowbreak 、\allowbreak 措\allowbreak 辞\allowbreak 问\allowbreak 题\allowbreak 的\allowbreak 正\allowbreak 式\allowbreak 因\allowbreak 子\allowbreak 设\allowbreak 计\allowbreak 、\allowbreak 匹\allowbreak 配\allowbreak 条\allowbreak 件\allowbreak 下\allowbreak 的\allowbreak 多\allowbreak 模\allowbreak 型\allowbreak 重\allowbreak 复\allowbreak 、\allowbreak 以\allowbreak 及\allowbreak 公\allowbreak 开\allowbreak 前\allowbreak 的\allowbreak 完\allowbreak 整\allowbreak 数\allowbreak 据\allowbreak 审\allowbreak 计\allowbreak 与\allowbreak 治\allowbreak 理\allowbreak 。}

\section{{\cjkfont 声\allowbreak 明}}
\textbf{{\cjkfont 数\allowbreak 据\allowbreak 可\allowbreak 用\allowbreak 性\allowbreak 。}} {\cjkfont 数\allowbreak 据\allowbreak 库\allowbreak 的\allowbreak 匿\allowbreak 名\allowbreak 只\allowbreak 读\allowbreak 快\allowbreak 照\allowbreak 、\allowbreak 字\allowbreak 段\allowbreak 字\allowbreak 典\allowbreak 、\allowbreak 来\allowbreak 源\allowbreak 清\allowbreak 单\allowbreak （\allowbreak 构\allowbreak 成\allowbreak 汇\allowbreak 总\allowbreak 见\allowbreak 表} 1{\cjkfont ）\allowbreak 与} CAICS {\cjkfont 操\allowbreak 作\allowbreak 定\allowbreak 义\allowbreak 将\allowbreak 在\allowbreak 审\allowbreak 稿\allowbreak 期\allowbreak 间\allowbreak 应\allowbreak 要\allowbreak 求\allowbreak 提\allowbreak 供\allowbreak ，\allowbreak 并\allowbreak 于\allowbreak 论\allowbreak 文\allowbreak 发\allowbreak 表\allowbreak 后\allowbreak 公\allowbreak 开\allowbreak 。\allowbreak 六\allowbreak 种\allowbreak 提\allowbreak 示\allowbreak 模\allowbreak 板\allowbreak 以\allowbreak 压\allowbreak 缩\allowbreak 形\allowbreak 式\allowbreak 见\allowbreak 附\allowbreak 录} A{\cjkfont ；\allowbreak 逐\allowbreak 字\allowbreak 提\allowbreak 示\allowbreak 包\allowbreak 含\allowbreak 于\allowbreak 复\allowbreak 现\allowbreak 包\allowbreak 。}

\textbf{{\cjkfont 代\allowbreak 码\allowbreak 可\allowbreak 用\allowbreak 性\allowbreak 。}} {\cjkfont 管\allowbreak 道\allowbreak 、\allowbreak 评\allowbreak 测\allowbreak 脚\allowbreak 本\allowbreak 与\allowbreak 实\allowbreak 验\allowbreak 运\allowbreak 行\allowbreak 器\allowbreak 将\allowbreak 于\allowbreak 论\allowbreak 文\allowbreak 发\allowbreak 表\allowbreak 后\allowbreak 以\allowbreak 开\allowbreak 源\allowbreak 许\allowbreak 可\allowbreak 证\allowbreak 随\allowbreak 数\allowbreak 据\allowbreak 发\allowbreak 布\allowbreak ；\allowbreak 审\allowbreak 稿\allowbreak 期\allowbreak 间\allowbreak 可\allowbreak 向\allowbreak 作\allowbreak 者\allowbreak 索\allowbreak 取\allowbreak 最\allowbreak 小\allowbreak 复\allowbreak 现\allowbreak 包\allowbreak 。}

\textbf{{\cjkfont 伦\allowbreak 理\allowbreak 与\allowbreak 数\allowbreak 据\allowbreak 治\allowbreak 理\allowbreak 。}} {\cjkfont 本\allowbreak 研\allowbreak 究\allowbreak 分\allowbreak 析\allowbreak 公\allowbreak 开\allowbreak 新\allowbreak 闻\allowbreak 报\allowbreak 道\allowbreak 与\allowbreak 政\allowbreak 府\allowbreak 发\allowbreak 布\allowbreak 的\allowbreak 事\allowbreak 故\allowbreak 通\allowbreak 报\allowbreak ；\allowbreak 结\allowbreak 构\allowbreak 化\allowbreak 记\allowbreak 录\allowbreak 中\allowbreak 不\allowbreak 涉\allowbreak 及\allowbreak 人\allowbreak 类\allowbreak 受\allowbreak 试\allowbreak 者\allowbreak 或\allowbreak 可\allowbreak 识\allowbreak 别\allowbreak 个\allowbreak 人\allowbreak 信\allowbreak 息\allowbreak 。\allowbreak 但\allowbreak 公\allowbreak 开\allowbreak 可\allowbreak 访\allowbreak 问\allowbreak 并\allowbreak 不\allowbreak 自\allowbreak 动\allowbreak 授\allowbreak 权\allowbreak 批\allowbreak 量\allowbreak 采\allowbreak 集\allowbreak 、\allowbreak 长\allowbreak 期\allowbreak 存\allowbreak 储\allowbreak 与\allowbreak 再\allowbreak 分\allowbreak 发\allowbreak 。\allowbreak 任\allowbreak 何\allowbreak 公\allowbreak 开\allowbreak 发\allowbreak 布\allowbreak 前\allowbreak 须\allowbreak 完\allowbreak 成\allowbreak ：\allowbreak 逐\allowbreak 来\allowbreak 源\allowbreak 许\allowbreak 可\allowbreak 或\allowbreak 合\allowbreak 理\allowbreak 使\allowbreak 用\allowbreak 评\allowbreak 估\allowbreak ；\allowbreak 为\allowbreak 每\allowbreak 条\allowbreak 记\allowbreak 录\allowbreak 保\allowbreak 留\allowbreak 原\allowbreak 始} URL {\cjkfont 与\allowbreak 抓\allowbreak 取\allowbreak 时\allowbreak 间\allowbreak ；\allowbreak 对\allowbreak 图\allowbreak 片\allowbreak 中\allowbreak 的\allowbreak 姓\allowbreak 名\allowbreak 与\allowbreak 可\allowbreak 识\allowbreak 别\allowbreak 面\allowbreak 部\allowbreak 做\allowbreak 移\allowbreak 除\allowbreak 或\allowbreak 遮\allowbreak 蔽\allowbreak ；\allowbreak 应\allowbreak 请\allowbreak 求\allowbreak 提\allowbreak 供\allowbreak 撤\allowbreak 回\allowbreak 删\allowbreak 除\allowbreak 机\allowbreak 制\allowbreak ；\allowbreak 对\allowbreak 创\allowbreak 伤\allowbreak 性\allowbreak 影\allowbreak 像\allowbreak 设\allowbreak 访\allowbreak 问\allowbreak 控\allowbreak 制\allowbreak 。\allowbreak 无\allowbreak 法\allowbreak 确\allowbreak 立\allowbreak 再\allowbreak 发\allowbreak 布\allowbreak 权\allowbreak 利\allowbreak 的\allowbreak 第\allowbreak 三\allowbreak 方\allowbreak 图\allowbreak 片\allowbreak 将\allowbreak 从\allowbreak 任\allowbreak 何\allowbreak 公\allowbreak 开\allowbreak 发\allowbreak 布\allowbreak 中\allowbreak 排\allowbreak 除\allowbreak 。}

\textbf{{\cjkfont 作\allowbreak 者\allowbreak 贡\allowbreak 献\allowbreak （}CRediT{\cjkfont ）\allowbreak 。}} {\cjkfont 孙\allowbreak 天\allowbreak 浩\allowbreak （\allowbreak 南\allowbreak 京\allowbreak 工\allowbreak 业\allowbreak 大\allowbreak 学\allowbreak ）\allowbreak ：\allowbreak 概\allowbreak 念\allowbreak 化\allowbreak 、\allowbreak 方\allowbreak 法\allowbreak 、\allowbreak 软\allowbreak 件\allowbreak 、\allowbreak 数\allowbreak 据\allowbreak 整\allowbreak 理\allowbreak 、\allowbreak 形\allowbreak 式\allowbreak 分\allowbreak 析\allowbreak 、\allowbreak 调\allowbreak 研\allowbreak 、\allowbreak 初\allowbreak 稿\allowbreak 撰\allowbreak 写\allowbreak 。}[{\cjkfont 导\allowbreak 师\allowbreak 姓\allowbreak 名}]{\cjkfont （\allowbreak 南\allowbreak 京\allowbreak 工\allowbreak 业\allowbreak 大\allowbreak 学\allowbreak ）\allowbreak ：\allowbreak 概\allowbreak 念\allowbreak 化\allowbreak 、\allowbreak 指\allowbreak 导\allowbreak 、\allowbreak 验\allowbreak 证\allowbreak 、\allowbreak 审\allowbreak 阅\allowbreak 与\allowbreak 修\allowbreak 订\allowbreak 。}

\textbf{{\cjkfont 利\allowbreak 益\allowbreak 冲\allowbreak 突\allowbreak 声\allowbreak 明\allowbreak 。}} {\cjkfont 作\allowbreak 者\allowbreak 声\allowbreak 明\allowbreak 无\allowbreak 利\allowbreak 益\allowbreak 冲\allowbreak 突\allowbreak 。}

\textbf{{\cjkfont 基\allowbreak 金\allowbreak 资\allowbreak 助\allowbreak 。}} {\cjkfont 本\allowbreak 研\allowbreak 究\allowbreak 未\allowbreak 获\allowbreak 得\allowbreak 公\allowbreak 共\allowbreak 、\allowbreak 商\allowbreak 业\allowbreak 或\allowbreak 非\allowbreak 营\allowbreak 利\allowbreak 部\allowbreak 门\allowbreak 资\allowbreak 助\allowbreak 机\allowbreak 构\allowbreak 的\allowbreak 专\allowbreak 项\allowbreak 资\allowbreak 助\allowbreak 。}

\textbf{AI {\cjkfont 工\allowbreak 具\allowbreak 使\allowbreak 用\allowbreak 声\allowbreak 明\allowbreak 。}} {\cjkfont 大\allowbreak 语\allowbreak 言\allowbreak 模\allowbreak 型\allowbreak 既\allowbreak 是\allowbreak 本\allowbreak 研\allowbreak 究\allowbreak 的\allowbreak 对\allowbreak 象\allowbreak ，\allowbreak 也\allowbreak 是\allowbreak 核\allowbreak 心\allowbreak 技\allowbreak 术\allowbreak 工\allowbreak 具\allowbreak ：\allowbreak 第} 3-4 {\cjkfont 节\allowbreak 报\allowbreak 告\allowbreak 的\allowbreak 全\allowbreak 部\allowbreak 抽\allowbreak 取\allowbreak 、\allowbreak 分\allowbreak 类\allowbreak 与\allowbreak 编\allowbreak 码\allowbreak 结\allowbreak 果\allowbreak 均\allowbreak 由\allowbreak 所\allowbreak 评\allowbreak 测\allowbreak 模\allowbreak 型\allowbreak 经\allowbreak 记\allowbreak 录\allowbreak 在\allowbreak 案\allowbreak 的\allowbreak 管\allowbreak 道\allowbreak 产\allowbreak 出\allowbreak 。\allowbreak 论\allowbreak 文\allowbreak 撰\allowbreak 写\allowbreak 过\allowbreak 程\allowbreak 中\allowbreak 使\allowbreak 用\allowbreak 了} LLM {\cjkfont 写\allowbreak 作\allowbreak 辅\allowbreak 助\allowbreak 进\allowbreak 行\allowbreak 语\allowbreak 言\allowbreak 润\allowbreak 色\allowbreak ；\allowbreak 全\allowbreak 部\allowbreak 科\allowbreak 学\allowbreak 内\allowbreak 容\allowbreak 、\allowbreak 实\allowbreak 验\allowbreak 设\allowbreak 计\allowbreak 与\allowbreak 数\allowbreak 据\allowbreak 出\allowbreak 自\allowbreak 作\allowbreak 者\allowbreak 工\allowbreak 作\allowbreak 并\allowbreak 经\allowbreak 作\allowbreak 者\allowbreak 核\allowbreak 验\allowbreak 。}

\section*{{\cjkfont 参\allowbreak 考\allowbreak 文\allowbreak 献}}
\begin{enumerate}
\item {\cjkfont （\allowbreak 与\allowbreak 英\allowbreak 文\allowbreak 主\allowbreak 源\allowbreak 一\allowbreak 致\allowbreak ，\allowbreak 列\allowbreak 表\allowbreak 见} eswa\_paper\_v2.md References {\cjkfont 节\allowbreak ）}
\end{enumerate}
Ahmadi, A., Sharif, S., Banad, Y. (2025). Improving Aviation Safety Analysis: Automated HFACS Classification Using Reinforcement Learning with Group Relative Policy Optimization. arXiv:2508.21201.

Ahmadi, E., Muley, S., Wang, C. (2024). Automatic Construction Accident Report Analysis Using Large Language Models (LLMs). Journal of Intelligent Construction, 3(1).

Ayhan, B.U., Tokdemir, O.B. (2019). Predicting the outcome of construction incidents. Safety Science, 113, 91--104.

Brown, T., Mann, B., Ryder, N., et al. (2020). Language Models are Few-Shot Learners. NeurIPS, 33, 1877--1901.

Devlin, J., Chang, M.-W., Lee, K., Toutanova, K. (2019). BERT: Pre-training of Deep Bidirectional Transformers for Language Understanding. NAACL-HLT 2019, 4171--4186.

Huang, Y., Yan, R., Zhang, Z. (2026). Automated knowledge extraction from marine accident reports using large language models. Ocean \& Coastal Management, 272, 108015.

Jiang, Y., Nie, B., Wang, J., et al. (2026). Large language model-driven knowledge graph and intelligent analysis for coal mine accidents. Measurement, 277, 121558.

Kim, D., Han, S. (2024). Systemic Analysis of Fall Accidents Caused by Worker Inattention at Construction Sites. KIBIM Magazine, 14(4), 23--34.

Kojima, T., Gu, S.S., Reid, M., et al. (2022). Large Language Models are Zero-Shot Reasoners. NeurIPS, 35, 22199--22213.

Kong, F., Ahn, S. (2024). Use of Knowledge Graphs for Construction Safety Management: A Systematic Literature Review. Information, 15(7), 390.

Liu, Q., Ding, Y., Luo, X. (2025b). Automated knowledge graph-based risk assessment for fall-from-height accidents in construction. Automation in Construction, 179, 106482.

Liu, Q., Li, F., Ng, K.K.H., et al. (2025a). Accident investigation via LLMs reasoning: HFACS-guided Chain-of-Thoughts enhance general aviation safety. Expert Systems with Applications, 269, 126422.

Liu, Q., Luo, X. (2025). Large language model-based knowledge graph construction of unstructured accident reports. EG-ICE 2025.

Onuchukwu, I.S., Gholizadeh, P., et al. (2024). Nature of Occupational Incidents among Roofing Contractors. Buildings, 14(3), 595.

OpenAI (2023). GPT-4 Technical Report. arXiv:2303.08774.

Rafindadi, A.D., Shafiq, N., Othman, I., et al. (2023). Data mining of the essential causes of different types of fatal construction accidents. Heliyon, 9(2), e13389.

Shappell, S.A., Wiegmann, D.A. (2000). The Human Factors Analysis and Classification System--HFACS. Report DOT/FAA/AM-00/7.

Shi, D., Gan, S., Zurada, J., et al. (2025). A multi-model approach to construction site safety. Expert Systems with Applications, 288, 127817.

Togan, V., Mostofi, F., Tokdemir, O.B., et al. (2025). Efficient Management of Safety Documents Using Text-Based Analytics. IEEE Access, 13, 99758--99777.

Wang, X., Wei, J., Schuurmans, D., et al. (2023). Self-Consistency Improves Chain of Thought Reasoning in Language Models. ICLR 2023.

Wei, J., Wang, X., Schuurmans, D., et al. (2022). Chain-of-Thought Prompting Elicits Reasoning in Large Language Models. NeurIPS, 35, 24824--24837.

Xiao, X., Chen, P., Qi, B., et al. (2026). KRAIL: A knowledge-driven framework for human reliability analysis. Reliability Engineering \& System Safety, 265, 111585.

Xuecai, X., Gui, F., Shifei, S., et al. (2024). Accident case data-accident causation model driven safety training method. Process Safety and Environmental Protection, 182, 1208--1226.

Yan, Y., Li, B., Lodewijks, G. (2025). UAV Accident Forensics via HFACS-LLM Reasoning. Drones, 9(10), 704.

Yang, J., Park, S. (2025). LLM-Assisted Knowledge Graph Construction for Managing Road Accident Information. IEEE Access, 13, 217846--217865.

Zhang, L., Hou, Y., Ren, F. (2025a). AIR Agent---A GPT-Based Subway Construction Accident Investigation Report Analysis Chatbot. Buildings, 15(4), 527.

Zhang, J., Zhang, W., Wei, Q., et al. (2025b). GRAG-ProSafe QAS. Expert Systems with Applications.

Zhou, D., Schärli, N., Hou, L., et al. (2023). Least-to-Most Prompting Enables Complex Reasoning in Large Language Models. ICLR 2023.

\appendix
\section*{{\cjkfont 提\allowbreak 示\allowbreak 模\allowbreak 板}}
{\cjkfont 六\allowbreak 种\allowbreak 提\allowbreak 示\allowbreak 策\allowbreak 略\allowbreak 概\allowbreak 述\allowbreak 如\allowbreak 下\allowbreak ；\allowbreak 逐\allowbreak 字\allowbreak 提\allowbreak 示\allowbreak 包\allowbreak 含\allowbreak 于\allowbreak 复\allowbreak 现\allowbreak 包\allowbreak 。}

\textbf{A.1 S0 {\cjkfont 基\allowbreak 础\allowbreak 零\allowbreak 样\allowbreak 本}}{\cjkfont ：}CAICS {\cjkfont 体\allowbreak 系\allowbreak （\allowbreak 十\allowbreak 九\allowbreak 类\allowbreak 定\allowbreak 义\allowbreak ）}+ JSON {\cjkfont 输\allowbreak 出\allowbreak 格\allowbreak 式\allowbreak ，\allowbreak 无\allowbreak 推\allowbreak 理\allowbreak 引\allowbreak 导\allowbreak 。}

\textbf{A.2 S1 {\cjkfont 基\allowbreak 础\allowbreak 少\allowbreak 样\allowbreak 本}}{\cjkfont ：}S0 {\cjkfont 加\allowbreak 三\allowbreak 个\allowbreak 专\allowbreak 家\allowbreak 标\allowbreak 注\allowbreak 示\allowbreak 例\allowbreak （\allowbreak 文\allowbreak 本} $\rightarrow$ {\cjkfont 带\allowbreak 证\allowbreak 据\allowbreak 编\allowbreak 码\allowbreak ）\allowbreak ，\allowbreak 取\allowbreak 自\allowbreak 银\allowbreak 标\allowbreak 准\allowbreak 。}

\textbf{A.3 S2 Auto-CoT}{\cjkfont ：}S0 {\cjkfont 加}"{\cjkfont 按} M$\rightarrow$F$\rightarrow$T$\rightarrow$H$\rightarrow$C {\cjkfont 因\allowbreak 果\allowbreak 链\allowbreak 先\allowbreak 自\allowbreak 组\allowbreak 织\allowbreak 推\allowbreak 理\allowbreak 再\allowbreak 编\allowbreak 码}"{\cjkfont 指\allowbreak 令\allowbreak 。}

\textbf{A.4 S3 {\cjkfont 计\allowbreak 划\allowbreak 与\allowbreak 求\allowbreak 解}}{\cjkfont ：}S0 {\cjkfont 加}"{\cjkfont 先\allowbreak 制\allowbreak 定\allowbreak 分\allowbreak 析\allowbreak 计\allowbreak 划\allowbreak 再\allowbreak 执\allowbreak 行}"{\cjkfont 指\allowbreak 令\allowbreak 。}

\textbf{A.5 S4 CAICS-CoT}{\cjkfont ：\allowbreak 共\allowbreak 测\allowbreak 试\allowbreak 两\allowbreak 版\allowbreak 。}v1 {\cjkfont 将\allowbreak 每\allowbreak 个\allowbreak 编\allowbreak 码\allowbreak 表\allowbreak 述\allowbreak 为\allowbreak 疑\allowbreak 问\allowbreak 式\allowbreak 条\allowbreak 件\allowbreak 检\allowbreak 查\allowbreak （}"{\cjkfont 文\allowbreak 本\allowbreak 提\allowbreak 到}……{\cjkfont 吗\allowbreak ？}"{\cjkfont ）\allowbreak 加\allowbreak 保\allowbreak 守\allowbreak 指\allowbreak 令\allowbreak 。}v2 {\cjkfont 以\allowbreak 两\allowbreak 步\allowbreak 判\allowbreak 断\allowbreak （\allowbreak 显\allowbreak 性\allowbreak 匹\allowbreak 配} + {\cjkfont 授\allowbreak 权\allowbreak 推\allowbreak 断\allowbreak ，\allowbreak 含\allowbreak 从\allowbreak 亡\allowbreak 人\allowbreak 结\allowbreak 果\allowbreak 推\allowbreak 断\allowbreak 监\allowbreak 管\allowbreak 缺\allowbreak 失\allowbreak 等\allowbreak 先\allowbreak 验\allowbreak 规\allowbreak 则\allowbreak ）\allowbreak 取\allowbreak 代\allowbreak 检\allowbreak 查\allowbreak ，\allowbreak 并\allowbreak 加\allowbreak 全\allowbreak 覆\allowbreak 盖\allowbreak 偏\allowbreak 好\allowbreak 。}

\textbf{A.6 S5 CAICS-CoT+}{\cjkfont ：\allowbreak 五\allowbreak 次\allowbreak 独\allowbreak 立\allowbreak 逐\allowbreak 层} API {\cjkfont 调\allowbreak 用\allowbreak （\allowbreak 层\allowbreak 间\allowbreak 摘\allowbreak 要\allowbreak 注\allowbreak 入\allowbreak ）\allowbreak ，\allowbreak 加\allowbreak 最\allowbreak 终\allowbreak 汇\allowbreak 总\allowbreak 调\allowbreak 用\allowbreak 。}

\textit{V2 {\cjkfont 修\allowbreak 订\allowbreak 说\allowbreak 明\allowbreak （\allowbreak 供\allowbreak 作\allowbreak 者\allowbreak 与\allowbreak 导\allowbreak 师\allowbreak 查\allowbreak 阅\allowbreak ，\allowbreak 投\allowbreak 稿\allowbreak 时\allowbreak 删\allowbreak 除\allowbreak 本\allowbreak 段\allowbreak ）\allowbreak ：\allowbreak 本\allowbreak 版\allowbreak 依\allowbreak 据\allowbreak 综\allowbreak 合\allowbreak 评\allowbreak 阅\allowbreak 意\allowbreak 见\allowbreak 完\allowbreak 成\allowbreak 以\allowbreak 下\allowbreak 文\allowbreak 字\allowbreak 级\allowbreak 修\allowbreak 订}------(1){\cjkfont 主\allowbreak 线\allowbreak 确\allowbreak 定\allowbreak 为}"{\cjkfont 探\allowbreak 索\allowbreak 性\allowbreak 数\allowbreak 据\allowbreak 工\allowbreak 程\allowbreak 系\allowbreak 统\allowbreak 与\allowbreak 数\allowbreak 据\allowbreak 资\allowbreak 源}"{\cjkfont ，\allowbreak 删\allowbreak 除}"{\cjkfont 首\allowbreak 次}/{\cjkfont 推\allowbreak 翻}/{\cjkfont 证\allowbreak 明\allowbreak 模\allowbreak 型\allowbreak 依\allowbreak 赖}/{\cjkfont 因\allowbreak 果\allowbreak 机\allowbreak 制}/{\cjkfont 专\allowbreak 家\allowbreak 级}/{\cjkfont 替\allowbreak 代\allowbreak 人\allowbreak 工}/{\cjkfont 全\allowbreak 部\allowbreak 编\allowbreak 码\allowbreak 有\allowbreak 证\allowbreak 据}"{\cjkfont 等\allowbreak 结\allowbreak 论\allowbreak ；}(2){\cjkfont 致\allowbreak 因\allowbreak 标\allowbreak 注\allowbreak 重\allowbreak 构\allowbreak 为} explicit/inferred/prior {\cjkfont 三\allowbreak 层\allowbreak 证\allowbreak 据\allowbreak 分\allowbreak 级\allowbreak （\allowbreak 对\allowbreak 应\allowbreak 数\allowbreak 据\allowbreak 库} evidence\_quality {\cjkfont 字\allowbreak 段\allowbreak ）\allowbreak ，\allowbreak 全\allowbreak 文\allowbreak 统\allowbreak 一\allowbreak 分\allowbreak 层\allowbreak 统\allowbreak 计\allowbreak 口\allowbreak 径\allowbreak ；}(3){\cjkfont 银\allowbreak 标\allowbreak 准\allowbreak 构\allowbreak 建\allowbreak 方\allowbreak 式\allowbreak 如\allowbreak 实\allowbreak 描\allowbreak 述\allowbreak 并\allowbreak 标\allowbreak 注\allowbreak 独\allowbreak 立\allowbreak 性\allowbreak 边\allowbreak 界\allowbreak ；}(4){\cjkfont 开\allowbreak 发\allowbreak 集\allowbreak 与\allowbreak 测\allowbreak 试\allowbreak 集\allowbreak 重\allowbreak 叠\allowbreak 风\allowbreak 险\allowbreak 在} 4.1 {\cjkfont 显\allowbreak 式\allowbreak 声\allowbreak 明\allowbreak ；}(5)2$\times$2 {\cjkfont 消\allowbreak 融\allowbreak 降\allowbreak 级\allowbreak 为}"{\cjkfont 措\allowbreak 辞\allowbreak 敏\allowbreak 感\allowbreak 性\allowbreak 观\allowbreak 察}+{\cjkfont 实\allowbreak 用\allowbreak 补\allowbreak 救}"{\cjkfont ，\allowbreak 删\allowbreak 除\allowbreak 机\allowbreak 制\allowbreak 性\allowbreak 归\allowbreak 因\allowbreak 与}"{\cjkfont 排\allowbreak 除\allowbreak 替\allowbreak 代\allowbreak 解\allowbreak 释}"{\cjkfont 表\allowbreak 述\allowbreak ；}(6){\cjkfont 人\allowbreak 机\allowbreak 一\allowbreak 致\allowbreak 性\allowbreak 限\allowbreak 定\allowbreak 为}"{\cjkfont 一\allowbreak 名\allowbreak 标\allowbreak 注\allowbreak 者\allowbreak 、\allowbreak 六\allowbreak 项\allowbreak 字\allowbreak 段\allowbreak 、\allowbreak 协\allowbreak 议\allowbreak 一\allowbreak 致\allowbreak 性}"{\cjkfont ，}$\kappa$/ICC/{\cjkfont 匹\allowbreak 配\allowbreak 率\allowbreak 分\allowbreak 列\allowbreak 不\allowbreak 再\allowbreak 统\allowbreak 称} Kappa{\cjkfont ；}(7){\cjkfont 全\allowbreak 部\allowbreak 计\allowbreak 数\allowbreak 对\allowbreak 账\allowbreak 闭\allowbreak 合\allowbreak （\allowbreak 类\allowbreak 型} 2,617+{\cjkfont 未\allowbreak 标\allowbreak 注}6{\cjkfont 、\allowbreak 等\allowbreak 级} 2,622+{\cjkfont 未\allowbreak 标\allowbreak 注}1{\cjkfont 、\allowbreak 表}2 {\cjkfont 补\allowbreak 非\allowbreak 空\allowbreak 分\allowbreak 母\allowbreak 列\allowbreak ）\allowbreak ；}(8)M1 {\cjkfont 数\allowbreak 字\allowbreak 全\allowbreak 文\allowbreak 统\allowbreak 一\allowbreak 为} 0.847/0.800{\cjkfont ；}(9){\cjkfont 表\allowbreak 引\allowbreak 用\allowbreak 勘\allowbreak 误\allowbreak （}4.5 {\cjkfont 节\allowbreak 表} 7$\rightarrow${\cjkfont 表} 8{\cjkfont ）\allowbreak 、\allowbreak 表} 6 {\cjkfont 脚\allowbreak 注\allowbreak 与\allowbreak 数\allowbreak 值\allowbreak 矛\allowbreak 盾\allowbreak 修\allowbreak 正\allowbreak ；}(10){\cjkfont 数\allowbreak 据\allowbreak 治\allowbreak 理\allowbreak 声\allowbreak 明\allowbreak 扩\allowbreak 充\allowbreak （\allowbreak 逐\allowbreak 源\allowbreak 许\allowbreak 可}/{\cjkfont 撤\allowbreak 回\allowbreak 机\allowbreak 制}/{\cjkfont 创\allowbreak 伤\allowbreak 影\allowbreak 像\allowbreak 访\allowbreak 问\allowbreak 控\allowbreak 制\allowbreak ）\allowbreak ；}(11){\cjkfont 局\allowbreak 限\allowbreak 性\allowbreak 扩\allowbreak 为} 8 {\cjkfont 项\allowbreak 分\allowbreak 级\allowbreak ；}(12){\cjkfont 新\allowbreak 增} 3.2 {\cjkfont 节}"{\cjkfont 数\allowbreak 据\allowbreak 来\allowbreak 源\allowbreak 与\allowbreak 语\allowbreak 料\allowbreak 构\allowbreak 建}"{\cjkfont （\allowbreak 含} 4 {\cjkfont 个\allowbreak 子\allowbreak 节\allowbreak 与\allowbreak 表} 1{\cjkfont ）\allowbreak ，\allowbreak 端\allowbreak 到\allowbreak 端\allowbreak 呈\allowbreak 现\allowbreak 采\allowbreak 集\allowbreak 、\allowbreak 清\allowbreak 洗\allowbreak 、\allowbreak 筛\allowbreak 选\allowbreak 与\allowbreak 整\allowbreak 合\allowbreak 流\allowbreak 程\allowbreak ，\allowbreak 原} 3.2--3.5 {\cjkfont 节\allowbreak 顺\allowbreak 延\allowbreak 为} 3.3--3.6 {\cjkfont 节\allowbreak 、\allowbreak 表} 1--9 {\cjkfont 顺\allowbreak 延\allowbreak 为\allowbreak 表} 2--10{\cjkfont ，\allowbreak 并\allowbreak 修\allowbreak 正\allowbreak 阶\allowbreak 段} 2 {\cjkfont 的\allowbreak 两\allowbreak 处\allowbreak 表\allowbreak 述\allowbreak （\allowbreak 文\allowbreak 本\allowbreak 类\allowbreak 型\allowbreak 实\allowbreak 为\allowbreak 六\allowbreak 类\allowbreak 而\allowbreak 非\allowbreak 三\allowbreak 类\allowbreak ；}895 {\cjkfont 条\allowbreak 数\allowbreak 字\allowbreak 化\allowbreak 记\allowbreak 录\allowbreak 经\allowbreak 独\allowbreak 立\allowbreak 的\allowbreak 文\allowbreak 献\allowbreak 切\allowbreak 分\allowbreak 路\allowbreak 径\allowbreak 入\allowbreak 库\allowbreak ，\allowbreak 并\allowbreak 非\allowbreak 在\allowbreak 该\allowbreak 阶\allowbreak 段\allowbreak 被\allowbreak 删\allowbreak 除\allowbreak ）\allowbreak 及} 5.3 {\cjkfont 节}(6){\cjkfont 枚\allowbreak 举\allowbreak 字\allowbreak 段\allowbreak 非\allowbreak 空\allowbreak 率\allowbreak 上\allowbreak 限\allowbreak （}52.3\%$\rightarrow$83.9\%{\cjkfont ）\allowbreak 。\allowbreak 六\allowbreak 张\allowbreak 插\allowbreak 图\allowbreak 已\allowbreak 由} `scripts/\allowbreak{}gen\_\allowbreak{}paper\_\allowbreak{}figures.\allowbreak{}py` {\cjkfont 生\allowbreak 成\allowbreak 并\allowbreak 插\allowbreak 入\allowbreak （\allowbreak 图} 1--6{\cjkfont ）\allowbreak ，\allowbreak 图\allowbreak 中\allowbreak 分\allowbreak 布\allowbreak 数\allowbreak 字\allowbreak 随\allowbreak 主\allowbreak 库\allowbreak 现\allowbreak 读\allowbreak ；\allowbreak 表} 2{\cjkfont （}5 {\cjkfont 层} 19 {\cjkfont 类\allowbreak 操\allowbreak 作\allowbreak 定\allowbreak 义\allowbreak ）\allowbreak 已\allowbreak 按\allowbreak 生\allowbreak 产\allowbreak 提\allowbreak 示\allowbreak 词\allowbreak 补\allowbreak 全\allowbreak ，\allowbreak 定\allowbreak 义\allowbreak 与\allowbreak 落\allowbreak 库\allowbreak 类\allowbreak 名\allowbreak 逐\allowbreak 条\allowbreak 核\allowbreak 对\allowbreak 一\allowbreak 致\allowbreak ，\allowbreak 表\allowbreak 注\allowbreak 给\allowbreak 出\allowbreak 实\allowbreak 测\allowbreak 的\allowbreak 各\allowbreak 层\allowbreak 推\allowbreak 断\allowbreak 占\allowbreak 比\allowbreak 。}}

\end{document}